% Connaissance de la loi nationale et internationale relative à l’éthique numérique — Informatique 3ème — Cours signé OnBuch+
% Programme officiel MINESEC. Compiler avec : tectonic N10-connaissance-loi-nationale-internationale-relative-ethi.tex
\newcommand{\DOCMATIERE}{Informatique}
\newcommand{\DOCNIVEAU}{3ème}
\newcommand{\DOCMODULE}{Informatique – classe de 3ème}
\newcommand{\DOCLECON}{N10}
\newcommand{\DOCTITRE}{Connaissance de la loi nationale et internationale relative à l’éthique numérique}
% OnBuch+ — préambule « cours signés » ENRICHI (compilé avec tectonic / XeLaTeX)
% Système visuel « Pop » d'OnBuch+ : fond crème (jamais blanc), bordures encre
% épaisses, ombres « dures » décalées, titres Archivo Black en capitales.
% Version MATHÉMATIQUES : graphes (pgfplots/tikz), boîtes pédagogiques
% (définition, propriété, méthode, attention, le savais-tu, exercice résolu,
% exercice, TP, bilan), page de garde et signature OnBuch+ ancrée en bas.
%
% Macros attendues dans main.tex AVANT \input{preamble} :
%   \DOCMATIERE  \DOCNIVEAU  \DOCMODULE  \DOCLECON (numéro)  \DOCTITRE
\documentclass[11pt]{article}

\usepackage{fontspec}
\usepackage[french]{babel}
\usepackage[a4paper,margin=2cm,top=3.4cm,bottom=2.8cm,headheight=1.4cm,footskip=1.3cm]{geometry}
\usepackage{amsmath,amssymb,amsfonts}
\usepackage{mathtools} % \xmapsto, \coloneqq... souvent utilisés par les modèles
\usepackage{cancel} % simplifications barrées (bilans de forces)
% \abs{z} (module, valeur absolue) et \norm{u} (norme), utilisés par les modèles
\providecommand{\abs}{}\let\abs\relax\DeclarePairedDelimiter{\abs}{\lvert}{\rvert}
\providecommand{\norm}{}\let\norm\relax\DeclarePairedDelimiter{\norm}{\lVert}{\rVert}
\usepackage[table]{xcolor} % [table] : fournit aussi \rowcolors (alternance de lignes)
\usepackage{pagecolor}
\usepackage{tikz}
\usetikzlibrary{arrows.meta,positioning,calc,shapes.geometric,shapes.misc,shapes.arrows,shapes.symbols,decorations.pathmorphing,decorations.pathreplacing,decorations.markings,patterns,shadows,mindmap,trees,fit,backgrounds,matrix,intersections,angles,quotes}
\usepackage{pgfplots}
\usepackage[version=4]{mhchem} % équations nucléaires \ce{^{235}_{92}U}
\usepackage[european,siunitx]{circuitikz} % schémas électriques
\pgfplotsset{compat=1.18}
\usepgfplotslibrary{fillbetween,groupplots} % aires entre courbes (calcul intégral)
\usepackage{enumitem}
\usepackage{fancyhdr}
% [most] charge toutes les bibliothèques tcolorbox (listings, minted, raster,
% documentation...) et gonfle inutilement la mémoire TeX sur un document long
% (~300 boîtes) : on ne charge que ce qui est réellement utilisé.
\usepackage[skins,breakable]{tcolorbox}
\usepackage{graphicx}
\usepackage{colortbl}
\usepackage{booktabs}
\usepackage{tabularx}
\usepackage{diagbox} % case d'en-tête diagonale des tableaux à double entrée
% Colonnes extensibles centrée / gauche / droite (C, L, R), souvent utilisées par les modèles
\newcolumntype{C}{>{\centering\arraybackslash}X}
\newcolumntype{L}{>{\raggedright\arraybackslash}X}
\newcolumntype{R}{>{\raggedleft\arraybackslash}X}
\usepackage{array}
\usepackage{multirow}
\usepackage{multicol}
\usepackage{siunitx}
% parse-numbers=false : accepte aussi \SI{2,5 \times 10^{-3}}{...}
\sisetup{locale=FR,output-decimal-marker={,},group-minimum-digits=4,parse-numbers=false}
\DeclareSIUnit\ohm{\ensuremath{\symup{\Omega}}} % U+2126 absent de la police
\DeclareSIUnit\division{div} % réglages d'oscilloscope (ms/div, V/div)
\DeclareSIUnit\tr{tr} % tours (vitesse de rotation, ex. \SI{1500}{\tr\per\minute})
\DeclareSIUnit\bit{bit}
\DeclareSIUnit\byte{o} % octet
\DeclareSIUnit\inch{po} % pouce
\DeclareSIUnit\ppp{ppp} % points par pouce (résolution d'image)
\usepackage{qrcode}
% Blocs de code source (PHP, C, HTML, JavaScript, SQL) — spécifique à la
% matière Informatique : listings + habillage aux couleurs OnBuch+ (voir le
% style \lstset ci-dessous, à côté des autres boîtes pédagogiques).
\usepackage{listings}
% \degree : commande fréquemment utilisée par les modèles pour un angle ou une
% température en dehors de \SI{}{} (non fournie par les paquets chargés).
\providecommand{\degree}{^{\circ}}
% \qed (fin de démonstration), utilisé par les modèles sans amsthm
\providecommand{\qed}{\hfill$\square$}
% \barre : double barre des valeurs interdites dans un tableau de variations (tkz-tab)
\providecommand{\barre}{\ensuremath{\|}}
% environnement proof (amsthm non chargé) utilisé par les modèles
\ifdefined\proof\else\newenvironment{proof}[1][Démonstration]{\par\noindent\textit{#1.}\ }{\qed\par}\fi
\usepackage{lastpage}
\usepackage{hyperref}
\hypersetup{hidelinks}

% ============================================================
% Palette « Pop » OnBuch+ — mêmes hex que la classe P (plus_ui.dart)
% ============================================================
\definecolor{poporange}{HTML}{F59321}
\definecolor{popdark}{HTML}{E07A0C}
\definecolor{popcream}{HTML}{FFF6E8}
\definecolor{popink}{HTML}{1C1714}
\definecolor{poppurple}{HTML}{7A5AE0}
\definecolor{popgreen}{HTML}{2E9E6A}
\definecolor{poppink}{HTML}{E0457B}
\definecolor{popblue}{HTML}{2F7DE1}
\definecolor{popgold}{HTML}{C98A12}
\definecolor{popmuted}{HTML}{8A7F76}
\definecolor{popline}{HTML}{EADFD2}
\definecolor{popgray}{HTML}{8A7F76}
\colorlet{popgrey}{popgray}
% Alias tolérants (le modèle nomme parfois les couleurs différemment)
\colorlet{popwhite}{white}
\colorlet{popblack}{popink}
\colorlet{popred}{poppink}
\colorlet{popyellow}{popgold}
% Teintes claires pour fonds de tableaux / zones
\colorlet{popblueL}{popblue!10!white}
\colorlet{poporangeL}{poporange!14!white}
\colorlet{popgreenL}{popgreen!12!white}
\colorlet{poppurpleL}{poppurple!12!white}
\colorlet{poppinkL}{poppink!12!white}
\colorlet{popgoldL}{popgold!14!white}

\pagecolor{popcream}

% ============================================================
% Typographie
% ============================================================
\defaultfontfeatures{Ligatures=TeX}
\setmainfont{PlusJakartaSans-Medium.ttf}[Path=../../../fonts/,BoldFont=PlusJakartaSans-Bold.ttf]
% Maths en sans-serif (Fira Math) avec chiffres et lettres droites de Plus
% Jakarta, pour que les formules se fondent dans le texte.
\usepackage[math-style=ISO,bold-style=ISO]{unicode-math}
\setmathfont{FiraMath-Regular.otf}[Path=../../../fonts/]
\setmathfont{PlusJakartaSans-Medium.ttf}[Path=../../../fonts/,range={up/num,up/{Latin,latin},\mathup}]
\usepackage{icomma} % virgule décimale sans espace en mode math
% Caractères mathématiques Unicode absents des polices de texte (Plus Jakarta,
% Archivo Black) : rendus par leur équivalent mathématique, en texte comme en
% formule — sinon ils s'affichent en carré vide (ex. « Arithmétique dans ℤ »).
\usepackage{newunicodechar}
\newunicodechar{ℝ}{\ensuremath{\mathbb{R}}}
\newunicodechar{ℤ}{\ensuremath{\mathbb{Z}}}
\newunicodechar{ℂ}{\ensuremath{\mathbb{C}}}
\newunicodechar{ℕ}{\ensuremath{\mathbb{N}}}
\newunicodechar{ℚ}{\ensuremath{\mathbb{Q}}}
\newunicodechar{𝔻}{\ensuremath{\mathbb{D}}}
\newunicodechar{α}{\ensuremath{\alpha}}
\newunicodechar{β}{\ensuremath{\beta}}
\newunicodechar{γ}{\ensuremath{\gamma}}
\newunicodechar{δ}{\ensuremath{\delta}}
\newunicodechar{ε}{\ensuremath{\varepsilon}}
\newunicodechar{θ}{\ensuremath{\theta}}
\newunicodechar{λ}{\ensuremath{\lambda}}
\newunicodechar{μ}{\ensuremath{\mu}}
\newunicodechar{σ}{\ensuremath{\sigma}}
\newunicodechar{τ}{\ensuremath{\tau}}
\newunicodechar{φ}{\ensuremath{\varphi}}
\newunicodechar{ψ}{\ensuremath{\psi}}
\newunicodechar{ω}{\ensuremath{\omega}}
\newunicodechar{Σ}{\ensuremath{\Sigma}}
% majuscules produites par \MakeUppercase dans les titres (α-aminés -> Α)
\newunicodechar{Α}{\ensuremath{\alpha}}
\newunicodechar{Β}{\ensuremath{\beta}}
\newunicodechar{Γ}{\ensuremath{\Gamma}}
\newunicodechar{Δ}{\ensuremath{\Delta}}
\newunicodechar{Ε}{\ensuremath{\varepsilon}}
\newunicodechar{Θ}{\ensuremath{\Theta}}
\newunicodechar{Λ}{\ensuremath{\Lambda}}
\newunicodechar{Μ}{\ensuremath{\mu}}
\newunicodechar{Τ}{\ensuremath{\tau}}
\newunicodechar{Φ}{\ensuremath{\Phi}}
\newunicodechar{Ψ}{\ensuremath{\Psi}}
\newunicodechar{Ω}{\ensuremath{\Omega}}
\newunicodechar{↦}{\ensuremath{\mapsto}}
\newunicodechar{∈}{\ensuremath{\in}}
\newunicodechar{∉}{\ensuremath{\notin}}
\newunicodechar{∧}{\ensuremath{\wedge}}
\newunicodechar{⇔}{\ensuremath{\Leftrightarrow}}
\newunicodechar{⟺}{\ensuremath{\Longleftrightarrow}}
\newunicodechar{⇒}{\ensuremath{\Rightarrow}}
\newunicodechar{⟹}{\ensuremath{\Longrightarrow}}
\newunicodechar{∘}{\ensuremath{\circ}}
\newunicodechar{∪}{\ensuremath{\cup}}
\newunicodechar{∩}{\ensuremath{\cap}}
\newunicodechar{≡}{\ensuremath{\equiv}}
\newunicodechar{⊂}{\ensuremath{\subset}}
\newunicodechar{⊥}{\ensuremath{\perp}}
\newunicodechar{∃}{\ensuremath{\exists}}
\newunicodechar{∀}{\ensuremath{\forall}}
\newunicodechar{∛}{\ensuremath{\sqrt[3]{\,}}}
\newunicodechar{∜}{\ensuremath{\sqrt[4]{\,}}}
\newunicodechar{⃗}{\ensuremath{{}^{\to}}}
\newunicodechar{ⁿ}{\ifmmode^{n}\else\textsuperscript{\ensuremath{n}}\fi}
\newunicodechar{ˣ}{\ifmmode^{x}\else\textsuperscript{\ensuremath{x}}\fi}
\newunicodechar{ᵇ}{\ifmmode^{b}\else\textsuperscript{\ensuremath{b}}\fi}
\newunicodechar{ʳ}{\ifmmode^{r}\else\textsuperscript{\ensuremath{r}}\fi}
\newunicodechar{ᵘ}{\ifmmode^{u}\else\textsuperscript{\ensuremath{u}}\fi}
\newunicodechar{ʸ}{\ifmmode^{y}\else\textsuperscript{\ensuremath{y}}\fi}
\newunicodechar{ᵃ}{\ifmmode^{a}\else\textsuperscript{\ensuremath{a}}\fi}
\newunicodechar{ᵉ}{\ifmmode^{e}\else\textsuperscript{\ensuremath{e}}\fi}
\newunicodechar{ᵏ}{\ifmmode^{k}\else\textsuperscript{\ensuremath{k}}\fi}
\newunicodechar{ᵖ}{\ifmmode^{p}\else\textsuperscript{\ensuremath{p}}\fi}
\newunicodechar{ᴺ}{\ifmmode^{N}\else\textsuperscript{\ensuremath{N}}\fi}
\newunicodechar{ᶜ}{\ifmmode^{c}\else\textsuperscript{\ensuremath{c}}\fi}
\newunicodechar{ʰ}{\ifmmode^{h}\else\textsuperscript{\ensuremath{h}}\fi}
\newunicodechar{ᵠ}{\ifmmode^{q}\else\textsuperscript{\ensuremath{q}}\fi}
\newunicodechar{ₙ}{\ifmmode_{n}\else\textsubscript{\ensuremath{n}}\fi}
\newunicodechar{ₐ}{\ifmmode_{a}\else\textsubscript{\ensuremath{a}}\fi}
\newunicodechar{ᵢ}{\ifmmode_{i}\else\textsubscript{\ensuremath{i}}\fi}
\newunicodechar{ᵣ}{\ifmmode_{r}\else\textsubscript{\ensuremath{r}}\fi}
\newunicodechar{ₖ}{\ifmmode_{k}\else\textsubscript{\ensuremath{k}}\fi}
\newunicodechar{ᵦ}{\ifmmode_{\beta}\else\textsubscript{\ensuremath{\beta}}\fi}
\newunicodechar{ₓ}{\ifmmode_{x}\else\textsubscript{\ensuremath{x}}\fi}
\newfontfamily\popdisplay{ArchivoBlack-Regular.ttf}[Path=../../../fonts/]
\newfontfamily\poplogo{SpaceGrotesk-Bold.ttf}[Path=../../../fonts/]
\newfontfamily\popsemi{PlusJakartaSans-SemiBold.ttf}[Path=../../../fonts/]
\newfontfamily\popxbold{PlusJakartaSans-ExtraBold.ttf}[Path=../../../fonts/]
\newfontfamily\popxboldls{PlusJakartaSans-ExtraBold.ttf}[Path=../../../fonts/,LetterSpace=3.0]

\setlist[enumerate]{leftmargin=1.7em,itemsep=5pt,topsep=4pt}
\setlist[itemize]{leftmargin=1.7em,itemsep=4pt,topsep=4pt,label={\color{poporange}\textbullet}}
\setlength{\parindent}{0pt}
\setlength{\parskip}{5pt}
\renewcommand{\arraystretch}{1.25}

% pgfplots : style Pop par défaut
\pgfplotsset{
  every axis/.append style={axis line style={popink,line width=1pt},tick style={popink},
    grid style={popline,line width=0.5pt},label style={font=\small},tick label style={font=\footnotesize}},
  popcurve/.style={poporange,line width=1.8pt,no markers,smooth},
}
% Même style disponible dans un \draw TikZ simple (hors pgfplots)
\tikzset{popcurve/.style={poporange,line width=1.8pt}}
\tikzset{popgrid/.style={popline,very thin}}
\colorlet{popcurve}{poporange} % le nom du style est parfois utilisé comme couleur

% ============================================================
% Pill / squiggle
% ============================================================
\newcommand{\poppill}[3]{%
  \tikz[baseline=-0.55ex]\node[draw=popink,line width=1.2pt,rounded corners=2.6pt,%
    fill=#2,inner xsep=6.5pt,inner ysep=3pt]{%
    \popxboldls\fontsize{7}{7}\selectfont\color{#3}\MakeUppercase{#1}};%
}
\newcommand{\popsquiggle}[1]{%
  \tikz[baseline=-0.4ex]{
    \draw[#1,line width=2.6pt,line cap=round]
      plot [smooth] coordinates {(0,0.09) (0.34,0.01) (0.68,0.21) (1.02,0.01) (1.36,0.10)};
  }%
}
% Mot-clé surligné (marqueur orange sous le texte)
\newcommand{\cle}[1]{{\popxbold\color{popdark}#1}}

% ============================================================
% En-tête / pied
% ============================================================
\pagestyle{fancy}
\fancyhf{}
\renewcommand{\headrulewidth}{0pt}
\renewcommand{\footrulewidth}{0pt}
\lhead{\raisebox{-0.15ex}{\poplogo\fontsize{13}{13}\selectfont\color{popink}OnBuch\color{poporange}+}}
\rhead{\poppill{\DOCMATIERE\ \textbullet\ \DOCNIVEAU\ \textbullet\ Leçon \DOCLECON}{white}{popink}}
\lfoot{\popxbold\fontsize{7.6}{7.6}\selectfont\color{popmuted}\MakeUppercase{Cours signé OnBuch+}\ \textbullet\ {\popsemi\DOCTITRE}}
\rfoot{\tikz[baseline=-0.6ex]\node[rounded corners=8.5pt,fill=popink,minimum height=17pt,minimum width=17pt,inner xsep=5pt,inner sep=0]{\popxbold\fontsize{8}{8}\selectfont\color{popcream}\thepage\,/\,\pageref*{LastPage}};}

% ============================================================
% Titres
% ============================================================
\newcommand{\chaptitre}[2]{%
  \begin{tcolorbox}[enhanced,colback=poporange,colframe=popink,boxrule=2.6pt,arc=11pt,
      drop shadow=popink,left=15pt,right=15pt,top=12pt,bottom=11pt,before skip=3pt,after skip=18pt]
    \poppill{#2}{popink}{popcream}\par\vspace{9pt}
    {\raggedright\hyphenpenalty=10000\exhyphenpenalty=10000\popdisplay\fontsize{22}{24}\selectfont\color{popcream}\MakeUppercase{#1}\par}
  \end{tcolorbox}
}
% Section numérotée : pastille orange avec le numéro + titre Archivo
\newcounter{popsec}
\newcommand{\coursec}[1]{%
  \stepcounter{popsec}\setcounter{popsub}{0}%
  \par\vspace{16pt}\noindent
  \begin{minipage}{\linewidth}
  \tikz[baseline=-0.7ex]\node[draw=popink,line width=1.6pt,fill=poporange,rounded corners=4pt,minimum size=22pt,inner sep=0]{\popdisplay\fontsize{12}{12}\selectfont\color{popcream}\thepopsec};%
  \hspace{8pt}{\popdisplay\fontsize{15}{17}\selectfont\color{popink}\MakeUppercase{#1}}\par
  \vspace{4pt}\noindent{\color{poporange}\rule{36pt}{3.6pt}}
  \end{minipage}\par\nopagebreak\vspace{8pt}
}
\newcounter{popsub}
\newcommand{\courssub}[1]{%
  \stepcounter{popsub}%
  \par\vspace{9pt}\noindent{\popxbold\fontsize{11.5}{13}\selectfont\color{popink}{\color{poporange}\thepopsec.\thepopsub}\hspace{6pt}#1}\par\nopagebreak\vspace{4pt}
}

% ============================================================
% Boîtes pédagogiques — même recette : fond blanc, bordure épaisse colorée,
% ombre dure encre, badge plein. Titre optionnel [#1].
% ============================================================
\tcbset{popbase/.style={breakable,enhanced,colback=white,boxrule=2.3pt,arc=9pt,
  shadow={3pt}{-3pt}{0pt}{popink},left=14pt,right=14pt,top=11pt,bottom=11pt,before skip=11pt,after skip=15pt,
  fontupper=\popsemi\fontsize{10.6}{14.5}\selectfont\color{popink},
  fontlower=\popsemi\fontsize{10.6}{14.5}\selectfont\color{popink}}}
% Coche dessinée (le glyphe ✓ n'existe pas dans les polices du document)
\newcommand{\popcheck}[1]{\tikz[baseline=-0.5ex]\draw[#1,line width=1.6pt,line cap=round,line join=round](-0.13,0.0)--(-0.04,-0.09)--(0.13,0.1);}
\providecommand{\coche}{\popcheck{popgreen}}
\providecommand{\resultat}[1]{\par\smallskip\noindent\fcolorbox{popgreen}{popgreenL}{\parbox{\dimexpr\linewidth-2\fboxsep-2\fboxrule\relax}{\textbf{Résultat.} #1}}\par\smallskip}
\newcommand{\popboxhead}[3]{\poppill{#1}{#2}{white}\ifx\relax#3\relax\else\hspace{6pt}{\popxbold\fontsize{10.6}{12}\selectfont\color{popink}#3}\fi\par\vspace{7pt}}

\newtcolorbox{aretenir}[1][]{popbase,colframe=popgreen,before upper={\popboxhead{À retenir}{popgreen}{#1}}}
\newtcolorbox{exemplebox}[1][]{popbase,colframe=poporange,before upper={\popboxhead{Exemple}{poporange}{#1}}}
\newtcolorbox{definition}[1][]{popbase,colframe=popblue,before upper={\popboxhead{Définition}{popblue}{#1}}}
\newtcolorbox{propriete}[1][]{popbase,colframe=poppurple,before upper={\popboxhead{Propriété}{poppurple}{#1}}}
\newtcolorbox{methode}[1][]{popbase,colframe=popgold,before upper={\popboxhead{Méthode}{popgold}{#1}}}
\newtcolorbox{attention}[1][]{popbase,colframe=poppink,before upper={\popboxhead{Attention}{poppink}{#1}}}
\newtcolorbox{experience}[1][]{popbase,colframe=popblue,colback=popblueL,before upper={\popboxhead{Expérience}{popblue}{#1}}}
% « Le savais-tu ? » : carte violette pleine (contexte, culture, Cameroun)
\newtcolorbox{savaistu}[1][]{popbase,colback=poppurple,colframe=popink,
  fontupper=\popsemi\fontsize{10.4}{14.2}\selectfont\color{white},
  before upper={\poppill{Le savais-tu\,?}{white}{poppurple}\ifx\relax#1\relax\else\hspace{6pt}{\popxbold\color{white}#1}\fi\par\vspace{7pt}}}
% Exercice résolu : énoncé en haut, \tcblower puis la solution sur fond crème
\newcounter{popexo}\newcounter{popexr}
\newtcolorbox{exoresolu}[1][]{popbase,colframe=poporange,colbacklower=popcream,
  segmentation style={popink,line width=1.2pt,dashed},
  before upper={\stepcounter{popexr}\popboxhead{Exercice résolu \thepopexr}{poporange}{#1}},
  before lower={\poppill{Solution}{popink}{popcream}\par\vspace{6pt}}}
% Exercice (non résolu en place) — la correction est dans la partie Corrigés
\newtcolorbox{exercice}[1][]{popbase,colframe=popink,
  before upper={\stepcounter{popexo}\popboxhead{Exercice \thepopexo}{popink}{#1}}}
% Corrigé
\newtcolorbox{corrige}[1][]{popbase,colframe=popgreen,colback=popgreenL,
  before upper={\popboxhead{Corrigé}{popgreen}{#1}}}
% Objectifs / prérequis / bilan
\newtcolorbox{objectifs}{popbase,colframe=popink,colback=poporangeL,
  before upper={\popboxhead{Objectifs de la leçon}{popink}{}}}
\newtcolorbox{prerequis}{popbase,colframe=popink,colback=popblueL,
  before upper={\popboxhead{Prérequis}{popblue}{}}}
\newtcolorbox{bilan}{popbase,colframe=popink,colback=white,boxrule=2.8pt,
  before upper={\popboxhead{Fiche bilan}{poporange}{L'essentiel à connaître pour l'examen}}}
% Figure encadrée avec légende
\newtcolorbox{popfigure}[1][]{enhanced,breakable=false,colback=white,colframe=popline,boxrule=1.4pt,arc=8pt,
  left=10pt,right=10pt,top=10pt,bottom=8pt,before skip=10pt,after skip=12pt,halign=center,
  lower separated=false,halign lower=center,
  fontlower=\popsemi\fontsize{9}{11.5}\selectfont\color{popmuted},
  #1}
\newcommand{\legende}[1]{\par\vspace{6pt}{\popsemi\fontsize{9}{11.5}\selectfont\color{popmuted}#1}}

% ============================================================
% Bloc de code source — \begin{lstlisting}[language=Php]...\end{lstlisting}
% (ou language=C, HTML, JavaScript, SQL ; option omise = pseudo-code brut).
% Même recette visuelle que les figures (\popfigure) : cadre encre épais,
% fond clair (popcream), légende possible juste après via \legende{...}
% comme pour une figure (listings ne sait pas arrondir ses coins : on garde
% un cadre rectangulaire, cohérent avec les autres tableaux du document).
% CONTRAT : les modèles ne doivent utiliser QUE \begin{lstlisting}[...] tel
% quel (pas de \begin{tcblisting}, pas de \verb pour un bloc multi-lignes).
% ============================================================
\lstdefinestyle{poplst}{
  backgroundcolor=\color{popcream},
  basicstyle=\popsemi\fontsize{8.6}{11}\selectfont\color{popink},
  keywordstyle=\bfseries\color{popblue},
  commentstyle=\itshape\color{popmuted},
  stringstyle=\color{popgreen},
  numberstyle=\tiny\color{popmuted},
  identifierstyle=\color{popink},
  frame=l,
  rulecolor=\color{popink},
  framerule=1.6pt,
  framesep=7pt,
  framexleftmargin=7pt,
  framexrightmargin=7pt,
  xleftmargin=2pt,
  xrightmargin=2pt,
  breaklines=true,
  breakatwhitespace=false,
  showstringspaces=false,
  showspaces=false,
  showtabs=false,
  tabsize=2,
  columns=flexible,
  keepspaces=true,
  numbers=none,
  aboveskip=10pt,
  belowskip=10pt,
  captionpos=b,
}
\lstset{style=poplst, language=PHP}
% La distribution TeX embarquée ne fournit pas le fichier de langue
% JavaScript de listings (lstlang*.dat) : on la redéfinit nous-mêmes
% pour que \begin{lstlisting}[language=JavaScript] fonctionne.
\lstdefinelanguage{JavaScript}{
  keywords={break,case,catch,class,const,continue,debugger,default,delete,do,
    else,export,extends,finally,for,function,if,import,in,instanceof,let,new,
    return,static,super,switch,this,throw,try,typeof,var,void,while,with,yield,
    async,await,of,get,set},
  keywordstyle=\bfseries\color{popblue},
  ndkeywords={true,false,null,undefined,NaN,Infinity},
  ndkeywordstyle=\bfseries\color{poporange},
  identifierstyle=\color{popink},
  sensitive=true,
  comment=[l]{//},
  morecomment=[s]{/*}{*/},
  string=[b]",
  morestring=[b]',
  morestring=[b]`,
}
% Même limitation pour CSS et les fichiers de configuration (apacheconf,
% nginx, ini...) : ils ne sont pas fournis. CSS et un style générique de
% type "config" (commentaires # ou ;, pas de mots-clés) couvrent les besoins.
\lstdefinelanguage{CSS}{
  keywords={color,background,background-color,margin,padding,border,
    border-radius,font,font-size,font-weight,font-family,width,height,
    display,position,top,left,right,bottom,float,clear,text-align,
    line-height,list-style,cursor,overflow,box-shadow,transition,
    flex,grid,align-items,justify-content,z-index},
  keywordstyle=\bfseries\color{popblue},
  identifierstyle=\color{popink},
  sensitive=false,
  comment=[s]{/*}{*/},
  string=[b]",
  morestring=[b]',
}
\lstdefinelanguage{apacheconf}{
  sensitive=false,
  morecomment=[l]{\#},
}
\lstdefinelanguage{Apache}{
  sensitive=false,
  morecomment=[l]{\#},
}
\lstdefinelanguage{text}{sensitive=false}
\lstdefinelanguage{powershell}{
  keywords={New-Object,New-SmbShare,Get-Acl,Set-Acl,Get-ChildItem,Set-Content,
    Get-Content,Write-Host,ForEach-Object,Where-Object,Select-Object,
    Test-Path,Remove-Item,Copy-Item,Move-Item,Start-Process,Stop-Process,
    If,Else,ElseIf,ForEach,While,Function,Return,Try,Catch},
  keywordstyle=\bfseries\color{popblue},
  identifierstyle=\color{popink},
  sensitive=false,
  comment=[l]{\#},
  string=[b]",
  morestring=[b]',
}
\lstdefinelanguage{ini}{
  sensitive=false,
  morecomment=[l]{;},
  morecomment=[l]{\#},
}

% ============================================================
% Signature OnBuch+
% ============================================================
\newcommand{\poplogoinline}[1]{{\poplogo\fontsize{#1}{#1}\selectfont\color{popink}OnBuch\color{poporange}+}}
\newcommand{\signatureonbuch}{%
  \begin{tcolorbox}[enhanced,colback=popink,colframe=popink,boxrule=2.2pt,arc=10pt,
      drop shadow=poporange,left=14pt,right=14pt,top=10pt,bottom=10pt,before skip=0pt,after skip=0pt]
    \tikz[baseline=-0.7ex]\node[circle,fill=popgreen,draw=popcream,line width=1.3pt,minimum size=22pt,inner sep=0]{\popcheck{white}};%
    \hspace{9pt}\begin{minipage}[c]{0.62\linewidth}
      {\popxbold\fontsize{12}{13}\selectfont\color{popcream}Signé\ \textbullet\ {\poplogo OnBuch\color{poporange}+}}\par\vspace{2pt}
      {\raggedright\popsemi\fontsize{8}{10.5}\selectfont\color{popline}Ludovic A.\\\MakeUppercase{Conforme au programme officiel MINESEC}\par}
    \end{minipage}\hfill
    {\poplogo\fontsize{22}{22}\selectfont\color{popcream}OnBuch\color{poporange}+}
  \end{tcolorbox}%
}

% ============================================================
% Page de garde
% #1 titre · #2 matière · #3 niveau · #4 module · #5 n° leçon · #6 durée ·
% #7 description / objectifs (texte ou itemize)
% ============================================================
\newcommand{\pagegarde}[7]{%
  \thispagestyle{empty}
  \newgeometry{a4paper,margin=1.7cm,top=1.6cm,bottom=1.4cm}%
  \begin{tikzpicture}[remember picture,overlay]
    \fill[poporange!18] ([xshift=-3.2cm,yshift=-3.6cm]current page.north east) circle (4.6cm);
    \fill[poppurple!14] ([xshift=2.4cm,yshift=7.6cm]current page.south west) circle (3.8cm);
  \end{tikzpicture}%
  {\poplogo\fontsize{30}{30}\selectfont\color{popink}OnBuch\color{poporange}+}\hfill
  \raisebox{0.6ex}{\poppill{Cours signé \textbullet\ Premium}{popink}{popcream}}\par
  \vspace{1pt}\popsquiggle{poppurple}\par
  \vspace{8pt}
  \begin{tcolorbox}[enhanced,colback=poporange,colframe=popink,boxrule=2.8pt,arc=13pt,
      drop shadow=popink,left=18pt,right=18pt,top=12pt,bottom=13pt,after skip=10pt]
    \poppill{#2 \textbullet\ #3}{popink}{popcream}\hspace{5pt}\poppill{#4}{white}{popink}\par\vspace{8pt}
    {\popdisplay\fontsize{14}{14}\selectfont\color{popink}LEÇON #5}\par\vspace{4pt}
    {\raggedright\hyphenpenalty=10000\exhyphenpenalty=10000\popdisplay\fontsize{25}{27}\selectfont\color{popcream}\MakeUppercase{#1}\par}
    \vspace{8pt}
    {\popxbold\fontsize{9.5}{9.5}\selectfont\color{popink}\MakeUppercase{Durée conseillée : #6}}
  \end{tcolorbox}
  \begin{tcolorbox}[enhanced,colback=white,colframe=popink,boxrule=2.2pt,arc=10pt,
      drop shadow=popink,left=16pt,right=16pt,top=10pt,bottom=10pt,after skip=8pt]
    \poppill{Dans cette leçon}{poppurple}{white}\par\vspace{5pt}
    \popsemi\fontsize{9.3}{12.6}\selectfont\color{popink}#7
  \end{tcolorbox}
  \noindent\begin{minipage}[t]{0.487\linewidth}
    \begin{tcolorbox}[enhanced,colback=popgreen,colframe=popink,boxrule=2.2pt,arc=10pt,
        drop shadow=popink,left=11pt,right=11pt,top=10pt,bottom=10pt,height=3.1cm,valign=center]
      \tcbox[colback=white,colframe=popink,boxrule=1.3pt,arc=5pt,boxsep=0pt,nobeforeafter,
        left=3pt,right=3pt,top=3pt,bottom=3pt,valign=center]{\qrcode[height=1.9cm]{https://play.google.com/store/apps/details?id=com.luvvix.onbuch&hl=fr}}%
      \hspace{8pt}\begin{minipage}[c]{0.5\linewidth}
        {\popdisplay\fontsize{11}{12.5}\selectfont\color{white}\MakeUppercase{Télécharge OnBuch}}\par\vspace{4pt}
        {\raggedright\popsemi\fontsize{8.4}{11}\selectfont\color{white}Gratuit \textbullet\ Google Play\\Tuteur IA, annales, concours, résultats\par}
      \end{minipage}
    \end{tcolorbox}
  \end{minipage}\hfill
  \begin{minipage}[t]{0.487\linewidth}
    \begin{tcolorbox}[enhanced,colback=poppurple,colframe=popink,boxrule=2.2pt,arc=10pt,
        drop shadow=popink,left=11pt,right=11pt,top=10pt,bottom=10pt,height=3.1cm,valign=center]
      \tikz[baseline=-0.7ex]\node[circle,fill=white,draw=popink,line width=1.3pt,minimum size=1.6cm,inner sep=0]{\popdisplay\fontsize{24}{24}\selectfont\color{poppurple}+};%
      \hspace{8pt}\begin{minipage}[c]{0.6\linewidth}
        {\popdisplay\fontsize{11}{12.5}\selectfont\color{white}\MakeUppercase{Passe à OnBuch+}}\par\vspace{4pt}
        {\raggedright\popsemi\fontsize{8.4}{11}\selectfont\color{white}Tous les cours signés, Tuteur IA illimité, fiches PDF — onglet OnBuch+ de l'app\par}
      \end{minipage}
    \end{tcolorbox}
  \end{minipage}
  \vspace{8pt}
  \popcontenu
  \vfill
  \signatureonbuch
  \restoregeometry
  \newpage
}

% Rangée « Ce cours contient » (page de garde)
\newcommand{\poptile}[3]{% #1 couleur #2 gros texte #3 légende
  \begin{tcolorbox}[enhanced,colback=#1,colframe=popink,boxrule=2pt,arc=9pt,shadow={2.5pt}{-2.5pt}{0pt}{popink},
      left=6pt,right=6pt,top=9pt,bottom=9pt,halign=center,valign=center,height=1.75cm,width=\linewidth,nobeforeafter]
    {\popdisplay\fontsize{12.5}{13}\selectfont\color{white}\MakeUppercase{#2}}\par\vspace{3pt}
    {\centering\popsemi\fontsize{7.8}{9.5}\selectfont\color{white}#3\par}
  \end{tcolorbox}}
\newcommand{\popcontenu}{%
  \par\noindent{\popxboldls\fontsize{7.5}{7.5}\selectfont\color{popmuted}CE COURS SIGNÉ CONTIENT}\par\vspace{7pt}
  \noindent\begin{minipage}[t]{0.235\linewidth}\poptile{popblue}{Cours}{complet, illustré, pas à pas}\end{minipage}\hfill
  \begin{minipage}[t]{0.235\linewidth}\poptile{poporange}{Méthodes}{et exercices résolus}\end{minipage}\hfill
  \begin{minipage}[t]{0.235\linewidth}\poptile{poppink}{Exercices}{type examen + corrigés}\end{minipage}\hfill
  \begin{minipage}[t]{0.235\linewidth}\poptile{popgreen}{Fiche bilan}{l'essentiel à retenir}\end{minipage}\par}

% Sommaire
\newtcolorbox{sommaire}{enhanced,colback=white,colframe=popink,boxrule=2.2pt,arc=10pt,
  drop shadow=popink,left=18pt,right=18pt,top=13pt,bottom=13pt,before skip=6pt,after skip=20pt,
  before upper={\poppill{Sommaire}{poppurple}{white}\par\vspace{9pt}\popsemi\fontsize{10.6}{14.5}\selectfont\color{popink}}}

% Signature de fin de cours — poussée tout en bas de la dernière page
\newcommand{\findecours}{%
  \par\vfill\nopagebreak
  \begin{center}\popsquiggle{poporange}\end{center}\vspace{4pt}
  \signatureonbuch
}
\AtBeginDocument{\def\llbracket{\mathopen{[\mkern-3.5mu[}}\def\rrbracket{\mathclose{]\mkern-3.5mu]}}} % intervalles d'entiers (glyphe ⟦ absent de la police)
% Glyphes absents de Fira Math : redéfinis à partir de glyphes disponibles
\AtBeginDocument{%
  \def\setminus{\mathbin{\backslash}}\let\smallsetminus\setminus
  \def\vdots{\mathord{\vbox{\baselineskip3.2pt\lineskiplimit0pt\kern4pt\hbox{.}\hbox{.}\hbox{.}}}}%
  \def\ddots{\mathinner{\mkern1mu\raise6.5pt\hbox{.}\mkern2mu\raise3.6pt\hbox{.}\mkern2mu\raise0.7pt\hbox{.}\mkern1mu}}%
  \def\triangle{\mathord{\scalebox{0.9}{$\symup{\Delta}$}}}%
  \def\nabla{\mathord{\raisebox{\depth}{\scalebox{1}[-1]{$\symup{\Delta}$}}}}%
  \def\checkmark{\popcheck{popdark}}%
  \def\star{\mathbin{\ast}}\def\bigstar{\mathord{\ast}}%
  \def\diamond{\mathbin{\scalebox{0.6}{$\Diamond$}}}%
  \def\bigcup{\mathop{\mathchoice{\raisebox{-0.3em}{\scalebox{1.7}{$\cup$}}}{\scalebox{1.25}{$\cup$}}{\cup}{\cup}}\displaylimits}%
  \def\bigcap{\mathop{\mathchoice{\raisebox{-0.3em}{\scalebox{1.7}{$\cap$}}}{\scalebox{1.25}{$\cap$}}{\cap}{\cap}}\displaylimits}%
  \def\lesseqqgtr{\mathrel{\lessgtr}}\def\bigoplus{\mathop{\oplus}\displaylimits}%
}
\providecommand{\ssi}{si et seulement si} % abréviation fréquente
\AtBeginDocument{\providecommand{\wideparen}{\overparen}} % arc de cercle

%% FIGFIX-BEGIN v3 — correctifs globaux des figures (docs/onbuch-generation/tools/figfix)
\usepackage{adjustbox}\usepackage{etoolbox}
% toute figure TikZ/pgfplots plus large que la ligne est réduite à la largeur disponible
\BeforeBeginEnvironment{tikzpicture}{\begin{adjustbox}{max width=\linewidth}}
\AfterEndEnvironment{tikzpicture}{\end{adjustbox}}
% légendes pgfplots lisibles par-dessus les courbes
\pgfplotsset{every axis/.append style={legend style={fill=white,fill opacity=0.92,text opacity=1,draw=black!15,inner sep=3pt}}}
% étiquettes d'arêtes (syntaxe edge["..."]) avec fond blanc
\tikzset{every edge quotes/.append style={fill=white,inner sep=1.2pt}}
% bulles de cartes mentales : texte à la ligne dans la bulle, bulle un peu plus grande
\tikzset{every concept/.append style={text width=2.0cm,align=center,minimum size=2.3cm,inner sep=2pt}}
%% FIGFIX-END
\begin{document}

\pagegarde{Connaissance de la loi nationale et internationale relative à l’éthique numérique}{Informatique}{3ème}{Informatique – classe de 3ème}{N10}{1 séance (fiche de progression harmonisée)}{Cette leçon présente les lois nationales (loi de 2010 sur la cybersécurité et la cybercriminalité au Cameroun) et internationales (Convention de Budapest, RGPD) relatives à l'éthique numérique, ainsi que les sanctions encourues. Elle apprend à analyser des situations concrètes de la vie courante et à justifier une réponse argumentée.
\vspace{4pt}
\begin{multicols}{2}\small
\begin{itemize}
\item À la fin de cette leçon, je sais citer les principaux textes de loi nationaux régissant le numérique au Cameroun
\item À la fin de cette leçon, je sais identifier les institutions nationales chargées de la régulation du numérique (ANTIC, ART, tribunaux)
\item À la fin de cette leçon, je sais énumérer les principales infractions numériques et leurs sanctions au niveau national
\item À la fin de cette leçon, je sais présenter les grands instruments internationaux de lutte contre la cybercriminalité (Convention de Budapest, RGPD, Interpol)
\item À la fin de cette leçon, je sais comparer les sanctions prévues au niveau national et international
\item À la fin de cette leçon, je sais appliquer les notions de l'unité à des situations concrètes de la vie courante
\end{itemize}\end{multicols}}

\begin{sommaire}
\begin{multicols}{2}
\begin{enumerate}
\item Cadre juridique national du numérique au Cameroun
\item Les institutions nationales de régulation et de sanction
\item Infractions numériques et sanctions au niveau national
\item Cadre juridique international de l'éthique numérique
\item Comparaison des sanctions nationales et internationales
\item Application à des situations concrètes et démarche argumentée
\item Activité d'intégration
\item Méthodes et exercices résolus
\item Exercices
\item Corrigés des exercices
\item Fiche bilan
\end{enumerate}
\end{multicols}
\end{sommaire}

\begin{objectifs}
\begin{itemize}
\item À la fin de cette leçon, je sais citer les principaux textes de loi nationaux régissant le numérique au Cameroun
\item À la fin de cette leçon, je sais identifier les institutions nationales chargées de la régulation du numérique (ANTIC, ART, tribunaux)
\item À la fin de cette leçon, je sais énumérer les principales infractions numériques et leurs sanctions au niveau national
\item À la fin de cette leçon, je sais présenter les grands instruments internationaux de lutte contre la cybercriminalité (Convention de Budapest, RGPD, Interpol)
\item À la fin de cette leçon, je sais comparer les sanctions prévues au niveau national et international
\item À la fin de cette leçon, je sais appliquer les notions de l'unité à des situations concrètes de la vie courante
\item À la fin de cette leçon, je sais rédiger et justifier une démarche, une réponse ou une conclusion argumentée sur un cas d'infraction numérique
\end{itemize}
\end{objectifs}

\begin{prerequis}
\begin{itemize}
\item Notions d'éthique et de déontologie numériques (leçons précédentes de l'unité)
\item Notion de cybercriminalité et exemples d'infractions numériques (fraude, piratage, diffamation en ligne)
\item Notion de droit et de sanction (enseignement de citoyenneté / ECM)
\item Connaissance des usages courants d'Internet et des réseaux sociaux
\item Notion de protection des données personnelles vue en début d'unité
\end{itemize}
\end{prerequis}

\newpage

\coursec{Cadre juridique national du numérique au Cameroun}

\begin{exemplebox}[Situation d'accroche]
En 2023, un jeune de Douala est poursuivi pour avoir publié sur Facebook de fausses informations accusant un commerçant de son quartier. Le même mois, une entreprise camerounaise voit ses données clients piratées et doit répondre devant l'ANTIC. À l'échelle internationale, un réseau de cybercriminels est démantelé grâce à la coopération d'Interpol.

\textbf{Problématique :} Que risque-t-on concrètement lorsqu'on enfreint l'éthique numérique, au Cameroun et dans le monde ? Qui sanctionne, et comment ?
\end{exemplebox}

Pour répondre à cette problématique, il faut d'abord comprendre que le numérique, comme la route ou le commerce, ne peut fonctionner sans règles. Cette première section présente le \cle{cadre juridique national} : l'ensemble des textes de droit camerounais qui organisent, protègent et sanctionnent les usages de l'informatique et d'Internet.

\courssub{Pourquoi le numérique a besoin de règles de droit}

Imagine un marché sans aucune règle : n'importe qui peut voler la marchandise du voisin, crier des mensonges sur un commerçant ou copier sa marque sans jamais être inquiété. Très vite, plus personne n'oserait y vendre. Internet, c'est exactement pareil : c'est un espace d'échanges (messages, argent, photos, données personnelles) qui ne peut fonctionner que si chacun y est protégé.

\begin{definition}[Règle de droit]
Une \cle{règle de droit} est une règle obligatoire, édictée par l'État, qui s'impose à tous et dont le non-respect est sanctionné par une \cle{peine} (amende, emprisonnement, etc.).
\end{definition}

Le droit numérique poursuit trois grandes missions de protection :
\begin{itemize}
\item \textbf{Protéger les personnes} : leur vie privée, leur honneur, leur image, leurs données personnelles (nom, numéro de téléphone, photos, localisation).
\item \textbf{Protéger les biens} : l'argent (mobile money, comptes bancaires), les œuvres (musique, textes, logiciels), les systèmes d'information des entreprises.
\item \textbf{Protéger l'État} : la sécurité nationale, les systèmes informatiques des administrations, la lutte contre la fraude et le terrorisme en ligne.
\end{itemize}

\begin{attention}[Internet n'est pas un espace sans loi]
Erreur fréquente : croire que « sur Internet, on peut tout faire » parce qu'on est derrière un écran, parfois sous un pseudonyme. C'est \textbf{faux} : toute infraction commise en ligne (insulte, diffamation, vol, escroquerie) est punissable \emph{comme dans la vie réelle}, et les traces numériques (adresse IP, journaux de connexion) permettent souvent d'identifier l'auteur.
\end{attention}

\courssub{Les grandes lois numériques camerounaises}

Le Cameroun s'est doté progressivement d'un arsenal juridique complet. Voici les textes essentiels à connaître.

\textbf{La Constitution camerounaise}

Au sommet de la hiérarchie se trouve la \cle{Constitution} (loi fondamentale du 18 février 1996). Son préambule garantit à tout citoyen des droits fondamentaux directement liés au numérique :
\begin{itemize}
\item la \textbf{protection de la vie privée} et le secret de la correspondance (qui s'étend aux courriels et messages) ;
\item la \textbf{liberté d'expression} et de communication, qui s'exerce aussi en ligne, dans le respect des droits d'autrui.
\end{itemize}
Toutes les autres lois doivent la respecter : c'est la norme suprême.

\textbf{La loi n° 2010/012 du 21 décembre 2010 : cybersécurité et cybercriminalité}

\begin{definition}[Cybersécurité]
La \cle{cybersécurité} est l'ensemble des moyens \textbf{techniques} (antivirus, chiffrement, pare-feu), \textbf{juridiques} (lois, sanctions) et \textbf{humains} (formation, vigilance des utilisateurs) destinés à protéger les systèmes d'information contre les attaques et les abus.
\end{definition}

\begin{definition}[Cybercriminalité]
La \cle{cybercriminalité} est l'ensemble des infractions pénales commises \textbf{par le biais} d'un système d'information (escroquerie en ligne, diffusion de fausses informations) ou \textbf{contre} un système d'information (piratage, destruction de données).
\end{definition}

\begin{propriete}[Objet et champ d'application de la loi n° 2010/012]
Cette loi fixe le cadre de la sécurité des réseaux et systèmes d'information au Cameroun. Elle s'applique :
\begin{itemize}
\item aux personnes physiques et morales utilisant les réseaux de communications électroniques sur le territoire national ;
\item aux opérateurs (CAMTEL, MTN, Orange...) et fournisseurs de services, tenus de coopérer avec la justice (conservation et communication des données de connexion) ;
\item à la répression des infractions commises via les TIC : accès frauduleux à un système, interception de données, atteinte à la vie privée, diffusion de fausses nouvelles, etc.
\end{itemize}
\end{propriete}

\textbf{La loi n° 2010/021 du 21 décembre 2010 : le commerce électronique}

Promulguée le même jour, cette loi encadre les \textbf{achats et ventes en ligne}. Elle définit notamment la valeur juridique de l'\cle{écrit électronique} et de la \cle{signature électronique}, les obligations du vendeur en ligne (identifier clairement son entreprise, informer le client) et la protection du consommateur. Grâce à elle, un contrat conclu par Internet a la même valeur qu'un contrat papier.

\textbf{Le Code pénal camerounais (loi de 2016 modifiant le code)}

Le \cle{Code pénal} est le texte qui définit les infractions et leurs peines. Réformé par la loi n° 2016/007 du 12 juillet 2016, il réprime des infractions très souvent commises aujourd'hui via les TIC :
\begin{itemize}
\item la \textbf{diffamation} : imputer à quelqu'un un fait qui porte atteinte à son honneur (ex. accuser faussement un commerçant sur Facebook) ;
\item l'\textbf{escroquerie} : obtenir de l'argent par des manœuvres frauduleuses (faux concours, fausses annonces, hameçonnage) ;
\item l'\textbf{atteinte à la vie privée} : capter, conserver ou diffuser l'image ou les paroles d'une personne sans son consentement ;
\item l'\textbf{atteinte à l'honneur} par voie de communication électronique.
\end{itemize}

\begin{savaistu}[Un pionnier en Afrique centrale]
Le Cameroun a été l'un des premiers pays d'Afrique centrale à se doter, dès 2010, d'une loi spécifique sur la cybersécurité et la cybercriminalité, et a créé l'ANTIC (Agence Nationale des Technologies de l'Information et de la Communication) pour veiller à son application. Cette avance a servi de modèle à plusieurs pays de la sous-région.
\end{savaistu}

\courssub{La pyramide des normes et la chronologie des textes}

Les règles de droit ne sont pas toutes au même niveau : elles s'organisent en \cle{hiérarchie des normes}. Une règle de niveau inférieur ne peut jamais contredire une règle de niveau supérieur.

\begin{popfigure}
\begin{center}
\begin{tikzpicture}[scale=1.0]
% Pyramide en 3 étages
\fill[popblue!85] (-1.6,4.6) -- (1.6,4.6) -- (0,6.2) -- cycle;
\fill[poporange!80] (-3.0,2.8) -- (3.0,2.8) -- (1.6,4.6) -- (-1.6,4.6) -- cycle;
\fill[popgreen!75] (-4.4,1.0) -- (4.4,1.0) -- (3.0,2.8) -- (-3.0,2.8) -- cycle;
% Textes
\node[white, font=\bfseries\small] at (0,5.15) {Constitution (1996)};
\node[white, font=\bfseries\small] at (0,3.7) {Lois de 2010 et Code pénal (2016)};
\node[white, font=\bfseries\small] at (0,1.9) {Décrets, arrêtés, décisions de l'ANTIC};
% Annotations latérales
\node[right, font=\footnotesize, text=popdark, align=left] at (5.0,5.15) {Norme suprême :\\ vie privée, libertés};
\node[right, font=\footnotesize, text=popdark, align=left] at (5.0,3.7) {Votées par\\ le Parlement};
\node[right, font=\footnotesize, text=popdark, align=left] at (5.0,1.9) {Précisent l'application\\ des lois};
\draw[->, thick, popmuted] (4.6,5.15) -- (1.8,5.15);
\draw[->, thick, popmuted] (4.6,3.7) -- (3.2,3.7);
\draw[->, thick, popmuted] (4.6,1.9) -- (4.5,1.9);
\end{tikzpicture}
\end{center}
\legende{La pyramide des normes juridiques du numérique au Cameroun : chaque niveau doit respecter le niveau supérieur.}
\end{popfigure}

\begin{popfigure}
\begin{center}
\begin{tikzpicture}[scale=1.0]
% Ligne du temps
\draw[->, very thick, popblue] (0,0) -- (13.2,0);
% Jalons
\foreach \x/\annee in {1.5/1996, 5/2010, 8.5/2016}{
  \draw[thick, popdark] (\x,-0.15) -- (\x,0.15);
  \node[below, font=\bfseries\small, text=popdark] at (\x,-0.2) {\annee};
}
% Événements alternés
\node[above, font=\footnotesize, text=popdark, align=center, text width=3.2cm] at (1.5,0.35) {Constitution :\\ protection de la vie privée et des libertés};
\draw[popmuted] (1.5,0.15) -- (1.5,0.35);
\node[above, font=\footnotesize, text=popdark, align=center, text width=4.2cm] at (5,0.35) {Lois n° 2010/012 (cybersécurité et cybercriminalité) et n° 2010/021 (commerce électronique)};
\draw[popmuted] (5,0.15) -- (5,0.35);
\node[above, font=\footnotesize, text=popdark, align=center, text width=3.4cm] at (8.5,0.35) {Loi n° 2016/007 :\\ Code pénal réprimant les infractions via les TIC};
\draw[popmuted] (8.5,0.15) -- (8.5,0.35);
\node[below, font=\footnotesize\itshape, text=popmuted, align=center] at (11.8,-0.2) {renforcement\\ progressif};
\end{tikzpicture}
\end{center}
\legende{Frise chronologique des grands textes juridiques du numérique au Cameroun.}
\end{popfigure}

\courssub{Vocabulaire juridique de base : infraction, peine, amende}

Pour lire un texte de loi ou un jugement, il faut maîtriser trois mots-clés.

\begin{definition}[Infraction, peine, amende]
\begin{itemize}
\item Une \cle{infraction} est un acte interdit par la loi et puni par elle (ex. pirater un compte, diffamer quelqu'un en ligne).
\item Une \cle{peine} est la sanction prononcée par un tribunal contre l'auteur d'une infraction : emprisonnement, amende, travaux d'intérêt général, etc.
\item Une \cle{amende} est une peine \textbf{pécuniaire} : le condamné doit payer une somme d'argent à l'État (en francs CFA au Cameroun).
\end{itemize}
\end{definition}

\begin{aretenir}[L'essentiel]
\begin{itemize}
\item Le droit numérique protège les \textbf{personnes}, les \textbf{biens} et l'\textbf{État}.
\item La \textbf{Constitution} (1996) garantit la vie privée et les libertés ; elle domine toutes les autres normes.
\item La \textbf{loi n° 2010/012} réprime la cybercriminalité ; la \textbf{loi n° 2010/021} encadre le commerce électronique ; le \textbf{Code pénal} (réformé en 2016) punit diffamation, escroquerie et atteinte à la vie privée, y compris en ligne.
\item Une infraction en ligne est une infraction comme une autre : elle expose son auteur à des peines de prison et/ou des amendes.
\end{itemize}
\end{aretenir}

\begin{exoresolu}[Identifier la règle de droit applicable]
\textbf{Énoncé.} Une élève de 3ème découvre qu'un camarade a créé un faux compte Facebook à son nom et y publie des messages mensongers la présentant comme voleuse. 1) Quelles infractions sont commises ? 2) Quels textes camerounais permettent de poursuivre l'auteur ? 3) Que risque-t-il ?
\tcblower
\textbf{Solution détaillée.}
\begin{enumerate}
\item \textbf{Infractions commises :} l'usurpation d'identité (création d'un faux compte au nom d'autrui), la diffamation (accusation mensongère de vol portant atteinte à l'honneur) et l'atteinte à la vie privée (utilisation de son identité sans consentement).
\item \textbf{Textes applicables :} la loi n° 2010/012 du 21 décembre 2010 relative à la cybersécurité et à la cybercriminalité, car les faits sont commis \emph{par le biais d'un système d'information} ; et le Code pénal camerounais (réformé en 2016), qui réprime la diffamation et l'atteinte à la vie privée, y compris par voie électronique. La victime peut déposer plainte ; les opérateurs et la plateforme peuvent être requis de fournir les données de connexion de l'auteur.
\item \textbf{Sanctions encourues :} l'auteur risque des \textbf{peines d'emprisonnement} et/ou des \textbf{amendes} prévues par ces textes, ainsi que des dommages et intérêts au profit de la victime. Le fait d'être mineur ne supprime pas la responsabilité : il relève alors de la justice des mineurs.
\end{enumerate}
\textbf{Conclusion :} même entre élèves et « pour rire », un faux compte diffamatoire est une infraction réelle, poursuivie sur la base de la loi de 2010 et du Code pénal.
\end{exoresolu}

\begin{exemplebox}[Les fausses nouvelles sur WhatsApp]
Publier ou relayer volontairement de \textbf{fausses nouvelles} sur WhatsApp (fausse alerte, rumeur accusant une personne, fausse annonce de décès) n'est pas anodin : l'auteur peut être poursuivi sur la base de la \textbf{loi n° 2010/012} (diffusion de fausses informations par voie électronique) et du \textbf{Code pénal} (diffamation, atteinte à l'honneur). Avant de partager un message, la bonne démarche est donc : vérifier la source, vérifier l'information sur un site fiable, et s'abstenir de relayer en cas de doute.
\end{exemplebox}

\begin{exercice}[À toi de jouer]
1) Définis la cybersécurité et la cybercriminalité en une phrase chacune. 2) Replace sur la pyramide des normes : un arrêté de l'ANTIC, la Constitution, la loi n° 2010/021. 3) Un vendeur en ligne refuse de livrer un article payé par mobile money : quelle loi de 2010 protège l'acheteur ? Justifie ta réponse.
\end{exercice}

\begin{corrige}[Exercice]
1) La cybersécurité est l'ensemble des moyens techniques, juridiques et humains destinés à protéger les systèmes d'information ; la cybercriminalité est l'ensemble des infractions pénales commises par le biais ou contre un système d'information. 2) Sommet : la Constitution ; milieu : la loi n° 2010/021 ; base : l'arrêté de l'ANTIC. 3) C'est la loi n° 2010/021 régissant le commerce électronique, car elle encadre les ventes en ligne et protège le consommateur (obligations du vendeur, valeur juridique du contrat électronique).
\end{corrige}

\coursec{Les institutions nationales de régulation et de sanction}

Après avoir étudié le cadre juridique qui pose les règles du jeu numérique au Cameroun (loi n°2010/012 sur la cybercriminalité, loi n°2010/013 sur les communications électroniques, loi n°2024/016 sur la protection des données personnelles), il est indispensable de comprendre \emph{qui} fait appliquer ces règles. Une loi sans organe de contrôle ni juridiction compétente reste une lettre morte. Le Cameroun a mis en place un écosystème institutionnel cohérent où chaque acteur a une mission précise : le MINPOSTEL pilote la politique, l'ANTIC sécurise et régule les systèmes d'information, l'ART régule les opérateurs télécoms, les forces de l'ordre enquêtent, et les tribunaux sanctionnent. Cette section présente ces institutions, leurs prérogatives et la manière dont elles s'articulent lorsqu'une infraction numérique est commise.

\courssub{Le MINPOSTEL : pilote de la politique numérique nationale}

Le Ministère des Postes et Télécommunications (MINPOSTEL) est l'autorité politique suprême en matière de numérique au Cameroun. Il ne gère pas directement les incidents techniques, mais il définit la stratégie, prépare les textes législatifs et réglementaires, et veille à leur mise en œuvre cohérente sur le territoire national.

\begin{definition}[MINPOSTEL]
Le Ministère des Postes et Télécommunications est l'administration centrale chargée de concevoir, d'élaborer et de suivre la mise en œuvre de la politique de l'État dans les domaines de la poste, des télécommunications et des technologies de l'information et de la communication. Il assure la tutelle technique des agences de régulation (ANTIC, ART) et représente le Cameroun dans les instances internationales (UIT, UA, CEMAC).
\end{definition}

Ses missions principales sont :
\begin{itemize}
    \item L'élaboration des projets de lois et décrets relatifs au secteur numérique (ex. : loi sur la cybercriminalité, décret sur la protection des données).
    \item La définition de la stratégie nationale « Cameroun Numérique 2030 » et le suivi de sa mise en œuvre.
    \item La coordination des actions des structures sous tutelle (ANTIC, ART, CAMTEL, SONATEL, etc.).
    \item La représentation de l'État dans les négociations internationales sur la gouvernance de l'Internet et les télécommunications.
\end{itemize}

\courssub{L'ANTIC : gardienne de la sécurité et de la régulation des systèmes d'information}

L'Agence Nationale des Technologies de l'Information et de la Communication (ANTIC) est l'organe technique central pour la sécurité des systèmes d'information (SSI) et la régulation des activités TIC. Créée par le décret n°2002/001 du 4 janvier 2002 et réorganisée par le décret n°2012/041 du 9 mars 2012, elle est placée sous la tutelle technique du MINPOSTEL.

\begin{definition}[ANTIC]
L'Agence Nationale des Technologies de l'Information et de la Communication est un établissement public à caractère administratif, doté de la personnalité morale et de l'autonomie financière, chargé de la régulation, du contrôle et du suivi des actions de l'État dans le domaine des technologies de l'information et de la communication, notamment en matière de sécurité des systèmes d'information, de certification électronique et de promotion de l'économie numérique.
\end{definition}

\begin{aretenir}[Missions clés de l'ANTIC]
\begin{enumerate}
    \item \textbf{Sécurité des systèmes d'information (SSI) :} Élaboration de la politique nationale de SSI, audit de sécurité, gestion des incidents (CERT-CM), certification des produits et prestataires de confiance.
    \item \textbf{Régulation et contrôle :} Attribution des noms de domaine \texttt{.cm}, agrément des fournisseurs de services de certification électronique (PSCE), contrôle de la conformité des systèmes d'information des administrations.
    \item \textbf{Promotion de l'économie numérique :} Appui aux startups, formation, sensibilisation du grand public (campagnes « Stop Cybercrime »).
    \item \textbf{Coopération internationale :} Point de contact national pour la cybercriminalité (réseau 24/7 de la Convention de Budapest).
\end{enumerate}
\end{aretenir}

\begin{savaistu}[L'ANTIC et la certification des administrations]
L'ANTIC délivre des \emph{attestations de conformité} et certifie la sécurité des systèmes d'information des administrations camerounaises. Depuis 2020, tout marché public impliquant un système d'information sensible doit présenter une attestation de l'ANTIC. L'agence gère également l'infrastructure à clés publiques (PKI) nationale via la \emph{Autorité de Certification Racine du Cameroun (ACR-CM)}, ce qui permet la signature électronique juridiquement reconnue (loi n°2010/012, art. 7).
\end{savaistu}

\courssub{L'ART : régulateur des communications électroniques et des opérateurs}

L'Agence de Régulation des Télécommunications (ART) a un périmètre distinct de l'ANTIC : elle se concentre sur le \emph{transport} de l'information (réseaux, fréquences, opérateurs) et non sur le \emph{contenu} ou la sécurité logique des systèmes d'information. Créée par la loi n°98/014 du 14 juillet 1998, renforcée par la loi n°2010/013, elle est l'équivalent camerounais de l'ARCEP française ou de la FCC américaine.

\begin{definition}[ART]
L'Agence de Régulation des Télécommunications est l'organe de régulation du secteur des communications électroniques au Cameroun. Elle veille au respect de la concurrence loyale, protège les consommateurs, gère le spectre des fréquences radioélectriques et les ressources de numérotation, et délivre les licences et autorisations aux opérateurs.
\end{definition}

\begin{attention}[Ne pas confondre ANTIC et ART]
C'est l'erreur la plus fréquente chez les élèves (et parfois dans la presse) :
\begin{itemize}
    \item \textbf{ANTIC} = Sécurité des \emph{systèmes d'information} (données, logiciels, certificats, noms de domaine \texttt{.cm}, cybercriminalité technique, CERT). \emph{Question clé : « Mon serveur a-t-il été piraté ? Mon certificat est-il valide ? »}
    \item \textbf{ART} = Régulation des \emph{communications électroniques} (opérateurs MTN/Orange/Camtel, qualité de service, tarifs, fréquences 4G/5G, numérotation, interconnexion). \emph{Question clé : « Mon forfait est-il respecté ? L'antenne relais perturbe-t-elle mon quartier ? »}
\end{itemize}
Elles collaborent (ex. : fraude sur Mobile Money = ART pour l'opérateur + ANTIC pour la fraude informatique + Justice pour la sanction), mais leurs prérogatives légales sont distinctes.
\end{attention}

\courssub{Les forces de l'ordre : enquêteurs de terrain sur le cyberespace}

Lorsque qu'une infraction numérique est commise (escroquerie, usurpation d'identité, diffusion de fausses nouvelles, piratage, etc.), ce sont les services de police judiciaire qui mènent l'enquête sous l'autorité du Procureur de la République. Deux structures sont particulièrement mobilisées :

\begin{itemize}
    \item \textbf{La Brigade des Investigations Spéciales (BIS) de la Police Judiciaire :} Unité d'élite de la Sûreté Nationale, spécialisée dans les enquêtes complexes (cybercriminalité, criminalité financière, terrorisme). Elle dispose d'experts en criminalistique numérique (analyse de disques durs, téléphones, journaux de connexion, traçage IP).
    \item \textbf{La Gendarmerie Nationale (notamment le Centre de Recherche Judiciaire) :} Compétente sur l'ensemble du territoire, y compris en zone rurale. Elle traite aussi les infractions numériques (ex. : arnaques Mobile Money en zone périurbaine) et dispose de sections de recherche spécialisées.
\end{itemize}

\begin{propriete}[Pouvoirs d'enquête en matière numérique (Code de procédure pénale + Loi n°2010/012)]
Les officiers de police judiciaire (OPJ) habilités peuvent, sur réquisition du Procureur ou sur instruction du Juge d'instruction :
\begin{enumerate}
    \item Procéder à la saisie et à la copie des supports numériques (ordinateurs, smartphones, serveurs).
    \item Exiger la remise des données de connexion et d'identification auprès des fournisseurs d'accès et opérateurs (conservation 1 an selon loi n°2010/013).
    \item Recourir à des experts assermentés (souvent formés par l'ANTIC) pour l'analyse forensique.
    \item Effectuer des interceptions de correspondances électroniques (sur autorisation écrite du Procureur Général, pour infractions graves, durée limitée).
\end{enumerate}
Toute preuve obtenue irrégulièrement (sans réquisition, hors procédure) est irrecevable au tribunal (principe de loyauté de la preuve).
\end{propriete}

\courssub{Les tribunaux camerounais : juge et sanction}

L'aboutissement de la chaîne institutionnelle est l'ordre judiciaire. Au Cameroun, les tribunaux de droit commun (Tribunaux de Première Instance - TPI, Tribunaux de Grande Instance - TGI) sont compétents pour juger les infractions de la loi n°2010/012 sur la cybercriminalité. Il n'existe \emph{pas} de tribunal spécialisé « cyber » distinct ; ce sont les juges du parquet et du siège classiques qui appliquent le droit pénal spécial.

\begin{aretenir}[Hiérarchie juridictionnelle pour le numérique]
\begin{itemize}
    \item \textbf{TPI / TGI (premier ressort) :} Jugement des contraventions et délits (escroquerie en ligne, atteintes aux systèmes, diffamation numérique). Peines : amendes, emprisonnement jusqu'à 10 ans pour certains délits.
    \item \textbf{Cours d'Appel (deuxième ressort) :} Examen des appels interjetés contre les jugements des TPI/TGI.
    \item \textbf{Cour Suprême (cassation) :} Contrôle de la bonne application de la loi (pas de réexamen des faits).
\end{itemize}
Le Parquet (Procureur de la République, Procureur Général) dirige l'action publique, décide des poursuites (classement sans suite, citation directe, ouverture d'information judiciaire) et requiert les peines.
\end{aretenir}

\courssub{La procédure simplifiée : de la plainte au jugement}

Comprendre le parcours d'une plainte permet à tout citoyen de savoir comment agir et à quel guichet s'adresser. Le schéma ci-dessous synthétise les étapes types pour une infraction numérique (ex. : arnaque Mobile Money, piratage de compte Facebook, diffusion de fausses nouvelles).

\begin{popfigure}
\begin{tikzpicture}[
    node distance=1.2cm and 1.5cm,
    box/.style={rectangle, draw=popdark, thick, fill=popcream, rounded corners=6pt, minimum width=3.2cm, minimum height=1.1cm, align=center, font=\small},
    arrow/.style={->, >=Stealth, thick, popdark},
    decision/.style={diamond, draw=poporange, thick, fill=poporangeL, aspect=2, minimum width=3.5cm, minimum height=1.3cm, align=center, font=\small},
    startstop/.style={ellipse, draw=popgreen, thick, fill=popgreenL, minimum width=2.5cm, minimum height=1cm, align=center, font=\small}
]
\node[startstop, align=center] (start){Victime\\constate l'infraction};
\node[box, below=of start, align=center] (preuve){1. Conservation des preuves\\ (captures écran, relevés,\\ logs, témoins)};
\node[box, below=of preuve, align=center] (plainte){2. Dépôt de plainte\\ (Commissariat / Brigade\\ de Gendarmerie / Procureur)};
\node[decision, below=of plainte, align=center] (competence){Compétence\\ territoriale /\\ matière ?};
\node[box, below left=of competence, xshift=-1.5cm, align=center] (enquete){3. Enquête policière\\ (BIS / Gendarmerie)\\ - Saisie supports\\ - Réquisitions opérateurs\\ - Expertise ANTIC si besoin};
\node[box, below right=of competence, xshift=1.5cm, align=center] (classement){Classement sans suite\\ (infraction non\\ déterminée, auteur\\ inconnu, prescription)};
\node[box, below=of enquete, align=center] (poursuite){4. Poursuites\\ (Citation directe /\\ Ouverture information\\ judiciaire par Procureur)};
\node[box, below=of poursuite, align=center] (jugement){5. Jugement\\ (TPI / TGI)\\ - Audience publique\\ - Débats contradictoires};
\node[startstop, below=of jugement, align=center] (sanction){Sanction prononcée\\ (Emprisonnement, Amende,\\ Dommages-intérêts,\\ Mesures complémentaires)};

\draw[arrow] (start) -- (preuve);
\draw[arrow] (preuve) -- (plainte);
\draw[arrow] (plainte) -- (competence);
\draw[arrow] (competence) -- node[left, font=\scriptsize, pos=0.6] {Oui} (enquete);
\draw[arrow] (competence) -- node[right, font=\scriptsize, pos=0.6] {Non} (classement);
\draw[arrow] (enquete) -- (poursuite);
\draw[arrow] (poursuite) -- (jugement);
\draw[arrow] (jugement) -- (sanction);
\draw[arrow, dashed, popmuted] (classement) |- node[above, font=\scriptsize, align=center]{Recours possible\\ devant Procureur Général} ++(-4,0) |- (plainte);
\end{tikzpicture}
\legende{Parcours type d'une plainte pour cyberinfraction au Cameroun : de la victime au jugement. Les flèches pleines indiquent le flux principal ; la flèche pointillée le recours en cas de classement sans suite.}
\end{popfigure}

\begin{methode}[Réagir face à une infraction numérique : les 4 réflexes indispensables]
Si vous êtes victime ou témoin d'une infraction numérique (arnaque, piratage, harcèlement, diffusion illicite) :
\begin{enumerate}
    \item \textbf{Ne touchez à rien / Conservez les preuves :} Faites des captures d'écran (\texttt{Windows+Maj+S} ou \texttt{Cmd+Maj+4}), notez les URLs, horodatages, numéros de téléphone, identifiants de transaction (Mobile Money). Ne supprimez pas les messages, ne formatez pas l'appareil.
    \item \textbf{Portez plainte immédiatement :} Rendez-vous au commissariat ou à la brigade de gendarmerie le plus proche. Exigez un récépissé de dépôt de plainte (PV). Vous pouvez aussi saisir directement le Procureur de la République par lettre recommandée avec AR.
    \item \textbf{Signalez à l'ANTIC si nécessaire :} Pour les atteintes aux systèmes (piratage, virus, défiguration de site), signalez sur \texttt{www.antic.cm} (rubrique « Signalement ») ou par email \texttt{cert@antic.cm}. L'ANTIC peut émettre des alertes, bloquer des domaines \texttt{.cm} malveillants, et assister l'enquête.
    \item \textbf{Suivez la procédure :} Notez le numéro de PV, le nom de l'OPJ en charge. Relancez poliment après 2-3 semaines. En cas de classement sans suite non motivé, saisissez le Procureur Général (recours hiérarchique) ou constituez-vous partie civile devant le Juge d'instruction (article 73 CPP).
\end{enumerate}
\end{methode}

\begin{exemplebox}[Escroquerie Mobile Money : parcours réel]
Mme Ngo, enseignante à Douala, reçoit un SMS : « Votre compte Mobile Money a été bloqué. Appelez le 6XX XX XX XX pour le débloquer. » Elle appelle, le faux agent lui demande son code PIN « pour vérification ». Elle le donne. 5 minutes plus tard, 150 000 FCFA disparaissent.
\begin{enumerate}
    \item \textbf{Preuves :} Elle fait une capture de l'SMS, note le numéro appelant, prend en photo l'historique des transactions (USSD \texttt{*126\#} ou appli) montrant le débit frauduleux.
    \item \textbf{Plainte :} Elle se rend à la Brigade de Gendarmerie de son arrondissement, dépose plainte contre X pour « escroquerie par moyen de communication électronique » (art. 6 loi n°2010/012). Elle remet copies des preuves.
    \item \textbf{Enquête :} L'OPJ réquisitionne l'opérateur (MTN/Orange) pour obtenir l'identité du titulaire du numéro fraudeur, la localisation de l'antenne relais, et le destinataire final des fonds (souvent un autre numéro ou un commerçant).
    \item \textbf{ANTIC :} L'OPJ peut demander l'appui du CERT-CM (ANTIC) pour analyser le téléphone de la victime (vérifier absence de malware) et tracer l'IP si l'arnaque a utilisé un lien web.
    \item \textbf{Jugement :} Si l'auteur est identifié et interpellé, le Procureur poursuit. Le TGI de Douala condamne l'auteur à 2 ans de prison ferme et 500 000 FCFA d'amende (art. 6 : 1-5 ans + amende 500 000-10 000 000 FCFA) + dommages-intérêts pour Mme Ngo.
\end{enumerate}
\end{exemplebox}

\courssub{Organigramme des acteurs nationaux : une vision d'ensemble}

Pour mémoriser les rôles et les liens de tutelle ou de collaboration, l'organigramme suivant positionne chaque institution par rapport aux autres.

\begin{popfigure}
\begin{tikzpicture}[
    node distance=1.5cm and 2cm,
    ministry/.style={rectangle, draw=popblue, thick, fill=popblueL, rounded corners=8pt, minimum width=3.5cm, minimum height=1.3cm, align=center, font=\small\bfseries},
    agency/.style={rectangle, draw=poppurple, thick, fill=poppurpleL, rounded corners=6pt, minimum width=3.2cm, minimum height=1.1cm, align=center, font=\small},
    justice/.style={rectangle, draw=popgreen, thick, fill=popgreenL, rounded corners=6pt, minimum width=3.2cm, minimum height=1.1cm, align=center, font=\small},
    police/.style={rectangle, draw=poporange, thick, fill=poporangeL, rounded corners=6pt, minimum width=3.2cm, minimum height=1.1cm, align=center, font=\small},
    arrow/.style={->, >=Stealth, thick, popdark},
    dasharrow/.style={<->, >=Stealth, dashed, thick, popmuted}
]
\node[ministry, align=center] (minpostel){MINPOSTEL\\ (Politique, Stratégie,\\ Tutelle technique,\\ Représentation int.)};
\node[agency, below left=of minpostel, align=center] (antic){ANTIC\\ (Sécurité SI, CERT-CM,\\ Certification, .cm,\\ Régulation contenus/\\ services TIC)};
\node[agency, below right=of minpostel, align=center] (art){ART\\ (Régulation opérateurs,\\ Fréquences, Numérotation,\\ Qualité service,\\ Concurrence, Consommateurs)};
\node[police, below=of antic, align=center] (police){Police Judiciaire / BIS\\ (Enquêtes cyber,\\ Saisies, Réquisitions,\\ Expertise forensique)};
\node[police, below=of art, align=center] (gendarmerie){Gendarmerie Nationale\\ (Enquêtes cyber zone rurale,\\ Saisies, Réquisitions,\\ Police judiciaire militaire)};
\node[justice, below=of police, align=center] (tribunaux){Tribunaux (TPI/TGI)\\ Cours d'Appel\\ Cour Suprême\\ (Jugement, Sanctions,\\ Dommages-intérêts)};

\draw[arrow] (minpostel) -- node[left, font=\scriptsize] {Tutelle} (antic);
\draw[arrow] (minpostel) -- node[right, font=\scriptsize] {Tutelle} (art);
\draw[dasharrow] (antic) -- node[above, font=\scriptsize, align=center]{Collaboration\\ technique} (art);
\draw[arrow] (antic) -- node[left, font=\scriptsize, align=center]{Appui expertise\\ CERT-CM} (police);
\draw[arrow] (art) -- node[right, font=\scriptsize, align=center]{Données abonnés\\ Réquisitions} (gendarmerie);
\draw[arrow] (police) -- node[left, font=\scriptsize, align=center]{Transmission\\ procès-verbal} (tribunaux);
\draw[arrow] (gendarmerie) -- node[right, font=\scriptsize, align=center]{Transmission\\ procès-verbal} (tribunaux);
\draw[dasharrow, bend left=20] (antic) to node[above, font=\scriptsize, align=center]{Signalement\\ incidents} (tribunaux);
\draw[dasharrow, bend right=20] (art) to node[below, font=\scriptsize, align=center]{Avis technique\\ réglementation} (tribunaux);
\end{tikzpicture}
\legende{Organigramme des acteurs nationaux de la régulation et de la sanction du numérique au Cameroun. Flèches pleines : lien de tutelle ou flux procédural principal. Flèches pointillées : collaboration technique ou avis consultatif.}
\end{popfigure}

\begin{exoresolu}[Identifier l'institution compétente]
\textbf{Énoncé :} Pour chacune des situations suivantes, indiquez l'institution camerounaise \emph{principalement} compétente pour agir en premier lieu (hors justice) et justifiez en une phrase.

\begin{enumerate}
    \item Un cybercafé à Yaoundé utilise un logiciel de gestion piraté (cracké) sur ses 20 postes.
    \item Un abonné MTN constate que son forfait « 1 Go/jour » est épuisé après 200 Mo utilisés ; le service client ne répond pas.
    \item Le site web d'une administration publique (\texttt{www.minxxx.gov.cm}) est défiguré par un hacker qui y affiche un message politique.
    \item Une entreprise souhaite devenir fournisseur de services de certification électronique (PSCE) pour vendre des certificats de signature électronique.
    \item Un particulier reçoit des appels répétés d'un numéro masqué proférant des menaces de mort.
\end{enumerate}

\textbf{Solution détaillée :}
\begin{enumerate}
    \item \textbf{ANTIC.} L'utilisation de logiciels piratés dans un cadre professionnel relève de la sécurité des systèmes d'information et de la propriété intellectuelle logicielle ; l'ANTIC mène des opérations de contrôle et de sensibilisation (art. 34 loi n°2010/012).
    \item \textbf{ART.} La qualité de service, le respect des offres tarifaires et le traitement des réclamations des consommateurs contre les opérateurs relèvent de la régulation des communications électroniques (loi n°2010/013, art. 24-26).
    \item \textbf{ANTIC (CERT-CM).} L'atteinte à un système d'information d'une administration (défiguration, intrusion) est un incident de sécurité majeur ; le CERT-CM de l'ANTIC est le point de contact national pour la réponse à incident et l'alerte.
    \item \textbf{ANTIC.} L'agrément des Prestataires de Services de Certification Électronique (PSCE) est une prérogative exclusive de l'ANTIC (loi n°2010/012, art. 7 et décret d'application).
    \item \textbf{Police / Gendarmerie (BIS / Brigade).} Les menaces de mort, même par téléphone, sont des infractions de droit commun (Code pénal) relevant de la police judiciaire ; l'ART peut être réquisitionnée pour identifier l'appelant (levée d'anonymat), mais l'enquête est policière.
\end{enumerate}
\end{exoresolu}

\begin{exercice}[Mise en situation : la chaîne institutionnelle]
M. Kamga, gérant d'une PME à Bafoussam, découvre que son ancien comptable, parti depuis 3 mois, a copié la base de données clients (fichier Excel + CRM en ligne) et l'utilise pour démarcher les clients au profit d'une entreprise concurrente. M. Kamga a les logs de connexion VPN (adresse IP, horaires) et des témoins.

\begin{enumerate}
    \item Quelles sont les infractions potentiellement constitutives (citez les articles de la loi n°2010/012) ?
    \item Décrivez, dans l'ordre chronologique, les 5 étapes que M. Kamga doit suivre pour faire valoir ses droits, en précisant à chaque étape l'institution saisie et l'action attendue.
    \item Pourquoi est-il crucial de ne \emph{pas} effacer les logs du serveur VPN ni de « nettoyer » le poste de l'ancien comptable avant l'arrivée des enquêteurs ?
\end{enumerate}
\end{exercice}

\begin{corrige}[Exercice : Mise en situation]
\begin{enumerate}
    \item \textbf{Infractions :}
    \begin{itemize}
        \item \textbf{Art. 4 :} Accès frauduleux / maintien dans un système d'information (l'ancien comptable n'avait plus d'habilitation).
        \item \textbf{Art. 5 :} Entrave au fonctionnement d'un système d'information (si suppression/modification de données).
        \item \textbf{Art. 6 :} Escroquerie / abus de confiance par moyen numérique (détournement de fichier clients pour gain personnel).
        \item \textbf{Art. 10 :} Violation du secret des correspondances / données à caractère personnel (fichier clients = données personnelles, loi n°2024/016 aussi applicable).
    \end{itemize}
    \item \textbf{Parcours en 5 étapes :}
    \begin{enumerate}
        \item \textbf{Conservation preuves (M. Kamga) :} Copie sécurisée des logs VPN, export CRM, attestations témoins. \emph{Ne rien modifier.}
        \item \textbf{Plainte (Gendarmerie Bafoussam / Commissariat) :} Dépôt de plainte contre X (ou nommé) pour « accès frauduleux, vol de données, concurrence déloyale numérique ». Remise copies preuves. Obtention récépissé (PV n°...).
        \item \textbf{Enquête (OPJ Gendarmerie + ANTIC) :} OPJ réquisitionne logs FAI/VPN, saisit poste comptable (si encore sur place), demande expertise ANTIC pour horodater les accès et certifier l'intégrité des preuves numériques (chaîne de custody).
        \item \textbf{Poursuites (Procureur TGI Bafoussam) :} Procureur ouvre information judiciaire ou cite direct. Juge d'instruction nommé si complexité. M. Kamga se constitue partie civile (demande dommages-intérêts).
        \item \textbf{Jugement (TGI Bafoussam) :} Audience, débats, jugement. Peines encourues : jusqu'à 5 ans prison + 10 000 000 FCFA amende (art. 4, 6) + dommages-intérêts civils.
    \end{enumerate}
    \item \textbf{Intégrité de la preuve numérique :} Effacer les logs ou « nettoyer » le poste détruit la \emph{chaîne de custody} (traçabilité ininterrompue de la preuve). L'expert ANTIC ne pourra plus certifier l'authenticité, l'horodatage et l'origine des accès. La défense de l'accusé arguera que les preuves ont été fabriquées \emph{a posteriori} $\rightarrow$ irrecevabilité des preuves $\rightarrow$ relaxe. \textbf{Règle d'or :} On fige l'état du système (image disque, hash SHA-256) avant toute manipulation.
\end{enumerate}
\end{corrige}

\begin{aretenir}[Synthèse : qui fait quoi ?]
\begin{center}
\begin{tabularx}{\linewidth}{|l|X|X|}
\hline
\rowcolor{popblueL}
\textbf{Institution} & \textbf{Mission cœur de métier} & \textbf{Exemple d'intervention type} \\
\hline
MINPOSTEL & Politique, lois, stratégie, tutelle & Préparation loi protection données ; représentation à l'UIT \\
\hline
ANTIC & Sécurité SI, certification, .cm, CERT & Audit sécurité ministère ; blocage site .cm malveillant ; alerte rançongiciel \\
\hline
ART & Opérateurs, fréquences, qualité, conso & Sanction MTN/Orange pour débit non tenu ; attribution licence 5G \\
\hline
Police / Gendarmerie & Enquête judiciaire, saisies, réquisitions & Interpellation arnaqueur Mobile Money ; saisie ordinateur pirate \\
\hline
Tribunaux (Parquet + Siège) & Poursuites, jugement, sanctions & Condamnation pirate à 3 ans ferme + 2 M FCFA amende + dommages \\
\hline
\end{tabularx}
\end{center}
\end{aretenir}

Cette architecture institutionnelle, bien que perfectible (manque de juges spécialisés, délais de procédure, ressources limitées des unités cyber), existe bel et bien et fonctionne. En tant que futurs citoyens numériques, développeurs ou gestionnaires de systèmes, vous devez savoir \emph{vers qui vous tourner} selon la nature du problème : \textbf{ART} pour la facture ou la couverture réseau, \textbf{ANTIC} pour le certificat expiré, le site piraté ou le nom de domaine volé, \textbf{Police/Gendarmerie} pour le délit pénal, \textbf{Justice} pour la réparation finale. La section suivante examinera les infractions précises et leurs sanctions, pour compléter ce tableau.

\coursec{Infractions numériques et sanctions au niveau national}

Dans la section précédente, nous avons découvert les institutions nationales chargées de réguler et de sanctionner les usages du numérique au Cameroun : l'ANTIC, le CONSUPE, les tribunaux et les forces de l'ordre spécialisées. Mais une question demeure : \emph{que risque concrètement} celui qui pirate un compte, escroque des internautes ou publie une fausse information sur WhatsApp ? C'est précisément l'objet de cette section : dresser le tableau des \cle{infractions numériques} et des \cle{sanctions} prévues par la loi camerounaise, en particulier la \textbf{loi n° 2010/012 du 21 décembre 2010} relative à la cybersécurité et à la cybercriminalité, complétée par le \textbf{code pénal camerounais} (loi n° 2016/007 du 12 juillet 2016).

L'intuition est simple : la loi fonctionne comme un « règlement du jeu » du monde numérique. Chaque mauvaise action en ligne correspond à une infraction précise, et chaque infraction correspond à une peine précise. Connaître cette correspondance, c'est savoir se protéger soi-même et éviter de commettre une infraction par ignorance.

\courssub{Accès frauduleux à un système d'information : le piratage}

\textbf{Intuition.} Imagine que quelqu'un force la serrure de la porte de ta maison et s'y installe sans ton autorisation. C'est une violation de domicile. Dans le monde numérique, le « domicile » est le \cle{système d'information} (un ordinateur, un serveur, un compte Facebook, la plateforme d'une banque), et « forcer la serrure » consiste à y accéder sans droit, par exemple en devinant ou en volant un mot de passe.

\begin{definition}[Accès frauduleux et maintien frauduleux]
L'\cle{accès frauduleux} est le fait de pénétrer, sans autorisation, dans tout ou partie d'un système d'information. Le \cle{maintien frauduleux} est le fait de rester dans ce système alors qu'on n'y a pas le droit, même si l'on y est entré par erreur ou avec une autorisation initiale qui a pris fin.
\end{definition}

\begin{exemplebox}[Exemple chiffré : que risque un pirate ?]
Selon l'article 7 de la loi n° 2010/012, l'accès frauduleux ou le maintien frauduleux dans un système d'information est puni d'un \textbf{emprisonnement de six mois à deux ans} et d'une \textbf{amende de 500\,000 à 5\,000\,000 de FCFA} (ou de l'une de ces deux peines seulement). Si le piratage s'accompagne d'une suppression ou d'une modification de données (art. 9), les peines sont \emph{alourdies} : l'emprisonnement est de \textbf{un à cinq ans} et l'amende de \textbf{1\,000\,000 à 10\,000\,000 de FCFA}. Autrement dit, « juste regarder » est déjà punissable ; « toucher aux données » aggrave considérablement la peine.
\end{exemplebox}

\begin{attention}[Piège fréquent : « je n'ai rien volé »]
Beaucoup d'élèves pensent qu'entrer dans le compte d'un camarade « juste pour voir » n'est pas grave. \textbf{C'est faux} : la loi punit l'accès lui-même, même sans vol, sans dégât et sans publication. La simple intrusion non autorisée constitue déjà l'infraction.
\end{attention}

\courssub{Atteintes aux données : interception, modification, destruction}

\textbf{Intuition.} Les données sont les « biens » du monde numérique : notes d'élèves, dossiers médicaux, comptes bancaires, photos. Trois attaques principales peuvent les viser :

\begin{itemize}
  \item l'\cle{interception} : capter des données pendant leur transmission (écouter une communication, récupérer un mot de passe qui circule sur un réseau non sécurisé) ;
  \item la \cle{modification} : altérer frauduleusement des données (changer une note dans le système d'un lycée, modifier un montant sur une facture électronique) ;
  \item la \cle{destruction} : supprimer ou endommager des données (effacer la base de données d'une entreprise, introduire un virus destructeur).
\end{itemize}

\begin{propriete}[Sanctions des atteintes aux données (Art. 9, 10, 11 de la loi 2010/012)]
La loi de 2010 punit l'interception, la modification et la destruction frauduleuses de données d'un \textbf{emprisonnement de un à cinq ans} et d'une \textbf{amende de 1\,000\,000 à 10\,000\,000 de FCFA}. Lorsque l'atteinte vise des données de l'État ou des données personnelles sensibles, les peines maximales sont appliquées plus sévèrement.
\end{propriete}

\courssub{Escroquerie et fraude en ligne : arnaques, faux sites, phishing}

\textbf{Intuition.} L'escroquerie est un mensonge destiné à soutirer de l'argent ou un bien. Sur Internet, elle prend des formes nouvelles : faux sites de vente qui encaissent sans livrer, messages annonçant un faux gain à un concours, ou encore le très répandu \cle{phishing}.

\begin{definition}[Phishing (hameçonnage)]
Le \cle{phishing} ou \cle{hameçonnage} est une technique frauduleuse qui consiste à soutirer des informations personnelles (mots de passe, numéros de carte bancaire, codes de mobile money) en se faisant passer pour un organisme de confiance : une banque, Orange Money, MTN Mobile Money, une administration ou un réseau social. La victime reçoit un message crédible contenant un lien vers un \emph{faux site} imitant le vrai, où elle saisit elle-même ses informations.
\end{definition}

\begin{exemplebox}[Une arnaque typique au Cameroun]
Amina reçoit un SMS : « Votre compte OM sera bloqué. Confirmez votre code secret sur om-verification.com ». Le site ressemble trait pour trait au site officiel. En saisissant son code, elle le transmet en réalité à un escroc qui vide son compte. L'escroc commet une \textbf{escroquerie informatique}, punie par le code pénal (art. 320) et la loi de 2010 d'un \textbf{emprisonnement pouvant atteindre dix ans} et d'une \textbf{amende pouvant aller jusqu'à 10\,000\,000 de FCFA}, voire plus si le préjudice est important ou en cas de récidive.
\end{exemplebox}

\courssub{Diffamation, injures et propagation de fausses nouvelles}

\textbf{Intuition.} Dire publiquement du mal de quelqu'un ou répandre un mensonge n'est pas devenu légal parce qu'on le fait derrière un écran. Au contraire : les réseaux sociaux amplifient la portée du propos, donc sa gravité.

\begin{definition}[Diffamation, injure, fausse nouvelle]
\begin{itemize}
  \item La \cle{diffamation} est l'accusation publique d'un fait précis qui porte atteinte à l'honneur d'une personne (ex. : « ce commerçant vole ses clients »).
  \item L'\cle{injure} est une expression outrageante, un terme de mépris ou une invective qui ne renferme l'imputation d'aucun fait précis (insultes publiques).
  \item La \cle{propagation de fausses nouvelles} est la diffusion, par quelque moyen que ce soit, d'informations mensongères, de bruits ou de rumeurs susceptibles de troubler l'ordre public (fausse alerte, rumeur de catastrophe, fausse annonce de décès).
\end{itemize}
\end{definition}

Ces infractions, commises par voie électronique (Facebook, WhatsApp, TikTok), sont punies par la loi de 2010 (art. 31, 32) et le code pénal d'\textbf{emprisonnement de six mois à deux ans} et d'\textbf{amendes de 500\,000 à 5\,000\,000 de FCFA}.

\begin{attention}[Relayer, c'est déjà risquer]
Partager une fausse information \emph{même sans l'avoir créée} peut engager ta responsabilité pénale. Le simple « transfert » d'une rumeur sur WhatsApp fait de toi un \textbf{relais} de la diffamation ou de la fausse nouvelle. L'article 31 de la loi 2010/012 punit quiconque « propage » de fausses nouvelles. Avant de partager : vérifie la source, la date et la crédibilité de l'information.
\end{attention}

\courssub{Atteintes à la vie privée et infractions aggravées}

\textbf{Intuition.} Ta vie privée — tes photos, tes messages, tes données personnelles — t'appartient. La publier sans ton accord, ou pire, harceler quelqu'un en ligne, sont des infractions sérieuses, encore plus sévèrement punies lorsque la victime est un mineur.

\begin{itemize}
  \item \textbf{Publication non autorisée d'images ou de données personnelles} (Art. 30 loi 2010/012) : diffuser l'image ou les données d'une personne sans son consentement est puni d'emprisonnement de six mois à deux ans et d'amende de 500\,000 à 5\,000\,000 de FCFA.
  \item \textbf{Cyberharcèlement} : le fait de harceler une personne de manière répétée par voie électronique (messages injurieux, menaces, humiliations publiques) est punissable (Code pénal, art. 241-1), avec des peines \emph{aggravées} (jusqu'à 5 ans de prison) lorsque la victime est un mineur ou une personne vulnérable.
  \item \textbf{Pornographie impliquant des mineurs} (Code pénal, art. 294-1 et loi 2010/012 art. 34) : détenir, produire, offrir, mettre à disposition ou diffuser de tels contenus est un \textbf{crime} très grave, puni de \textbf{réclusion criminelle de 10 à 20 ans} et d'une amende de 5\,000\,000 à 50\,000\,000 de FCFA, pouvant aller jusqu'à la \textbf{réclusion à perpétuité} si la victime est âgée de moins de 15 ans ou en cas de récidive.
\end{itemize}

\begin{propriete}[Échelle générale des sanctions au Cameroun]
Les sanctions pour cyberinfractions s'échelonnent ainsi :
\begin{itemize}
  \item \textbf{Amendes} : de 500\,000 FCFA (infractions simples : injure, accès simple) à 50\,000\,000 FCFA (crimes graves) ;
  \item \textbf{Emprisonnement / Réclusion} : de six mois (infractions simples) à 20 ans (crimes), voire la perpétuité pour les cas les plus graves (pornographie juvénile aggravée, cyberterrorisme).
\end{itemize}
Les peines d'amende et d'emprisonnement sont le plus souvent \emph{cumulées} : le juge peut prononcer les deux simultanément.
\end{propriete}

\begin{propriete}[Circonstances aggravantes (Art. 42 loi 2010/012 et Code pénal)]
Certaines situations \emph{alourdissent} automatiquement la peine encourue :
\begin{enumerate}
  \item L'infraction est commise contre un \textbf{système d'information de l'État} ou d'une administration (site d'un ministère, plateforme des impôts, système hospitalier) ;
  \item La victime est un \textbf{mineur} (moins de 18 ans) ou une personne vulnérable ;
  \item L'infraction est commise \textbf{en bande organisée} (plusieurs personnes coordonnées) ;
  \item L'auteur est un \textbf{professionnel} (informaticien, agent de banque, fonctionnaire) qui a abusé de ses fonctions ou de son accès légitime.
\end{enumerate}
\end{propriete}

\courssub{Tableau récapitulatif et comparaison des amendes}

\begin{popfigure}
\begin{center}
\begin{tabularx}{\linewidth}{|>{\raggedright\arraybackslash}X|>{\raggedright\arraybackslash}X|}
\hline
\rowcolor{popblue}
\textbf{\color{white}Infraction} & \textbf{\color{white}Sanction encourue (Loi 2010/012 / Code pénal)} \\
\hline
Accès frauduleux à un système d'information (piratage simple) & 6 mois à 2 ans de prison ; amende 500\,000 à 5\,000\,000 FCFA \\
\hline
Modification ou destruction de données & 1 à 5 ans de prison ; amende 1\,000\,000 à 10\,000\,000 FCFA \\
\hline
Escroquerie en ligne, phishing & Jusqu'à 10 ans de prison ; amende jusqu'à 10\,000\,000 FCFA (ou plus selon préjudice) \\
\hline
Diffamation, injures, fausses nouvelles sur les réseaux sociaux & 6 mois à 2 ans de prison ; amende 500\,000 à 5\,000\,000 FCFA \\
\hline
Pornographie impliquant des mineurs (détention/diffusion) & 10 à 20 ans de réclusion (perpétuité si < 15 ans) ; amende 5\,000\,000 à 50\,000\,000 FCFA \\
\hline
\end{tabularx}
\end{center}
\legende{Correspondance entre cinq infractions numériques types et les sanctions encourues au Cameroun (montants et durées indicatifs selon les textes en vigueur).}
\end{popfigure}

\begin{popfigure}
\begin{center}
\begin{tikzpicture}
\begin{axis}[
  ybar, bar width=0.6cm,
  width=0.95\linewidth, height=7cm,
  ymin=0, ymax=55,
  ylabel={Amende maximale indicative (millions de FCFA)},
  symbolic x coords={Injure/Diffamation, Piratage simple, Destruction de donnees, Escroquerie, Pedopornographie},
  xtick=data,
  x tick label style={rotate=25, anchor=east, font=\footnotesize, align=center},
  ytick={0,10,20,30,40,50},
  nodes near coords,
  every node near coord/.append style={font=\footnotesize, font=\bfseries},
  axis line style={draw=popline},
  ymajorgrids=true, grid style={popline!40},
  enlarge x limits=0.25,
]
\addplot[fill=popblue, draw=popdark] coordinates {
  (Injure/Diffamation, 5)
  (Piratage simple, 5)
  (Destruction de donnees, 10)
  (Escroquerie, 10)
  (Pedopornographie, 50)
};
\end{axis}
\end{tikzpicture}
\end{center}
\legende{Ordres de grandeur des amendes maximales selon la gravité de l'infraction : la sanction financière croît avec la gravité de l'atteinte.}
\end{popfigure}

On observe une logique claire : la sanction croît avec la \textbf{gravité de l'atteinte} (aux personnes, aux biens numériques, à l'État) et avec le \textbf{préjudice} causé.

\courssub{Analyser un cas concret : la démarche en quatre étapes}

Face à une situation de la vie courante (un exercice au BEPC, un fait réel), il faut raisonner comme un juriste, méthodiquement.

\begin{methode}[Analyser un cas de cyberinfraction en 4 étapes]
\begin{enumerate}
  \item \textbf{Identifier les faits} : qui a fait quoi, par quel moyen numérique, avec quelles conséquences ?
  \item \textbf{Qualifier l'infraction} : à quelle catégorie correspondent les faits (accès frauduleux, escroquerie, diffamation, atteinte à la vie privée...) ?
  \item \textbf{Citer le texte applicable} : loi n° 2010/012 du 21 décembre 2010 relative à la cybersécurité et à la cybercriminalité, et/ou code pénal camerounais.
  \item \textbf{Indiquer la sanction encourue} : nature de la peine (amende et/ou emprisonnement), ordre de grandeur, et circonstances aggravantes éventuelles.
\end{enumerate}
\end{methode}

\begin{exoresolu}[Le compte piraté de la boutique de M. Etoundi]
\textbf{Énoncé.} Junior, 17 ans, devine le mot de passe du compte WhatsApp Business de la boutique de M. Etoundi. Il y reste connecté plusieurs jours, lit les messages des clients, puis envoie à certains d'entre eux un faux numéro de mobile money pour « payer leurs commandes ». Deux clients versent au total \num{150000} FCFA sur ce numéro. Analyse cette situation en suivant la démarche en quatre étapes.
\tcblower
\textbf{Solution détaillée.}
\begin{enumerate}
  \item \textbf{Faits} : Junior accède sans autorisation au compte professionnel de M. Etoundi (système d'information), s'y maintient plusieurs jours (maintien frauduleux), lit des messages (interception/violation de correspondance), puis détourne \num{150000} FCFA en trompant deux clients via un faux numéro (escroquerie).
  \item \textbf{Qualification} : on relève (a) un \emph{accès frauduleux} et un \emph{maintien frauduleux} dans un système d'information (Art. 7, 8 loi 2010/012) ; (b) une \emph{escroquerie en ligne} (Art. 320 Code pénal, Art. 11 loi 2010/012) ; (c) potentiellement une \emph{violation du secret des correspondances} (Art. 29 loi 2010/012).
  \item \textbf{Textes applicables} : la loi n° 2010/012 du 21 décembre 2010 relative à la cybersécurité et à la cybercriminalité (accès, maintien, escroquerie, secret des correspondances) et le code pénal camerounais (escroquerie, recel).
  \item \textbf{Sanctions encourues} :
  \begin{itemize}
    \item Accès/Maintien frauduleux : 6 mois à 2 ans de prison + amende 500\,000 à 5\,000\,000 FCFA.
    \item Escroquerie : jusqu'à 10 ans de prison + amende (jusqu'à 10\,000\,000 FCFA ou plus).
    \item Les peines peuvent se \emph{cumuler} (concours réel d'infractions).
  \end{itemize}
  Junior étant mineur (17 ans), il relève de la \textbf{justice des mineurs} (Ordonnance 1962 modifiée) : le juge pour enfants peut prononcer des mesures éducatives ou des peines atténuées (peine maximale divisée par deux en principe), mais sa responsabilité pénale reste engagée. Ses parents (ou tuteurs) sont civilement responsables du préjudice (Art. 1384 Code civil) et devront rembourser les \num{150000} FCFA aux victimes.
\end{enumerate}
\textbf{Conclusion} : une « blague » en ligne peut se transformer en plusieurs infractions cumulées aux conséquences très lourdes, même pour un mineur.
\end{exoresolu}

\begin{savaistu}[Même la tentative est punissable]
Au Cameroun, l'article 44 de la loi n° 2010/012 dispose que « la tentative des infractions prévues par la présente loi est punie des mêmes peines ». La simple \textbf{tentative} de piratage, d'escroquerie ou de destruction de données est donc punissable, même si elle échoue. Un pirate qui essaie de pénétrer le système d'une banque sans y parvenir peut quand même être poursuivi et condamné. Mieux vaut donc ne jamais « essayer pour voir » !
\end{savaistu}

\begin{aretenir}[L'essentiel de la section]
\begin{itemize}
  \item La loi n° 2010/012 du 21 décembre 2010 et le code pénal punissent les principales cyberinfractions : piratage (accès/maintien frauduleux), atteintes aux données (interception, modification, destruction), escroquerie et phishing, diffamation/injures/fausses nouvelles, atteintes à la vie privée, cyberharcèlement, pornographie juvénile.
  \item Les sanctions vont de 6 mois de prison et 500\,000 FCFA d'amende (infractions simples) jusqu'à 20 ans de réclusion (voire perpétuité) et 50\,000\,000 FCFA d'amende pour les crimes les plus graves.
  \item Les circonstances aggravantes (système de l'État, victime mineure, bande organisée, abus de fonction) alourdissent les peines.
  \item Relayer une fausse information ou tenter une infraction, même sans succès, engage la responsabilité pénale.
  \item Méthode d'analyse d'un cas : \textbf{Faits $\rightarrow$ Qualification $\rightarrow$ Texte applicable $\rightarrow$ Sanction encourue}.
\end{itemize}
\end{aretenir}

\begin{exercice}[À toi de juger]
Dans un groupe WhatsApp de classe, Sandra publie la photo d'une camarade en la traitant de « voleuse », sans aucune preuve. Le message est transféré par dix autres élèves.
\begin{enumerate}
  \item Identifie les infractions commises par Sandra.
  \item Les élèves qui ont transféré le message risquent-ils quelque chose ? Justifie.
  \item Cite le texte applicable et indique les sanctions encourues.
\end{enumerate}
\end{exercice}

\begin{corrige}[Exercice : À toi de juger]
\begin{enumerate}
  \item Sandra commet une \textbf{diffamation} (allégation d'un fait précis — le vol — portant atteinte à l'honneur, sans preuve, Art. 31 loi 2010/012) et une \textbf{atteinte à la vie privée} par publication non autorisée de l'image de sa camarade (Art. 30 loi 2010/012). Si les propos sont simplement insultants sans fait précis, c'est de l'\textbf{injure} (Art. 31).
  \item \textbf{Oui} : les dix élèves sont des \textbf{relais} (coauteurs ou complices) de la diffamation/injure. L'article 31 punit quiconque « propage » l'information diffamatoire. Partager (transférer) une fausse information, même sans l'avoir créée, constitue une propagation et engage leur responsabilité pénale.
  \item \textbf{Textes} : Loi n° 2010/012 du 21 décembre 2010 (Art. 30, 31) et Code pénal. \textbf{Sanctions} : Diffamation/Injure en ligne : 6 mois à 2 ans d'emprisonnement et amende de 500\,000 à 5\,000\,000 FCFA. Atteinte à la vie privée (image) : 6 mois à 2 ans d'emprisonnement et amende de 500\,000 à 5\,000\,000 FCFA. La victime peut aussi demander des \textbf{dommages-intérêts} au civil pour le préjudice moral subi.
\end{enumerate}
\end{corrige}

\coursec{Cadre juridique international de l'éthique numérique}

Dans la section précédente, nous avons vu que le Cameroun dispose de lois nationales (notamment la loi n° 2010/012 du 21 décembre 2010 relative à la cybersécurité et à la cybercriminalité) qui définissent les infractions numériques et leurs sanctions sur le territoire camerounais. Mais une question se pose immédiatement : que se passe-t-il lorsque le criminel se trouve à Douala, que le serveur attaqué est hébergé en France et que la victime habite au Canada ? Aucune loi nationale, à elle seule, ne peut résoudre ce problème. C'est précisément pour cela qu'existe un \cle{cadre juridique international} de l'éthique numérique, que nous allons étudier dans cette section.

\courssub{Pourquoi des règles internationales ?}

\begin{exemplebox}[Une situation qui dépasse les frontières]
Imaginons la situation suivante. Un cybercriminel installé à Yaoundé crée un faux site de banque. Il envoie des messages frauduleux (hameçonnage) à des clients d'une banque américaine. Les informations bancaires volées transitent par un serveur loué en Allemagne. L'argent détourné est ensuite transféré sur un compte dans un troisième pays.
\par
Questions : quelle police peut enquêter ? Quelle loi s'applique ? Quel tribunal peut juger ? La police camerounaise ne peut pas aller saisir un serveur en Allemagne, et la police américaine ne peut pas venir arrêter quelqu'un à Yaoundé sans autorisation.
\end{exemplebox}

Cet exemple montre l'idée fondamentale : \textbf{la cybercriminalité ignore les frontières}. Sur Internet, un criminel, un serveur et une victime peuvent se trouver dans trois pays différents. Or la police et la justice d'un État n'ont de pouvoir que sur leur propre territoire : c'est le principe de \cle{souveraineté nationale}. Sans règles internationales, il suffirait à un criminel d'agir depuis un pays étranger pour échapper à toute sanction.

\begin{definition}[Traité international]
Un \cle{traité international} (ou \cle{convention internationale}) est un accord écrit conclu entre des États. Il est \textbf{contraignant} (obligatoire) pour les États qui l'ont \cle{ratifié}, c'est-à-dire approuvé officiellement selon leur propre procédure nationale. Un État qui n'a pas ratifié un traité n'est pas tenu de l'appliquer.
\end{definition}

\begin{aretenir}[L'idée à retenir]
Les règles internationales du numérique servent à trois choses :
\begin{enumerate}
\item \textbf{Harmoniser les lois} : que les mêmes actes (piratage, fraude en ligne, vol de données) soient reconnus comme des infractions dans un maximum de pays ;
\item \textbf{Organiser la coopération} : permettre aux polices et aux justices de différents pays de s'entraider (échange d'informations, saisie de données, arrestations) ;
\item \textbf{Protéger les personnes partout} : garantir un niveau minimal de protection des données et des droits, quel que soit le pays.
\end{enumerate}
\end{aretenir}

\begin{popfigure}
\centering
\begin{tikzpicture}[font=\small, >=Stealth]
% Zones géographiques stylisées
\draw[fill=popgreenL, rounded corners=6pt] (-5.2,-2.6) rectangle (-1.6,2.6);
\draw[fill=popblueL, rounded corners=6pt] (-1.0,-2.6) rectangle (2.6,2.6);
\draw[fill=poporangeL, rounded corners=6pt] (3.2,-2.6) rectangle (6.8,2.6);
\node[font=\small\bfseries, text=popdark] at (-3.4,2.2) {Afrique};
\node[font=\small\bfseries, text=popdark] at (0.8,2.2) {Europe};
\node[font=\small\bfseries, text=popdark] at (5.0,2.2) {Amérique};
% Acteurs
\node[draw=popdark, fill=white, rounded corners=3pt, align=center, inner sep=4pt] (crim) at (-3.4,0.6) {\textbf{Criminel}\\ (Yaoundé)};
\node[draw=popdark, fill=white, rounded corners=3pt, align=center, inner sep=4pt] (serv) at (0.8,0.6) {\textbf{Serveur}\\ (Allemagne)};
\node[draw=popdark, fill=white, rounded corners=3pt, align=center, inner sep=4pt] (vict) at (5.0,0.6) {\textbf{Victime}\\ (Canada)};
% Flux de l'attaque
\draw[->, very thick, poppink] (crim) -- node[above, font=\footnotesize, text=poppink]{attaque} (serv);
\draw[->, very thick, poppink] (serv) -- node[above, font=\footnotesize, text=poppink]{données volées} (vict);
% Flux de coopération
\node[draw=popblue, fill=white, rounded corners=3pt, align=center, inner sep=4pt, font=\footnotesize] (pol1) at (-3.4,-1.4) {Police\\ Cameroun};
\node[draw=popblue, fill=white, rounded corners=3pt, align=center, inner sep=4pt, font=\footnotesize] (pol2) at (0.8,-1.4) {Europol /\\ Police allemande};
\node[draw=popblue, fill=white, rounded corners=3pt, align=center, inner sep=4pt, font=\footnotesize] (pol3) at (5.0,-1.4) {Interpol /\\ Police canadienne};
\draw[<->, thick, popblue, dashed] (pol1) -- node[below, font=\footnotesize, text=popblue]{coopération} (pol2);
\draw[<->, thick, popblue, dashed] (pol2) -- node[below, font=\footnotesize, text=popblue]{entraide} (pol3);
\draw[<->, thick, popblue, dashed] (pol1.south) to[bend right=25] node[below, font=\footnotesize, text=popblue]{demande d'extradition} (pol3.south);
\end{tikzpicture}
\legende{Schéma d'une attaque transfrontalière : le criminel, le serveur et la victime sont dans trois pays différents. En rose, le flux de l'attaque ; en bleu pointillé, les flux de coopération policière et judiciaire entre les États.}
\end{popfigure}

\courssub{La Convention de Budapest (2001) : le premier traité mondial contre la cybercriminalité}

Face à la montée de la cybercriminalité à la fin des années 1990, les États ont compris qu'il fallait un texte commun. C'est ainsi qu'est née la \cle{Convention de Budapest}, signée en 2001 à Budapest (Hongrie) sous l'égide du \textbf{Conseil de l'Europe} (une organisation européenne de défense des droits humains, différente de l'Union européenne).

\begin{definition}[Convention de Budapest]
La \cle{Convention de Budapest} (2001) est le \textbf{premier traité international} consacré à la lutte contre la cybercriminalité. Elle demande aux États qui l'ont ratifiée de :
\begin{itemize}
\item punir dans leur loi nationale les mêmes infractions : accès illégal à un système informatique (piratage), interception illégale de données, atteinte aux données et aux systèmes, fraude informatique, etc. ;
\item coopérer entre eux : échange rapide d'informations, préservation et saisie des données stockées à l'étranger, entraide pour les enquêtes.
\end{itemize}
\end{definition}

\begin{propriete}[Principe de coopération de la Convention de Budapest]
Lorsqu'une infraction numérique implique plusieurs pays, chaque État partie à la Convention doit pouvoir :
\begin{enumerate}
\item recevoir une \textbf{demande d'entraide} d'un autre État (par exemple : « conservez les données de ce serveur avant qu'elles ne soient effacées ») ;
\item agir rapidement, car les données numériques disparaissent vite ;
\item transmettre les preuves recueillies pour le procès.
\end{enumerate}
Condition d'application : l'État sollicité doit avoir ratifié la Convention.
\end{propriete}

\begin{savaistu}[Un traité ouvert au monde entier]
Bien qu'elle ait été préparée par le Conseil de l'Europe, la Convention de Budapest est \textbf{ouverte à tous les pays du monde}, et pas seulement aux pays européens. Plusieurs pays africains y ont adhéré, comme le Sénégal, le Ghana, le Cap-Vert ou encore Maurice. C'est aujourd'hui le texte de référence le plus utilisé au monde pour la coopération contre la cybercriminalité : plus de soixante-dix États y sont parties.
\end{savaistu}

\courssub{Le RGPD (2018) : la protection des données personnelles en Europe}

Les entreprises du numérique (réseaux sociaux, boutiques en ligne, applications mobiles) collectent d'énormes quantités de \textbf{données personnelles} : nom, âge, photos, localisation, habitudes d'achat... Ces données peuvent être revendues ou utilisées de manière abusive. Pour protéger les citoyens européens, l'Union européenne a adopté le \cle{RGPD}.

\begin{definition}[RGPD]
Le \cle{RGPD} (Règlement Général sur la Protection des Données) est un règlement de l'Union européenne entré en application en \textbf{2018}. Il impose à toute entreprise ou organisation qui traite des données personnelles de personnes se trouvant en Europe de :
\begin{itemize}
\item obtenir le \textbf{consentement} clair de la personne avant de collecter ses données ;
\item expliquer clairement à quoi servent les données collectées ;
\item garantir le \textbf{droit d'accès}, de \textbf{rectification} et d'\textbf{effacement} (droit à l'oubli) ;
\item sécuriser les données et signaler les fuites de données.
\end{itemize}
En cas de violation, les amendes peuvent atteindre \textbf{20 millions d'euros ou 4 \% du chiffre d'affaires mondial annuel} de l'entreprise, le montant le plus élevé étant retenu.
\end{definition}

\begin{exemplebox}[Un exemple chiffré : des amendes records]
Une grande entreprise technologique américaine a été condamnée en Europe, en 2023, à une amende d'environ \textbf{1,2 milliard d'euros} pour avoir transféré des données personnelles d'utilisateurs européens vers les États-Unis sans garanties suffisantes, en violation du RGPD. D'autres entreprises (réseaux sociaux, compagnies aériennes, chaînes d'hôtels) ont reçu des amendes de plusieurs dizaines ou centaines de millions d'euros. Ces montants montrent que la protection des données n'est pas un simple conseil : c'est une obligation légale assortie de sanctions financières très lourdes.
\end{exemplebox}

\begin{attention}[Le RGPD peut concerner le Cameroun]
Le RGPD s'applique même aux entreprises situées \textbf{hors de l'Europe} dès lors qu'elles traitent des données de personnes se trouvant en Europe. Une application développée à Douala qui collecte les données d'utilisateurs français doit respecter le RGPD. Le numérique rend les lois « voyageuses » !
\end{attention}

\courssub{La Convention de Malabo (2014) : la réponse de l'Afrique}

L'Afrique a également voulu se doter de son propre texte, adapté à ses réalités. En 2014, l'\textbf{Union africaine} (organisation qui regroupe les pays africains) a adopté à Malabo (Guinée équatoriale) la \cle{Convention de l'Union africaine sur la cybersécurité et la protection des données personnelles}, appelée \textbf{Convention de Malabo}.

\begin{definition}[Convention de Malabo]
La \cle{Convention de Malabo} (2014) est un traité de l'Union africaine qui fixe des règles communes aux pays africains dans trois domaines :
\begin{itemize}
\item les \textbf{transactions électroniques} (commerce en ligne, signatures électroniques) ;
\item la \textbf{protection des données personnelles} (chaque pays doit créer une autorité de protection des données) ;
\item la \textbf{cybersécurité et la cybercriminalité} (définir les infractions et organiser la coopération entre pays africains).
\end{itemize}
Elle n'entre en vigueur dans un pays qu'après sa ratification par ce pays.
\end{definition}

\begin{aretenir}
La Convention de Malabo est le « pendant africain » de la Convention de Budapest : elle montre que l'Afrique, dont le Cameroun, participe activement à la construction du droit international du numérique.
\end{aretenir}

\courssub{Les organisations de coopération : Interpol, Europol, l'ONU et l'UIT}

Les traités fixent des règles, mais il faut aussi des \textbf{organisations} pour les mettre en œuvre au quotidien.

\begin{definition}[Interpol et Europol]
\begin{itemize}
\item \cle{Interpol} (Organisation internationale de police criminelle) est une organisation qui relie les polices de près de 200 pays. Elle ne fait pas d'arrestations elle-même : elle \textbf{facilite l'échange d'informations} entre polices nationales, diffuse des notices de recherche de criminels et coordonne des opérations internationales contre la cybercriminalité.
\item \cle{Europol} est l'agence de police de l'Union européenne. Elle aide les polices des pays européens à coopérer, notamment contre les réseaux de cybercriminels (centres d'appel frauduleux, rançongiciels, pédocriminalité en ligne).
\end{itemize}
\end{definition}

\begin{definition}[ONU et UIT]
\begin{itemize}
\item L'\cle{ONU} (Organisation des Nations Unies) adopte des \textbf{recommandations} et organise des négociations entre tous les pays du monde pour renforcer la sécurité numérique et lutter contre la cybercriminalité.
\item L'\cle{UIT} (Union Internationale des Télécommunications) est une agence spécialisée de l'ONU, chargée des télécommunications et des réseaux. Elle publie des normes techniques, aide les pays en développement (dont les pays africains) à renforcer leur cybersécurité et publie un classement mondial de l'engagement des États en matière de cybersécurité.
\end{itemize}
\end{definition}

\begin{attention}[Recommandation n'est pas loi]
Les recommandations de l'ONU et de l'UIT ne sont \textbf{pas contraignantes} comme un traité ratifié : ce sont des orientations que les États sont invités à suivre. C'est une différence importante à ne pas confondre à l'examen.
\end{attention}

\courssub{Extradition et entraide judiciaire}

Dernier point essentiel : comment un criminel qui a fui à l'étranger peut-il être jugé ?

\begin{definition}[Extradition et entraide judiciaire]
\begin{itemize}
\item L'\cle{extradition} est la remise, par un État à un autre État, d'une personne recherchée pour y être jugée ou y purger une peine. Elle nécessite en général un \textbf{traité d'extradition} entre les deux pays et le respect de conditions (l'acte doit être un délit dans les deux pays).
\item L'\cle{entraide judiciaire} est l'ensemble des services que les justices de deux pays se rendent : transmission de preuves, auditions par visioconférence, saisies de données, exécution de demandes d'enquête.
\end{itemize}
\end{definition}

\begin{attention}[Agir depuis le Cameroun ne protège pas]
Une infraction commise \textbf{depuis le Cameroun} contre une victime située à l'étranger peut entraîner des \textbf{poursuites internationales} et une demande d'\textbf{extradition}. Un élève ou un adulte qui pirate un compte étranger ou escroque une victime européenne depuis un cybercafé de Bafoussam peut être poursuivi à la fois au Cameroun et à l'étranger. Internet n'est pas une zone de non-droit !
\end{attention}

\begin{popfigure}
\centering
\begin{tikzpicture}[font=\small]
% Axe chronologique
\draw[->, very thick, popline] (0,0) -- (13.2,0);
\foreach \x/\annee in {1.5/2001, 6.5/2014, 11.5/2018}{
  \draw[thick, popdark] (\x,-0.12) -- (\x,0.12);
  \node[below, font=\small\bfseries, text=popdark] at (\x,-0.15) {\annee};
}
% Budapest 2001 (au-dessus)
\node[draw=popblue, fill=popblueL, rounded corners=4pt, align=center, inner sep=5pt, text width=3.6cm, anchor=south] at (1.5,0.35) {\textbf{Convention de Budapest}\\ Premier traité mondial contre la cybercriminalité (Conseil de l'Europe)};
% Malabo 2014 (en dessous)
\node[draw=popgreen, fill=popgreenL, rounded corners=4pt, align=center, inner sep=5pt, text width=3.6cm, anchor=north] at (6.5,-1.1) {\textbf{Convention de Malabo}\\ Cybersécurité et protection des données en Afrique (Union africaine)};
% RGPD 2018 (au-dessus)
\node[draw=poporange, fill=poporangeL, rounded corners=4pt, align=center, inner sep=5pt, text width=3.6cm, anchor=south] at (11.5,0.35) {\textbf{RGPD}\\ Protection des données personnelles, amendes jusqu'à 4 \% du chiffre d'affaires mondial (Union européenne)};
\end{tikzpicture}
\legende{Frise chronologique des grands textes internationaux relatifs à l'éthique numérique : Budapest (2001), Malabo (2014), RGPD (2018).}
\end{popfigure}

\courssub{Exercice résolu : appliquer le cadre international à une situation concrète}

\begin{exoresolu}[Une fraude transfrontalière]
\textbf{Énoncé.} Depuis son quartier à Douala, un individu crée un faux site de vente de téléphones et escroque une vingtaine de clients situés en France, pour un préjudice total de 15 000 euros. Le site est hébergé sur un serveur aux Pays-Bas.
\begin{enumerate}
\item Pourquoi la seule loi camerounaise ne suffit-elle pas à traiter cette affaire ?
\item Cite deux mécanismes internationaux qui peuvent aider à l'enquête.
\item L'individu peut-il être inquiété alors qu'il n'a jamais quitté le Cameroun ? Justifie.
\end{enumerate}
\tcblower
\textbf{Solution détaillée.}
\begin{enumerate}
\item La loi camerounaise ne s'applique que sur le territoire du Cameroun. Or ici, les \textbf{victimes} sont en France et le \textbf{serveur} (preuve essentielle) est aux Pays-Bas. La police camerounaise ne peut ni saisir le serveur aux Pays-Bas ni entendre les victimes en France sans coopération internationale.
\item Deux mécanismes possibles :
\begin{itemize}
\item la \textbf{coopération prévue par la Convention de Budapest} : demande de préservation puis de saisie des données du serveur aux Pays-Bas, échange d'informations entre polices ;
\item l'action d'\textbf{Interpol} (et d'Europol côté européen) pour relier les polices camerounaise, française et néerlandaise et coordonner l'enquête. On peut aussi citer l'\textbf{entraide judiciaire} (transmission des plaintes et des preuves) et l'\textbf{extradition}.
\end{itemize}
\item Oui, il peut être inquiété. D'une part, il peut être \textbf{poursuivi au Cameroun} selon la loi nationale sur la cybercriminalité, car l'infraction a été commise depuis le territoire camerounais. D'autre part, la France peut demander son \textbf{extradition} pour le juger, car ses victimes sont françaises. Le fait de n'avoir jamais quitté le Cameroun ne le met donc pas à l'abri : la cybercriminalité transfrontalière est punissable des deux côtés.
\end{enumerate}
\end{exoresolu}

\begin{exercice}[À toi de jouer]
Une entreprise camerounaise développe une application mobile utilisée par des élèves en France. Elle collecte leurs noms, photos et localisations sans leur demander leur accord.
\begin{enumerate}
\item Quel texte international cette entreprise risque-t-elle de violer ?
\item Quel type de sanction encourt-elle ?
\item Cite un droit que les élèves français peuvent exercer sur leurs données.
\end{enumerate}
\end{exercice}

\begin{corrige}[Exercice : application mobile et données personnelles]
\begin{enumerate}
\item L'entreprise risque de violer le \textbf{RGPD} (Règlement Général sur la Protection des Données, Union européenne, 2018), car elle traite des données personnelles de personnes se trouvant en Europe sans consentement.
\item Elle encourt une \textbf{amende administrative} pouvant atteindre 20 millions d'euros ou 4 \% de son chiffre d'affaires mondial annuel (le montant le plus élevé étant retenu).
\item Les élèves peuvent exercer leur \textbf{droit d'accès} (savoir quelles données sont détenues), leur \textbf{droit de rectification} (corriger des erreurs) ou leur \textbf{droit à l'effacement} (faire supprimer leurs données).
\end{enumerate}
\end{corrige}

\begin{aretenir}[Bilan de la section]
\begin{itemize}
\item La cybercriminalité ignore les frontières : des \textbf{traités internationaux} sont indispensables pour harmoniser les lois et organiser la coopération.
\item La \textbf{Convention de Budapest} (2001) est le premier traité mondial contre la cybercriminalité ; la \textbf{Convention de Malabo} (2014) est son équivalent africain ; le \textbf{RGPD} (2018) protège les données personnelles en Europe avec des amendes très lourdes.
\item \textbf{Interpol} et \textbf{Europol} coordonnent les polices ; l'\textbf{ONU} et l'\textbf{UIT} émettent des recommandations (non contraignantes).
\item L'\textbf{extradition} et l'\textbf{entraide judiciaire} permettent de poursuivre un criminel même s'il agit depuis l'étranger.
\end{itemize}
\end{aretenir}

\coursec{Comparaison des sanctions nationales et internationales}

Dans la section précédente, nous avons découvert le cadre juridique international de l'éthique numérique : la Convention de Budapest sur la cybercriminalité, le RGPD européen, les résolutions de l'ONU et le rôle d'Interpol. Nous savons aussi, depuis le début de cette leçon, que le Cameroun dispose de ses propres lois (loi n° 2010/012 du 21 décembre 2010 relative à la cybersécurité et à la cybercriminalité, loi de 2010 sur le commerce électronique, etc.) avec ses propres sanctions. Une question se pose alors naturellement : \emph{que se passe-t-il quand une même infraction numérique peut être jugée à la fois par la justice camerounaise et par des mécanismes internationaux ?} Qu'est-ce qui ressemble, qu'est-ce qui diffère, et comment ces deux niveaux travaillent-ils ensemble ? C'est l'objet de cette section.

\courssub{Pourquoi comparer les deux niveaux ?}

Imaginons une situation concrète. Un jeune pirate installé à Douala s'introduit dans le serveur d'une banque située en France et vole les données bancaires de clients européens. Qui peut le punir ? Le Cameroun, parce que l'infraction a été commise depuis son territoire ? La France, parce que les victimes sont françaises ? Les deux ? Pour répondre, il faut comprendre que le droit national et le droit international ne sont pas concurrents : ils sont \cle{complémentaires}. Chacun a son territoire d'action, ses outils et ses limites.

\begin{propriete}[Cumul possible des compétences (admise)]
Une même infraction numérique peut relever \emph{simultanément} du droit national (par exemple la loi camerounaise de 2010) et d'instruments internationaux (par exemple la Convention de Budapest ou le RGPD). On admet ce principe en s'appuyant sur la \cle{coopération judiciaire} : les États signataires s'engagent à s'entraider pour identifier, poursuivre et juger les cybercriminels, sans qu'un seul niveau de droit ne bloque l'autre.
\end{propriete}

\courssub{Les points communs entre les deux niveaux}

Avant d'opposer les deux niveaux, remarquons d'abord qu'ils poursuivent les \emph{mêmes objectifs}. Au niveau national comme au niveau international, on réprime :

\begin{itemize}
\item le \cle{piratage} : l'accès non autorisé à un système informatique (intrusion dans un serveur, vol de mot de passe) ;
\item la \cle{fraude} informatique : escroqueries en ligne, faux sites de paiement, hameçonnage (\emph{phishing}), arnaques aux transferts d'argent ;
\item les \cle{atteintes aux données} : vol, destruction, modification ou divulgation de données personnelles sans autorisation ;
\item les \cle{atteintes aux personnes} : cyberharcèlement, diffamation en ligne, usurpation d'identité, publication d'images intimes sans consentement.
\end{itemize}

Autrement dit, un acte interdit à Yaoundé est presque toujours interdit aussi à Paris, à Nairobi ou à Washington. Les deux niveaux partagent la même philosophie : \textbf{protéger les systèmes, les données et les personnes}.

\courssub{Les différences de nature des sanctions}

C'est sur la \emph{nature} des sanctions que les deux niveaux se distinguent le plus nettement.

\textbf{Au niveau national}, les sanctions sont celles du droit pénal classique, prononcées par les tribunaux camerounais contre des \emph{personnes physiques} :

\begin{itemize}
\item des \cle{peines de prison} (de quelques mois à plusieurs années, pouvant atteindre 20 ans, voire la réclusion à perpétuité pour les formes les plus graves de cybercriminalité selon la loi de 2010) ;
\item des \cle{amendes} fixées en francs CFA (de quelques centaines de mille à plusieurs millions de FCFA selon la gravité) ;
\item éventuellement des peines complémentaires : confiscation du matériel, interdiction d'exercer une activité.
\end{itemize}

\textbf{Au niveau international}, les outils sont différents :

\begin{itemize}
\item des \cle{amendes proportionnelles au chiffre d'affaires} : c'est l'arme principale du RGPD contre les entreprises. Une société qui protège mal les données de ses clients européens peut payer jusqu'à 4 \% de son chiffre d'affaires mondial annuel, ou 20 millions d'euros (le montant le plus élevé étant retenu) ;
\item la \cle{coopération judiciaire} : échange de preuves, enquêtes conjointes, extradition, entraide entre polices via Interpol. L'objectif n'est pas de punir directement, mais de permettre aux États de punir efficacement.
\end{itemize}

\begin{exemplebox}[Une amende proportionnelle, ça change tout]
Une petite boutique de Douala qui fraude risque une amende fixe prévue par la loi camerounaise. En revanche, un géant du numérique qui gagne des milliards d'euros par an ne serait pas impressionné par une amende fixe : c'est pourquoi le RGPD prévoit des amendes en \emph{pourcentage du chiffre d'affaires mondial}. Pour une entreprise réalisant 10 milliards d'euros de chiffre d'affaires, une amende de 4 \% représente 400 millions d'euros : de quoi faire réfléchir sérieusement sa direction !
\end{exemplebox}

\courssub{Les différences d'échelle (champ d'application)}

La deuxième grande différence concerne \emph{l'espace} sur lequel chaque droit s'applique.

\begin{itemize}
\item Le \cle{droit national} camerounais s'applique sur le \textbf{territoire camerounais} : il punit les infractions commises au Cameroun, depuis le Cameroun, ou contre des systèmes et des victimes camerounais.
\item Les \cle{instruments internationaux} (Convention de Budapest, RGPD, coopération Interpol) visent la \textbf{coopération entre États} : ils n'ont pas de territoire propre, ils relient les territoires entre eux. Leur rôle est de faire en sorte qu'un criminel ne soit pas en sécurité simplement parce que sa victime se trouve dans un autre pays.
\end{itemize}

\begin{attention}[Le RGPD ne s'applique pas automatiquement à tout le monde]
Ne dis jamais que « le RGPD protège tous les citoyens du monde ». Le RGPD est un règlement \emph{européen} : il protège les données des personnes se trouvant dans l'\textbf{Union européenne}. Un citoyen camerounais vivant à Bafoussam n'est pas directement protégé par le RGPD ; c'est la loi camerounaise (notamment celle de 2010) qui le protège. En revanche, dès qu'une entreprise — même camerounaise — traite les données de résidents européens, elle doit respecter le RGPD.
\end{attention}

\courssub{Tableau comparatif synthétique}

Le schéma suivant résume la comparaison selon quatre critères : c'est exactement la démarche à reproduire dans une copie d'examen.

\begin{popfigure}
\begin{center}
\begin{tikzpicture}[
  cell/.style={draw=popline, fill=white, text width=3.6cm, minimum height=1.1cm, align=center, font=\small, inner sep=4pt},
  crit/.style={draw=popline, fill=popgoldL, text width=2.6cm, minimum height=1.1cm, align=center, font=\small\bfseries, inner sep=4pt},
  head/.style={draw=popline, fill=popblue, text=white, text width=3.6cm, minimum height=0.9cm, align=center, font=\small\bfseries, inner sep=4pt},
  headc/.style={draw=popline, fill=popblue, text=white, text width=2.6cm, minimum height=0.9cm, align=center, font=\small\bfseries, inner sep=4pt}
]
% En-têtes
\node[headc] at (0,0) {Critère};
\node[head]  at (3.25,0) {Niveau national};
\node[head]  at (6.95,0) {Niveau international};
% Ligne 1
\node[crit] at (0,-1.05) {Champ d'application};
\node[cell] at (3.25,-1.05) {Territoire camerounais};
\node[cell] at (6.95,-1.05) {Coopération entre États};
% Ligne 2
\node[crit] at (0,-2.15) {Infractions visées};
\node[cell] at (3.25,-2.15) {Piratage, fraude, atteintes aux données et aux personnes};
\node[cell] at (6.95,-2.15) {Les mêmes, dès qu'elles traversent les frontières};
% Ligne 3
\node[crit] at (0,-3.45) {Sanctions};
\node[cell] at (3.25,-3.45) {Prison, amendes en FCFA, confiscations};
\node[cell] at (6.95,-3.45) {Amendes en \% du chiffre d'affaires, entraide judiciaire};
% Ligne 4
\node[crit] at (0,-4.75) {Institutions compétentes};
\node[cell] at (3.25,-4.75) {Tribunaux camerounais, ANTIC, police et gendarmerie};
\node[cell] at (6.95,-4.75) {Interpol, CNIL (France), États parties à la Convention de Budapest};
\end{tikzpicture}
\end{center}
\legende{Tableau comparatif des dispositifs juridiques national (Cameroun) et international.}
\end{popfigure}

\begin{methode}[Comment comparer deux dispositifs juridiques]
Pour comparer proprement deux systèmes de sanctions (par exemple droit camerounais et instruments internationaux), suis toujours les quatre étapes suivantes :
\begin{enumerate}
\item Comparer le \textbf{champ d'application} : sur quel territoire, pour quelles personnes ?
\item Comparer les \textbf{infractions visées} : quels actes sont punis ?
\item Comparer les \textbf{sanctions} : nature (prison, amende fixe, amende proportionnelle) et montants.
\item Comparer les \textbf{institutions compétentes} : qui enquête, qui juge, qui sanctionne ?
\end{enumerate}
Cette méthode en quatre points te garantit une comparaison complète et bien structurée.
\end{methode}

\courssub{Complémentarité : comment les deux niveaux travaillent ensemble}

La vraie force du système apparaît quand les deux niveaux \emph{se combinent}. La Convention de Budapest permet par exemple à un État d'obtenir des \cle{preuves situées à l'étranger} : si le serveur contenant les traces d'une attaque se trouve dans un autre pays signataire, la justice camerounaise peut demander officiellement la conservation et la transmission de ces données.

\begin{exemplebox}[Un hacker camerounais attaque une banque européenne]
Un pirate installé à Yaoundé s'introduit dans le système d'une banque française et détourne de l'argent. Deux scénarios sont possibles :
\begin{itemize}
\item il est \textbf{poursuivi au Cameroun}, car l'infraction a été commise depuis le territoire national, et la loi de 2010 s'applique ;
\item il est \textbf{extradé} vers la France, si un accord d'extradition le permet, pour y être jugé au plus près des victimes.
\end{itemize}
Dans les deux cas, la coopération judiciaire (partage des preuves, entraide policière via Interpol) rend la poursuite possible. Sans instruments internationaux, le pirate pourrait se croire intouchable du simple fait que sa victime est à l'étranger.
\end{exemplebox}

Le diagramme de Venn ci-dessous montre visuellement que les deux niveaux sanctionnent en grande partie les \emph{mêmes} infractions, chacun avec ses spécificités.

\begin{popfigure}
\begin{center}
\begin{tikzpicture}[font=\small]
\draw[fill=popblue!25, draw=popblue, thick] (-1.6,0) circle (2.4);
\draw[fill=poporange!25, draw=poporange, thick] (1.6,0) circle (2.4);
\node[font=\small\bfseries, text=popblue] at (-2.6,2.1) {Niveau national};
\node[font=\small\bfseries, text=poporange] at (2.6,2.1) {Niveau international};
\node[align=center, text width=2.1cm] at (-2.9,0) {Prison, amendes en FCFA, jugement par les tribunaux camerounais};
\node[align=center, text width=2.1cm] at (2.9,0) {Amendes en \% du chiffre d'affaires, extradition, entraide};
\node[align=center, text width=2.3cm, font=\small\bfseries] at (0,0.4) {Piratage, fraude, atteintes aux données et aux personnes};
\node[align=center, text width=2.3cm] at (0,-0.9) {(infractions communes)};
\end{tikzpicture}
\end{center}
\legende{Diagramme de Venn : infractions et sanctions aux niveaux national et international.}
\end{popfigure}

\courssub{Le cas particulier des entreprises multinationales}

Une entreprise multinationale (réseau social, opérateur de téléphonie, plateforme de commerce en ligne) est soumise \emph{en même temps} :

\begin{itemize}
\item aux \textbf{lois nationales} de chaque pays où elle exerce (au Cameroun : loi de 2010, contrôle de l'ANTIC) ;
\item aux \textbf{règlements étrangers} dès qu'elle traite les données de résidents étrangers. Ainsi, une application camerounaise qui collecte les données d'utilisateurs vivant en France doit respecter le RGPD, même si son siège est à Douala.
\end{itemize}

C'est une conséquence directe de la différence d'échelle : le territoire n'est plus seulement géographique, il devient \emph{le territoire des données}. Dès que des données de résidents européens sont traitées, le droit européen « voyage » jusqu'à l'entreprise, où qu'elle soit.

\courssub{Les limites de la coopération internationale}

Il serait faux de croire que ce système est parfait. Trois limites importantes doivent être connues :

\begin{enumerate}
\item \textbf{Les pays non signataires} : tous les États n'ont pas signé la Convention de Budapest. Si le criminel ou le serveur se trouve dans un pays non coopératif, l'entraide judiciaire devient très difficile, voire impossible.
\item \textbf{La lenteur des procédures d'entraide} : obtenir une preuve à l'étranger peut prendre des mois, alors qu'un criminel peut effacer ses traces en quelques secondes. C'est pourquoi la Convention prévoit des procédures de \emph{conservation urgente} des données.
\item \textbf{L'anonymat des criminels} : grâce aux réseaux privés virtuels (VPN), aux fausses identités et aux serveurs relais, un cybercriminel peut cacher sa véritable localisation, ce qui complique énormément l'enquête.
\end{enumerate}

\begin{savaistu}[Interpol et l'Afrique contre la cybercriminalité]
Certaines grandes cyberattaques internationales ont été déjouées grâce à des \textbf{opérations conjointes d'Interpol} impliquant des pays africains. Par exemple, des opérations coordonnées sur tout le continent ont permis d'identifier des réseaux d'escrocs spécialisés dans la fraude bancaire et l'hameçonnage, conduisant à des arrestations simultanées dans plusieurs pays. Le Cameroun participe à cette coopération policière internationale : la cybercriminalité ne connaît pas les frontières, la réponse des polices non plus !
\end{savaistu}

\begin{exoresolu}[Comparer et argumenter]
\textbf{Énoncé.} Une entreprise camerounaise de vente en ligne collecte les données personnelles de clients vivant au Cameroun et en Belgique. Elle néglige totalement la sécurité de sa base de données, qui est piratée. (1) Devant quels dispositifs juridiques cette entreprise peut-elle répondre de ses actes ? (2) Quels types de sanctions encourt-elle à chaque niveau ?
\tcblower
\textbf{Solution détaillée.} Appliquons la méthode de comparaison en quatre points.
\begin{enumerate}
\item \emph{Champ d'application} : les clients camerounais relèvent du droit national ; les clients belges, résidents de l'Union européenne, relèvent du RGPD.
\item \emph{Infraction visée} : dans les deux cas, il s'agit d'une atteinte aux données personnelles (protection insuffisante puis vol de données).
\item \emph{Sanctions} : au niveau national, l'entreprise encourt des amendes en FCFA et ses responsables des peines pouvant aller jusqu'à la prison, selon la loi camerounaise de 2010 ; au niveau européen, elle encourt une amende pouvant atteindre 4 \% de son chiffre d'affaires mondial.
\item \emph{Institutions} : les tribunaux camerounais et l'ANTIC d'un côté ; les autorités européennes de protection des données de l'autre, avec une coopération judiciaire possible.
\end{enumerate}
\emph{Conclusion} : une même négligence expose l'entreprise à une double sanction, nationale et internationale, car elle traite des données de résidents de deux espaces juridiques différents.
\end{exoresolu}

\begin{aretenir}[L'essentiel de la section]
\begin{itemize}
\item Les deux niveaux répriment les \textbf{mêmes infractions} : piratage, fraude, atteintes aux données et aux personnes.
\item Le droit \textbf{national} punit par des peines de \textbf{prison} et des \textbf{amendes fixes} ; le niveau \textbf{international} utilise des \textbf{amendes proportionnelles au chiffre d'affaires} et la \textbf{coopération judiciaire}.
\item Le droit national s'applique sur le \textbf{territoire} ; les instruments internationaux organisent la \textbf{coopération entre États}.
\item Une entreprise traitant des données de résidents étrangers est soumise \textbf{à la fois} aux lois nationales et aux règlements étrangers (RGPD).
\item Limites : pays non signataires, lenteur de l'entraide, anonymat des criminels.
\end{itemize}
\end{aretenir}

\coursec{Application à des situations concrètes et démarche argumentée}

Dans la section précédente, nous avons comparé les sanctions prévues au niveau national (loi camerounaise de 2010 sur la cybersécurité et la cybercriminalité) et au niveau international (notamment la Convention de Budapest). Mais connaître les textes ne suffit pas : au BEPC comme dans la vie courante, il faut savoir les \emph{appliquer} à des situations concrètes. Face à un fait réel — un compte piraté, une fausse boutique en ligne, une fuite de données —, le citoyen numérique doit être capable de répondre à trois questions : de quelle infraction s'agit-il ? Quel texte la punit ? Quelle sanction est encourue ? C'est toute la démarche que nous allons apprendre ici, pas à pas, à travers quatre études de cas tirées de la vie quotidienne au Cameroun.

\courssub{La démarche d'analyse juridique d'un cas concret}

Pourquoi une méthode ? Parce que devant une situation, on a tendance à réagir avec les émotions : « c'est mal », « il mérite la prison ». Or le droit ne fonctionne pas ainsi. Le juge, et toi dans ta copie du BEPC, doit raisonner de manière ordonnée : partir du \cle{fait}, lui donner un nom juridique (la \cle{qualification}), citer le \cle{texte de loi} applicable, en déduire la \cle{sanction} encourue, et terminer par un conseil de \cle{prévention}. Cette démarche en cinq étapes est un outil universel : elle fonctionne pour tous les cas, nationaux ou internationaux.

\begin{popfigure}
\begin{center}
\begin{tikzpicture}[
    node distance=0.55cm,
    etape/.style={rectangle, rounded corners=3pt, draw=popblue, thick, fill=popblueL,
        text width=4.6cm, minimum height=1.05cm, align=center, font=\small},
    fleche/.style={-{Stealth[length=2.5mm]}, thick, popdark}
]
\node[etape, align=center] (e1){\textbf{1. Fait}\\ Décrire objectivement ce qui s'est passé};
\node[etape, below=of e1, align=center] (e2){\textbf{2. Qualification juridique}\\ Nommer l'infraction (accès frauduleux, escroquerie, diffamation\ldots)};
\node[etape, below=of e2, align=center] (e3){\textbf{3. Texte applicable}\\ Citer la loi (loi n° 2010/012, Convention de Budapest\ldots)};
\node[etape, below=of e3, align=center] (e4){\textbf{4. Sanction encourue}\\ Peines de prison et/ou amendes prévues};
\node[etape, below=of e4, fill=popgreenL, draw=popgreen, align=center] (e5){\textbf{5. Prévention}\\ Conseil éthique et citoyen pour éviter le cas};
\draw[fleche] (e1) -- (e2);
\draw[fleche] (e2) -- (e3);
\draw[fleche] (e3) -- (e4);
\draw[fleche] (e4) -- (e5);
\end{tikzpicture}
\legende{Les cinq étapes de la démarche d'analyse juridique d'un cas concret.}
\end{center}
\end{popfigure}

\begin{methode}[Rédiger et justifier une conclusion argumentée]
Pour rédiger une réponse ou une conclusion complète et bien justifiée, suis ces étapes :
\begin{enumerate}
    \item \textbf{Annoncer la qualification} : « Les faits décrits constituent l'infraction de\ldots »
    \item \textbf{Citer le texte} : «\ldots prévue et punie par la loi n° 2010/012 du 21 décembre 2010 relative à la cybersécurité et à la cybercriminalité au Cameroun\ldots »
    \item \textbf{Conclure sur la sanction} : « L'auteur encourt une peine d'emprisonnement de\ldots et une amende de\ldots »
    \item \textbf{Proposer une attitude éthique} : « Pour se prémunir, il convient de\ldots »
\end{enumerate}
Une bonne copie ne dit jamais seulement « c'est interdit » : elle dit \emph{pourquoi} c'est interdit (le texte) et \emph{ce que risque} l'auteur (la sanction).
\end{methode}

\begin{attention}[Faute morale ou infraction pénale ?]
Erreur fréquente : confondre la \cle{faute morale} et l'\cle{infraction pénale}. Dire du mal d'un camarade dans la cour est moralement condamnable, mais n'est pas forcément un délit. En droit, \textbf{seule une infraction prévue par un texte est sanctionnée par la justice} : c'est le principe de légalité (« pas de peine sans loi »). Dans ta copie, vérifie toujours qu'un texte (loi nationale ou convention internationale) punit bien le fait décrit avant d'annoncer une sanction.
\end{attention}

\courssub{Étude de cas 1 : piratage du compte WhatsApp d'un camarade}

\textbf{Les faits.} Junior, élève de 3ème, récupère le code de vérification reçu par SMS sur le téléphone de sa camarade Amina laissé sans surveillance. Il prend le contrôle de son compte WhatsApp et diffuse ses photos personnelles dans plusieurs groupes de la classe.

Appliquons la démarche. Le fait est clair : Junior s'est introduit dans un compte qui ne lui appartient pas et a publié des données privées. La \emph{qualification} est double : \cle{accès frauduleux} à un système de traitement de données (le compte) et \cle{atteinte à la vie privée} par diffusion de données personnelles sans consentement. Le \emph{texte applicable} est la loi n° 2010/012 du 21 décembre 2010 relative à la cybersécurité et à la cybercriminalité, qui punit l'accès frauduleux à un système informatique ainsi que les atteintes aux données personnelles. La \emph{sanction encourue} peut atteindre des peines d'emprisonnement et des amendes de plusieurs millions de francs CFA, aggravées lorsque des données intimes sont divulguées. Enfin, la \emph{prévention} : ne jamais communiquer son code de vérification, activer la double authentification et verrouiller son téléphone.

\begin{exoresolu}[Rédaction guidée d'un paragraphe argumenté]
Énoncé : rédige un paragraphe de 5 à 8 lignes analysant juridiquement le cas de Junior, en suivant la démarche en cinq étapes.
\tcblower
\textbf{Solution.} « En prenant le contrôle du compte WhatsApp d'Amina grâce à son code de vérification, puis en diffusant ses photos personnelles, Junior a commis deux infractions : l'accès frauduleux à un système de traitement de données et l'atteinte à la vie privée d'autrui. Ces infractions sont prévues et punies par la loi n° 2010/012 du 21 décembre 2010 relative à la cybersécurité et à la cybercriminalité au Cameroun. Junior encourt donc des peines d'emprisonnement ainsi que des amendes pouvant atteindre plusieurs millions de francs CFA, d'autant que les données divulguées sont intimes. Pour éviter de telles situations, chaque élève doit protéger ses comptes par un mot de passe fort et la double authentification, et ne jamais partager ses codes de vérification. Respecter la vie privée des autres sur Internet est un devoir de tout citoyen numérique. »
\end{exoresolu}

\courssub{Étude de cas 2 : le faux site de vente en ligne d'un commerçant de Yaoundé}

\textbf{Les faits.} Un commerçant de Yaoundé crée un site imitant une grande boutique en ligne, avec de fausses promotions (téléphones à moitié prix). Des clients paient par mobile money et ne reçoivent jamais rien. Le préjudice total s'élève à 3 millions de francs CFA.

\begin{itemize}
    \item \textbf{Qualification} : il s'agit d'une \cle{escroquerie informatique} (obtenir de l'argent par un procédé frauduleux utilisant un système informatique), associée à l'usurpation d'identité de la vraie boutique et à la création d'un faux site.
    \item \textbf{Texte applicable} : la loi n° 2010/012 de 2010 punit expressément la fraude et l'escroquerie commises au moyen des technologies de l'information ; le Code pénal camerounais punit également l'escroquerie en général.
    \item \textbf{Sanction encourue} : emprisonnement et amendes lourdes, proportionnés au préjudice ; la dimension internationale (si des victimes ou des serveurs sont à l'étranger) peut engager la coopération prévue par la Convention de Budapest.
    \item \textbf{Prévention} : avant tout achat en ligne, vérifier l'adresse exacte du site, l'existence de mentions légales, les avis de clients, et se méfier des prix anormalement bas.
\end{itemize}

\courssub{Étude de cas 3 : diffusion de données clients sans consentement}

\textbf{Les faits.} Une entreprise camerounaise de vente à crédit transmet à des démarcheurs, sans accord de ses clients, leurs noms, numéros de téléphone et adresses. Certains clients sont ensuite harcelés d'appels publicitaires.

Ici, la responsabilité est double. Au niveau \textbf{national}, la loi de 2010 et les textes sur la protection des données personnelles imposent le \cle{consentement} de la personne avant toute collecte ou transmission de ses données ; l'ANTIC (Agence Nationale des Technologies de l'Information et de la Communication) veille à la sécurité des systèmes d'information et peut intervenir. L'entreprise encourt des sanctions pénales et financières. Au niveau \textbf{international}, si les données concernent des personnes situées à l'étranger ou transitent par des serveurs étrangers, les normes internationales de protection des données (principes de la Convention de Budapest et des règlements comme le RGPD européen pour les échanges avec l'Europe) peuvent s'appliquer et engager la responsabilité de l'entreprise au-delà des frontières. La prévention : toute entreprise doit informer ses clients, recueillir leur consentement et sécuriser ses bases de données.

\courssub{Étude de cas 4 : fausse information sur Facebook ayant provoqué des troubles}

\textbf{Les faits.} Une rumeur annonçant la fermeture d'un marché de Douala est publiée sur Facebook. Partagée des milliers de fois, elle provoque une panique et des affrontements sur place.

Qui est responsable ? D'abord \textbf{l'auteur} de la publication : la loi de 2010 punit la diffusion de fausses informations (fake news) par voie électronique lorsqu'elle trouble l'ordre public ; il encourt des peines d'emprisonnement et des amendes. Mais attention : \textbf{ceux qui relaient} massivement une information qu'ils savent fausse peuvent aussi voir leur responsabilité engagée, car le partage volontaire contribue à la diffusion. La prévention est ici essentiellement citoyenne : avant de partager, vérifier la source, la date, et croiser l'information avec des médias fiables. Un bon réflexe : « si je ne suis pas sûr, je ne partage pas ».

\begin{savaistu}[L'ANTIC et les jeunes]
La sensibilisation à l'éthique numérique fait partie des missions officielles de l'ANTIC. L'agence organise régulièrement des campagnes et des journées de sensibilisation à destination des jeunes et des établissements scolaires, afin d'apprendre aux élèves à se protéger en ligne et à adopter un comportement responsable sur les réseaux sociaux. Ton collège peut donc un jour accueillir une équipe de l'ANTIC !
\end{savaistu}

\courssub{Conseils de prévention et de bonne conduite numérique}

La meilleure sanction, c'est celle qu'on évite. Voici les attitudes d'un citoyen numérique responsable :

\begin{itemize}
    \item \textbf{Protéger ses comptes} : mots de passe forts et différents, double authentification, ne jamais partager ses codes.
    \item \textbf{Respecter autrui} : ne jamais publier la photo ou les données de quelqu'un sans son accord.
    \item \textbf{Vérifier avant de partager} : une information non vérifiée ne se relaie pas.
    \item \textbf{Se méfier des offres trop belles} : un prix anormalement bas cache souvent une escroquerie.
    \item \textbf{Connaître ses recours} : en cas d'infraction, conserver les preuves (captures d'écran) et saisir les services compétents ou l'ANTIC.
\end{itemize}

\courssub{Synthèse de la leçon}

\begin{popfigure}
\begin{center}
\small
\begin{tabularx}{\linewidth}{|X|l|X|X|}
\hline
\rowcolor{popblueL} \textbf{Texte} & \textbf{Institution} & \textbf{Infractions couvertes} & \textbf{Sanctions} \\
\hline
Loi n° 2010/012 du 21/12/2010 (cybersécurité et cybercriminalité) & ANTIC, tribunaux camerounais & Accès frauduleux, atteinte à la vie privée, escroquerie en ligne, fake news, atteinte aux données & Emprisonnement et amendes (jusqu'à plusieurs millions de FCFA) \\
\hline
Code pénal camerounais & Tribunaux & Escroquerie, diffamation, injures & Prison et amendes selon la gravité \\
\hline
Convention de Budapest (2001) & Coopération entre États & Cybercriminalité transfrontalière (fraude, accès illégal, contenus illicites) & Harmonisation des peines entre pays signataires \\
\hline
\end{tabularx}
\legende{Tableau récapitulatif : textes, institutions, infractions et sanctions.}
\end{center}
\end{popfigure}

\begin{aretenir}[L'essentiel pour le BEPC]
Dans toute réponse argumentée sur l'éthique numérique, sache citer \textbf{au moins deux textes} — la loi n° 2010/012 du 21 décembre 2010 et la Convention de Budapest — et \textbf{au moins une institution} — l'ANTIC. Structure toujours ta réponse en cinq étapes : fait, qualification, texte, sanction, prévention. C'est cette méthode qui fait la différence entre une copie moyenne et une excellente copie.
\end{aretenir}

\begin{exercice}[À toi de jouer]
Un élève crée un faux profil Facebook au nom de son professeur et publie des messages insultants en son nom. En suivant la démarche en cinq étapes, rédige un paragraphe argumenté (qualification, texte applicable, sanction encourue, conseil de prévention).
\end{exercice}

\begin{corrige}[Exercice]
\textbf{Fait} : création d'un faux profil et publication de messages insultants au nom d'autrui. \textbf{Qualification} : usurpation d'identité numérique et diffamation (atteinte à l'honneur). \textbf{Texte} : loi n° 2010/012 du 21 décembre 2010 (atteintes aux personnes via les TIC) et Code pénal pour la diffamation. \textbf{Sanction} : peines d'emprisonnement et amendes prévues par ces textes. \textbf{Prévention} : ne jamais utiliser l'identité d'autrui en ligne ; la victime doit conserver les captures d'écran et signaler le compte à la plateforme et aux autorités compétentes.
\end{corrige}

\newpage

\coursec{Activité d'intégration}

\coursec{Leçon 10 : Connaissance de la loi nationale et internationale relative à l'éthique numérique}

\courssub{Objectifs de l'activité}

À l'issue de cette activité d'intégration, tu seras capable de :

\begin{itemize}
\item \textbf{Qualifier une infraction numérique} à partir des faits décrits dans une situation concrète de la vie courante ;
\item \textbf{Citer le texte de loi applicable} au niveau national (loi n° 2010/012 du 21 décembre 2010 relative à la cybersécurité et à la cybercriminalité au Cameroun, code pénal camerounais) ;
\item \textbf{Proposer une sanction réaliste et proportionnée}, en tenant compte de la gravité des faits et de la situation du contrevenant (ici, un mineur) ;
\item \textbf{Rédiger et justifier une conclusion juridique argumentée} sous la forme d'un jugement de 5 à 8 lignes ;
\item \textbf{Défendre oralement une position} lors d'un débat contradictoire, dans le respect de la parole d'autrui.
\end{itemize}

Cette activité mobilise les deux savoir-faire de la leçon : appliquer les notions de l'unité à une situation concrète, et rédiger une démarche justifiée.

\courssub{Situation}

\textbf{Le cas de la fausse annonce au Collège Bilingue de Nkol-Eton.}

Le lundi 12 janvier 2026, à 5 h 30 du matin, une publication apparaît sur Facebook. Elle émane d'un compte intitulé « Collège Bilingue de Nkol-Eton -- Page Officielle », orné du logo de l'établissement et de la photo du principal, copiés sur la véritable page de l'école. Le message annonce : « En raison d'une épidémie de méningite, les cours sont suspendus du 12 au 16 janvier. Les élèves sont priés de rester à domicile. -- Le Principal. »

En deux heures, la publication est partagée 340 fois et commentée plus de 500 fois. À 7 h 30, seuls 87 élèves sur les 1 240 que compte l'établissement se présentent au portail. Les enseignants attendent dans des salles vides. Des parents inquiets affluent vers le collège ; deux d'entre eux, affolés, conduisent leurs enfants à l'hôpital de district pour des examens inutiles. Le principal, alerté, dément l'information à 9 h et saisit le commissariat de police du quartier.

L'enquête de la brigade de gendarmerie, appuyée par l'Agence Nationale des Technologies de l'Information et de la Communication (ANTIC), remonte jusqu'à l'adresse IP utilisée : la fausse page a été créée le dimanche 11 janvier à 21 h 14 depuis un téléphone portable. Le propriétaire de la ligne est identifié, puis l'utilisateur réel : \textbf{M. Junior E., élève de 3\textsuperscript{ème} 4 du même collège, âgé de 14 ans}. Interrogé en présence de ses parents, il reconnaît les faits : il voulait « échapper à l'interrogation de mathématiques du lundi matin » et pensait que ce n'était « qu'une blague ». La fausse page a été supprimée le lundi à 14 h, après 21 heures d'existence.

La classe de 3\textsuperscript{ème} 4 est chargée de simuler le procès de cette affaire. La classe est divisée en quatre groupes :

\begin{itemize}
\item \textbf{Groupe 1 -- le Ministère public (procureur)} : il accuse et requiert une sanction ;
\item \textbf{Groupe 2 -- la Défense (avocat de Junior)} : il plaide les circonstances atténuantes ;
\item \textbf{Groupe 3 -- le Juge} : il écoute les deux parties et rédige le jugement ;
\item \textbf{Groupe 4 -- le Jury populaire} : il délibère et rend un avis sur la culpabilité et la peine.
\end{itemize}

\courssub{Consignes}

Chaque groupe prépare sa contribution en répondant, par écrit, aux questions suivantes.

\begin{enumerate}
\item \textbf{Qualification des faits.} Quelles infractions numériques Junior a-t-il commises ? Pour chacune, indique l'élément matériel (ce qu'il a fait concrètement) et l'élément moral (son intention).
\item \textbf{Textes applicables.} Cite au moins deux textes nationaux applicables : un article de la loi n° 2010/012 du 21 décembre 2010 relative à la cybersécurité et à la cybercriminalité, et une disposition du code pénal camerounais (par exemple sur l'usurpation d'identité ou la diffusion de fausses nouvelles).
\item \textbf{Sanctions encourues.} D'après la loi de 2010, quelles peines (emprisonnement, amendes) encourent un auteur d'usurpation d'identité numérique et de publication de fausses informations ?
\item \textbf{Proportionnalité.} Junior a 14 ans : c'est un mineur. Comment la justice des mineurs adapte-t-elle les sanctions ? Propose une sanction réaliste et éducative, en la justifiant.
\item \textbf{Rédaction du jugement.} Le groupe « Juge » rédige un jugement argumenté de 5 à 8 lignes comportant : le rappel des faits, la qualification juridique, le texte appliqué, la sanction prononcée et sa justification.
\item \textbf{Débat.} Chaque groupe présente sa position en 3 minutes. Le jury délibère 5 minutes et rend son avis. La classe compare ensuite les décisions des différents groupes.
\item \textbf{Ouverture internationale.} Si la fausse annonce avait été hébergée sur un serveur situé à l'étranger ou avait visé un établissement d'un autre pays, quels instruments internationaux auraient pu s'appliquer ? Cite au moins un texte international relatif à la cybercriminalité.
\end{enumerate}

\courssub{Ressources à mobiliser}

Pour mener à bien cette activité, rappelle-toi les notions essentielles de la leçon.

\begin{aretenir}[Les textes nationaux de référence]
\begin{itemize}
\item \textbf{Loi n° 2010/012 du 21 décembre 2010} relative à la cybersécurité et à la cybercriminalité au Cameroun : elle définit et punit notamment l'\cle{usurpation d'identité numérique} (art. 7), l'\cle{accès frauduleux} à un système informatique, la \cle{publication de fausses informations} par voie électronique (art. 9), la diffamation et l'atteinte à la vie privée en ligne. Les peines vont de l'amende (de 200 000 à plusieurs millions de FCFA) à l'emprisonnement (de quelques mois à plusieurs années) selon la gravité.
\item \textbf{Code pénal camerounais} : il punit la diffusion de fausses nouvelles (article 240), l'usurpation de fonctions et d'identité, la diffamation et l'injure.
\item \textbf{Justice des mineurs} : un mineur de moins de 18 ans relève d'une juridiction spéciale ; les mesures éducatives et la réparation priment sur la prison (Ordonnance n° 62/OF/18 du 12 mars 1962 modifiée).
\end{itemize}
\end{aretenir}

\begin{aretenir}[Les repères internationaux]
\begin{itemize}
\item \textbf{Convention de Budapest} (Conseil de l'Europe, 2001) : premier traité international sur la cybercriminalité, que le Cameroun a rejointe par adhésion en 2022 ; elle harmonise la répression de l'usurpation d'identité et favorise la coopération entre États (entraide judiciaire, extradition).
\item \textbf{Convention de Malabo} (Union africaine, 2014) : protection des données personnelles et lutte contre la cybercriminalité en Afrique.
\item \textbf{Règlement général sur la protection des données (RGPD)} : référence européenne pour la protection des données personnelles, applicable aux traitements concernant des résidents de l'UE.
\end{itemize}
\end{aretenir}

\begin{methode}[Démarche d'analyse juridique en 4 étapes]
\begin{enumerate}
\item \textbf{Constater les faits} : qui a fait quoi, quand, avec quel outil, quelles conséquences ?
\item \textbf{Qualifier juridiquement} : à quelle infraction définie par la loi correspondent ces faits ?
\item \textbf{Identifier le texte applicable} : article de loi national ou international, peines prévues.
\item \textbf{Proposer et justifier la sanction} : proportionnée aux faits, adaptée à l'auteur (âge, intention, récidive).
\end{enumerate}
\end{methode}

\courssub{Démarche et corrigé commenté}

\begin{exoresolu}[Corrigé commenté du tribunal numérique simulé]
\textbf{Énoncé.} À partir du cas de la fausse annonce au Collège Bilingue de Nkol-Eton, qualifier les infractions commises, citer les textes applicables, proposer une sanction réaliste et rédiger un jugement argumenté.

\tcblower

\textbf{Étape 1 : constatation des faits.}

Junior E., 14 ans, a : (a) créé un compte Facebook reprenant le nom, le logo et la photo du principal de son collège ; (b) publié une fausse annonce de fermeture pour épidémie ; (c) provoqué l'absence de 1 153 élèves sur 1 240, un mouvement de panique des parents et deux consultations médicales inutiles. L'intention est établie : il voulait échapper à une interrogation.

\textbf{Étape 2 : qualification juridique.}

\begin{itemize}
\item \textbf{Usurpation d'identité numérique} : créer un compte au nom d'une institution et d'une personne (le principal) sans autorisation est puni par l'article 7 de la loi n° 2010/012 du 21 décembre 2010. L'élément matériel est la création du faux compte ; l'élément moral est la volonté de se faire passer pour l'école.
\item \textbf{Publication de fausses informations par voie électronique} : la fausse annonce de fermeture, diffusée volontairement, relève de l'article 9 de la loi de 2010 et de l'article 240 du code pénal sur la diffusion de fausses nouvelles susceptibles de troubler l'ordre public.
\item Le \textbf{trouble causé à l'ordre public} (panique, absentéisme massif, dépenses médicales inutiles) constitue une circonstance aggravante.
\end{itemize}

\textbf{Étape 3 : sanctions encourues.}

Pour un majeur, l'usurpation d'identité numérique (art. 7) et la diffusion de fausses informations (art. 9) sont passibles, selon la loi de 2010, d'un emprisonnement de 6 mois à 5 ans et d'amendes de 500 000 à 10 000 000 FCFA (cumuls possibles). Mais Junior est \textbf{mineur} : la justice des mineurs privilégie les mesures éducatives.

\textbf{Étape 4 : proposition de sanction réaliste (défense et procureur).}

\begin{itemize}
\item \emph{Réquisitions du procureur} : reconnaître la culpabilité pour usurpation d'identité et diffusion de fausses nouvelles ; rappel à la loi ; obligation de présenter des excuses publiques au principal et aux élèves ; travail éducatif (exposé sur la cybercriminalité devant toute l'école) ; confiscation temporaire du téléphone par les parents ; suivi par le conseiller d'orientation.
\item \emph{Plaidoirie de la défense} : âge de l'auteur (14 ans), absence de casier, reconnaissance des faits, mobile non crapuleux (peur d'une interrogation), suppression rapide du compte : demander l'indulgence et des mesures purement éducatives.
\end{itemize}

\textbf{Étape 5 : jugement rédigé (groupe « Juge »).}

« Le Tribunal, après avoir entendu le Ministère public, la Défense et l'avis du Jury, constate que Junior E., 14 ans, a créé le 11 janvier 2026 un compte Facebook usurpant l'identité du Collège Bilingue de Nkol-Eton et de son principal, et y a publié une fausse annonce de fermeture ayant provoqué l'absentéisme de 1 153 élèves et un trouble à l'ordre public. Ces faits sont qualifiés d'usurpation d'identité numérique (art. 7) et de publication de fausses informations par voie électronique (art. 9), infractions prévues et punies par la loi n° 2010/012 du 21 décembre 2010 relative à la cybersécurité et à la cybercriminalité, ainsi que par l'article 240 du code pénal. Attendu que le prévenu est mineur, qu'il a reconnu les faits et exprimé des regrets, le Tribunal le déclare coupable et le condamne à un rappel à la loi, à des excuses publiques à l'établissement, et à la réalisation d'une campagne de sensibilisation à l'éthique numérique au sein du collège, sous le contrôle de ses parents et du conseiller d'orientation. »

\textbf{Étape 6 : ouverture internationale.}

Si le serveur ou la victime se trouvaient à l'étranger, la \textbf{Convention de Budapest} permettrait la coopération entre polices des deux pays (échange de preuves, entraide judiciaire, art. 23 à 35), et la \textbf{Convention de Malabo} encadrerait la protection des données personnelles usurpées. La cybercriminalité ne connaît pas les frontières : c'est pourquoi les États harmonisent leurs sanctions.
\end{exoresolu}

\courssub{Bilan de l'activité}

Ce tribunal numérique simulé t'a permis de mettre en évidence plusieurs acquis essentiels.

\begin{itemize}
\item \textbf{Une infraction numérique est une infraction comme une autre} : créer un faux compte ou diffuser une fausse nouvelle sur Internet n'est pas « une blague » ; c'est un délit défini et puni par la loi camerounaise de 2010 et par le code pénal.
\item \textbf{La sanction est proportionnée} : elle dépend de la gravité des faits, de l'intention et de l'âge de l'auteur ; pour un mineur, les mesures éducatives et la réparation priment sur la prison.
\item \textbf{La démarche juridique est rigoureuse} : on ne juge pas « au sentiment » ; on constate les faits, on qualifie, on cite le texte, puis on motive la sanction. C'est exactement la démarche de rédaction argumentée exigée au BEPC.
\item \textbf{Le droit national s'inscrit dans un cadre international} : la Convention de Budapest et la Convention de Malabo montrent que la lutte contre la cybercriminalité dépasse les frontières et exige la coopération des États.
\item \textbf{Sur le plan du savoir-être} : tu as appris à écouter une position contraire, à argumenter sans agresser, et à délibérer collectivement — des attitudes de citoyen numérique responsable.
\end{itemize}

\begin{attention}[Erreurs fréquentes à éviter]
\begin{itemize}
\item Confondre la \textbf{qualification des faits} (ce qui est interdit par la loi) avec l'\textbf{appréciation morale} (« c'est méchant ») : un jugement cite toujours un texte de loi.
\item Oublier que l'\textbf{intention} compte autant que l'acte : sans volonté de nuire ou de tromper, la qualification peut changer.
\item Proposer une sanction irréaliste (prison ferme pour un mineur de 14 ans primo-délinquant) : la peine doit rester proportionnée et conforme à la justice des mineurs.
\item Croire que « sur Internet, on est anonyme » : l'enquête de l'ANTIC a remonté l'adresse IP en quelques heures.
\end{itemize}
\end{attention}

\begin{savaistu}[L'ANTIC, gardienne du cyberespace camerounais]
Créée en 2002, l'Agence Nationale des Technologies de l'Information et de la Communication (ANTIC) est l'organe chargé de la sécurité des systèmes d'information au Cameroun. Elle assiste la police et la gendarmerie dans les enquêtes numériques, délivre les certifications de sécurité et sensibilise les jeunes à l'usage responsable d'Internet. Le Cameroun a par ailleurs adhéré en 2022 à la Convention de Budapest, devenant l'un des premiers pays africains à rejoindre ce traité international contre la cybercriminalité.
\end{savaistu}

\newpage

\coursec{Méthodes et exercices résolus}

\coursec{Leçon 10 : Connaissance de la loi nationale et internationale relative à l'éthique numérique}

\courssub{Cadre juridique national du numérique au Cameroun}

\begin{definition}[Les trois piliers législatifs camerounais]
    Le cadre juridique national repose sur trois textes fondamentaux que tout citoyen numérique doit connaître :
    \begin{itemize}
        \item \cle{Loi n°2010/012 du 21 décembre 2010} relative à la cybersécurité et à la cybercriminalité. C'est le texte de référence principal qui définit les infractions spécifiques au numérique (piratage, escroquerie en ligne, atteintes aux données, contenus illicites) et leurs sanctions.
        \item \cle{Loi n°2010/013 du 21 décembre 2010} régissant les communications électroniques. Elle organise le secteur des télécoms, pose le principe de protection des données à caractère personnel (Art. 118) et confie la régulation à l'ARTEL.
        \item \cle{Code pénal (Loi n°2016/007)} et Code de procédure pénale. Ils répriment les infractions classiques commises par voie numérique (vol, escroquerie, diffamation, atteinte à la vie privée) et fixent les règles de poursuite et de jugement.
    \end{itemize}
\end{definition}

\begin{propriete}[Principes directeurs du droit numérique camerounais]
    \begin{itemize}
        \item \textbf{Légalité des peines :} Nul ne peut être puni pour un acte qui n'était pas prévu par la loi au moment où il a été commis (Art. 1er Code pénal, Art. 1er Loi 2010/012).
        \item \textbf{Proportionnalité :} La sanction doit être adaptée à la gravité de l'infraction.
        \item \textbf{Protection de la vie privée :} Le secret des correspondances et la protection des données personnelles sont des droits fondamentaux (Art. 1er Loi 2010/012, Art. 118 Loi 2010/013).
        \item \textbf{Libre expression encadrée :} La liberté d'opinion s'exerce dans le respect de la loi, de l'ordre public et des droits d'autrui (pas d'injure, diffamation, incitation à la haine, fausses nouvelles).
    \end{itemize}
\end{propriete}

\begin{methode}[Identifier le texte applicable face à une situation]
    \textbf{Objectif :} Savoir quel texte citer pour qualifier une infraction numérique.
    \begin{enumerate}
        \item \textbf{Infraction spécifique au système ou aux données (piratage, virus, rançongiciel) :} \rightarrow \cle{Loi 2010/012} (Arts 7 à 18).
        \item \textbf{Infraction liée au contenu (harcèlement, fausses nouvelles, pédopornographie, apologie terrorisme) :} \rightarrow \cle{Loi 2010/012} (Arts 19 à 26).
        \item \textbf{Infraction classique « numérisée » (escroquerie, diffamation, vol) :} \rightarrow \cle{Code pénal} + \cle{Loi 2010/012 Art. 21} (procédure de preuve numérique).
        \item \textbf{Problème de régulation télécoms / protection données personnelles (fuite données, spam) :} \rightarrow \cle{Loi 2010/013} + \cle{ARTEL / APDP}.
    \end{enumerate}
\end{methode}

\courssub{Les institutions nationales de régulation et de sanction}

\begin{definition}[ANTIC -- Agence Nationale des Technologies de l'Information et de la Communication]
    \textbf{Tutelle :} MINPOSTEL. \textbf{Base légale :} Art. 5 Loi 2010/012, Décret 2012/164.
    \textbf{Missions clés :}
    \begin{itemize}
        \item Sécurité des systèmes d'information de l'État et infrastructures critiques.
        \item Certification et homologation des solutions de sécurité.
        \item \textbf{Répression administrative :} Constat d'infractions, mise en demeure, suspension de service, saisine du Procureur de la République.
        \item Gestion du domaine \texttt{.cm} et CERT-CM (Centre de Veille et de Réponse aux Attaques Informatiques).
    \end{itemize}
\end{definition}

\begin{definition}[ARTEL -- Agence de Régulation des Télécommunications]
    \textbf{Tutelle :} MINPOSTEL. \textbf{Base légale :} Loi 2010/013.
    \textbf{Missions clés :} Régulation économique et technique (licences, interconnexion, qualité de service, tarifs), protection des consommateurs, gestion du spectre de fréquences.
    \textbf{Sanctions administratives :} Amendes, injonctions, retrait/suspension de licences.
\end{definition}

\begin{definition}[Autorité de Protection des Données Personnelles (APDP)]
    \textbf{Base légale :} Art. 118 Loi 2010/013. Instituée par la loi, sa mise en place opérationnelle complète est en cours (décret d'organisation récent). Elle veillera au respect de la loi sur les données personnelles.
    \textbf{Sanctions prévues :} Avertissement, injonction, amende administrative, saisine du Procureur.
    \textit{Note :} En attendant sa pleine opérationnalité, l'ANTIC assure une partie de ces missions.
\end{definition}

\begin{propriete}[Distinction fondamentale : Voie administrative vs Voie pénale]
    \begin{center}
    \begin{tabularx}{\linewidth}{|l|X|X|}
    \hline
    \rowcolor{popblueL}
    \textbf{Critère} & \textbf{Voie Administrative (ANTIC, ARTEL, APDP)} & \textbf{Voie Pénale (Procureur, Tribunaux)} \\
    \hline
    \textbf{Autorité} & Autorité administrative indépendante & Autorité judiciaire (Procureur, Juge) \\
    \hline
    \textbf{But} & Faire cesser le manquement, réguler le secteur, protéger l'intérêt général & Punir l'auteur, réparer le préjudice de la victime, maintenir l'ordre public \\
    \hline
    \textbf{Sanctions types} & Amende administrative, mise en demeure, suspension service, retrait licence, publication sanction & Emprisonnement, amende pénale (Trésor Public), confiscation, interdiction de droits, dommages-intérêts \\
    \hline
    \textbf{Procédure} & Contradictoire, rapide, sans juge & Respect des droits de la défense, instruction, audience publique, appel possible \\
    \hline
    \textbf{Lien} & L'administration \textbf{saisit le Procureur} si les faits révèlent un crime/délit (Art. 40 CPP) & Le juge peut ordonner des mesures administratives complémentaires \\
    \hline
    \end{tabularx}
    \end{center}
    \textbf{Règle d'or :} Les deux voies sont \textbf{cumulatives}. Une entreprise peut payer une amende administrative \textbf{ET} être condamnée pénalement pour les mêmes faits.
\end{propriete}

\courssub{Infractions numériques et sanctions au niveau national}

\begin{propriete}[Principales infractions de la Loi 2010/012 et sanctions encourues (Adultes)]
    \begin{center}
    \begin{tabularx}{\linewidth}{|l|l|X|}
    \hline
    \rowcolor{popblueL}
    \textbf{Infraction (Libellé simplifié)} & \textbf{Article} & \textbf{Peines maximales encourues (Prison / Amende FCFA)} \\
    \hline
    Accès frauduleux à un système (STDAD) & Art. 7 & 5 ans / 10 000 000 \\
    Maintien frauduleux dans un système & Art. 8 & 5 ans / 10 000 000 \\
    \textbf{Modification / Suppression de données (par accès frauduleux)} & \textbf{Art. 9} & \textbf{7 ans / 20 000 000} \\
    Interception illicite de communications & Art. 10 & 5 ans / 10 000 000 \\
    Entrave au fonctionnement d'un système (virus, DoS) & Art. 11 & 5 ans / 10 000 000 \\
    Introduction frauduleuse de données (faux numérique) & Art. 13 & 5 ans / 10 000 000 \\
    \hline
    Atteinte à l'intégrité de données (virus, corruption sans accès) & Art. 15 & 5 ans / 10 000 000 \\
    \hline
    Pédopornographie (production, diffusion, détention image mineur) & Art. 20 & 10 ans / 50 000 000 \\
    Escroquerie par voie numérique (arnaque 419, faux site) & Art. 21 (renvoie Art. 316 CPP) & 5 ans / 5 000 000 \\
    Usurpation d'identité numérique & Art. 23 & 3 ans / 5 000 000 \\
    Harcèlement par voie électronique & Art. 24 & 2 ans / 2 000 000 \\
    Diffusion de fausses nouvelles (trouble à l'ordre public) & Art. 25 & 2 ans / 5 000 000 \\
    \hline
    \end{tabularx}
    \end{center}
    \textbf{Peines complémentaires fréquentes (Art. 27 Loi 2010/012, Art. 318 CPP) :} Confiscation du matériel (PC, téléphone, serveur), interdiction d'exercer une activité, publication du jugement, interdiction de séjour.
\end{propriete}

\begin{attention}[Régime spécial des mineurs (Élèves de 3ème : 14-15 ans)]
    \textbf{Âge de la responsabilité pénale :} 10 ans (Ordonnance 1962).
    \textbf{De 10 à 18 ans :} Le \textbf{Tribunal pour Enfants} est compétent.
    \begin{itemize}
        \item \textbf{Principes :} Primauté de la mesure éducative sur la sanction (protection, surveillance, éducation).
        \item \textbf{Mesures :} Réprimande, remise à parents, liberté surveillée, placement en centre éducatif (ouvert ou fermé).
        \item \textbf{Peine de prison :} Exceptionnelle, seulement si les mesures éducatives ont échoué ou pour crimes graves. \textbf{Maximum divisé par 2} par rapport à l'adulte (ex: max 5 ans au lieu de 10 pour pédopornographie).
        \item \textbf{Amendes :} Prononcées mais souvent converties en travaux d'intérêt général ou payées par les responsables légaux.
    \end{itemize}
    \textbf{Au BEPC :} Toute analyse d'un cas impliquant un mineur \textbf{doit} mentionner ce régime spécifique.
\end{attention}

\begin{exemplebox}[Cas concret : Piratage du serveur de notes (Kevin, 15 ans)]
    \textbf{Faits :} Kevin accède au serveur du collège avec le mot de passe d'un prof (shoulder surfing) et modifie ses notes.
    \textbf{Qualification :}
    \begin{enumerate}
        \item Accès frauduleux (Art. 7 Loi 2010/012).
        \item Modification frauduleuse de données (Art. 9 Loi 2010/012) -- \textbf{Infraction principale la plus grave}.
    \end{enumerate}
    \textbf{Sanctions encourues (Adulte) :} 1 à 7 ans prison + 1 000 000 à 20 000 000 FCFA (Art. 9).
    \textbf{Application au mineur (15 ans) :} Tribunal pour Enfants. Mesures éducatives prioritaires (liberté surveillée, placement). Peine de prison possible mais max 3,5 ans. Amende possible (parents). Réparation du préjudice (remise en état notes, dommages-intérêts pour le collège).
\end{exemplebox}

\courssub{Cadre juridique international de l'éthique numérique}

\begin{definition}[Convention de Budapest (2001) -- Conseil de l'Europe]
    \textbf{Statut au Cameroun :} Signée 2018, Ratifiée 2022, Entrée en vigueur \textbf{1er mars 2023}.
    \textbf{Portée :} \cle{Seul traité international contraignant} sur la cybercriminalité.
    \textbf{3 piliers :}
    \begin{enumerate}
        \item \textbf{Harmonisation des infractions :} Accès illicite, interception, atteintes aux données, infractions liées au contenu (pédopornographie), infractions au droit d'auteur.
        \item \textbf{Pouvoirs d'enquête adaptés :} Conservation données, perquisition informatique, interception en temps réel.
        \item \textbf{Coopération internationale efficace :} Entraide judiciaire (Autorités centrales : MINJUSTICE au Cameroun), \textbf{Réseau 24/7} (Points de contact permanents : CERT-CM/ANTIC au Cameroun) pour geler les preuves d'urgence à l'étranger.
    \end{enumerate}
    \textbf{Principe clé :} \cle{Double incrimination} -- L'acte doit être infraction dans les deux pays pour que l'entraide fonctionne.
\end{definition}

\begin{definition}[Convention de Malabo (2014) -- Union Africaine]
    \textbf{Statut au Cameroun :} Ratifiée (dépôt 2019).
    \textbf{Spécificité :} Adaptée au contexte africain. Couvre : 1) Protection données personnelles (droits étendus), 2) Sécurité des SI, 3) Cybercriminalité (alignée sur Budapest), 4) Coopération au niveau UA (Mécanisme de suivi).
    \textbf{Complémentarité :} Budapest = standard mondial répression/entraide ; Malabo = cadre panafricain données + cyber + coopération UA.
\end{definition}

\begin{aretenir}[Autres références internationales]
    \begin{itemize}
        \item \textbf{RGPD (UE 2016/679) :} Non contraignant directement au Cameroun, mais \textbf{standard de référence « Gold »} pour la protection des données (DPO, DPIA, droit à l'oubli, portabilité). S'applique aux entreprises camerounaises ciblant des résidents UE (extraterritorialité).
        \item \textbf{Résolutions ONU / Normes UIT (X.1205, X.1303) :} Cadres techniques et politiques pour la sécurité et la lutte contre la cybercriminalité.
    \end{itemize}
\end{aretenir}

\courssub{Comparaison des sanctions nationales et internationales}

\begin{methode}[Réussir un exercice de comparaison (Type BEPC)]
    \textbf{Consigne type :} « Comparez la sanction prévue par la loi camerounaise et l'exigence de la Convention de Budapest pour l'atteinte à l'intégrité des données. »
    \textbf{Méthode :} Construire un tableau à 4 colonnes : \textbf{Critère} | \textbf{Droit Camerounais} | \textbf{Standard International (Budapest/Malabo)} | \textbf{Analyse (Convergence / Divergence)}.
    \textbf{Critères obligatoires à traiter :}
    \begin{enumerate}
        \item \textbf{Champ d'application :} Territorialité stricte (National) vs Territorialité + Extraterritorialité (Nationalité, Intérêt protégé -- Budapest Art. 22).
        \item \textbf{Definition de l'infraction :} Comparer les libellés (ex: Art. 9 \& 15 Loi 2010/012 vs Art. 4 Budapest).
        \item \textbf{Niveau des peines (Prison) :} Max national vs Exigence Budapest (Art. 13 : « peines privatives de liberté pouvant aller jusqu'à un maximum d'au moins 2 ans » pour les infractions Art. 2-6).
        \item \textbf{Niveau des amendes :} Montants fixes CFA vs Critère « effectives, proportionnées, dissuasives » (Budapest Art. 13).
        \item \textbf{Tentative :} Punissable dans les deux (Art. 33 Loi 2010/012 = Art. 7 Budapest).
        \item \textbf{Responsabilité personnes morales :} Prévue dans les deux (Art. 34 Loi 2010/012 = Art. 12 Budapest).
    \end{enumerate}
\end{methode}

\begin{exoresolu}[Comparaison : Atteinte à l'intégrité des données (Art. 9/15 vs Budapest Art. 4/13)]
    \textbf{Énoncé :} Comparez les sanctions pour « atteinte à l'intégrité des données » (effacement, altération, virus).
    \tcblower
    \textbf{Solution détaillée :}
    \begin{center}
    \begin{tabularx}{\linewidth}{|l|X|X|X|}
    \hline
    \rowcolor{popblueL}
    \textbf{Critère} & \textbf{Droit Camerounais (Loi 2010/012)} & \textbf{Convention de Budapest (Art. 4 \& 13)} & \textbf{Analyse} \\
    \hline
    \textbf{Actes réprimés} & Art. 9 : Modif/suppression \emph{par accès frauduleux}. Art. 15 : Endommagement \emph{sans accès frauduleux} (virus, corruption). & Art. 4 §1 : Endommagement etc. \emph{par accès}. Art. 4 §2 : Endommagement etc. \emph{sans accès} (virus). & \textbf{Parallélisme parfait.} Le droit camerounais distingue les deux cas comme Budapest. \\
    \hline
    \textbf{Peine Prison (Max)} & Art. 9 : \textbf{7 ans}. Art. 15 : \textbf{5 ans}. & Art. 13 : Max \textbf{au moins 2 ans} (pour infractions Art. 2-6). & \textbf{Conforme + Protecteur.} Max national (7 ans) $>$ Standard min Budapest (2 ans). \\
    \hline
    \textbf{Peine Prison (Min)} & Art. 9 : 1 an. Art. 15 : 6 mois. & Pas de minimum imposé par la Convention. & \textbf{Conforme.} Laisser au juge la marge d'appréciation. \\
    \hline
    \textbf{Amendes (Max)} & Art. 9 : \textbf{20 000 000 FCFA}. Art. 15 : \textbf{10 000 000 FCFA}. & « Effectives, proportionnées, dissuasives ». Pas de montant fixe. & \textbf{Conforme.} Montants significatifs (jusqu'à 20M $\approx$ 30k €), critères respectés. \\
    \hline
    \textbf{Tentative} & Art. 33 : Punie comme infraction consommée. & Art. 7 : Punissable. & \textbf{Identique.} \\
    \hline
    \textbf{Personnes Morales} & Art. 34 : Amende x5 (max 100M FCFA), peines compl. (exclusion marchés). & Art. 12 : Responsabilité, sanctions effectives. & \textbf{Convergence forte.} Cameroun plus précis sur le quantum (facteur 5). \\
    \hline
    \end{tabularx}
    \end{center}
    \textbf{Conclusion :} Le droit camerounais est \textbf{pleinement conforme} aux exigences minimales de Budapest. Il va même au-delà sur les maximums de peine (prison et amende) et la précision de la responsabilité des personnes morales. La ratification de 2023 a scellé cette conformité.
\end{exoresolu}

\courssub{Application à des situations concrètes et démarche argumentée}

\begin{methode}[La méthode « F.A.R.C. » pour le cas pratique (Niveau 3ème)]
    \textbf{Objectif :} Structurer sa réponse pour « Appliquer les notions à une situation concrète » et « Justifier sa démarche ».
    \begin{enumerate}
        \item \textbf{F -- Faits juridiques :} Écrire \textbf{uniquement} les faits qui constituent l'infraction (Qui ? Quoi ? Comment ? Quand ? Moyen numérique ?). Pas de sentiments.
        \item \textbf{A -- Articles de loi (Règles) :} Citer \textbf{précisément} le numéro de l'article et la loi (ex: « Article 24 de la Loi n°2010/012 du 21 décembre 2010 »).
        \item \textbf{R -- Raisonnement (Application) :} \textbf{Étape clé.} Faire le lien : « L'élément matériel est constitué car [fait précis]... L'élément moral (intention) est constitué car [fait précis]... »
        \item \textbf{C -- Conclusion / Sanctions :} Dire : « L'auteur encourt [peine prison] et [amende] ». \textbf{Ajouter :} « Comme l'auteur est mineur (15 ans), le Tribunal pour Enfants appliquera des mesures éducatives (liberté surveillée, placement) et la peine de prison max sera divisée par 2. »
    \end{enumerate}
    \textbf{Conseil :} Une phrase = une idée. Vocabulaire précis : \cle{constitue}, \cle{caractérise}, \cle{encourt}, \cle{relevé de}.
\end{methode}

\begin{exoresolu}[Cas pratique : Harcèlement et usurpation d'identité sur réseau social (Niveau 3ème)]
    \textbf{Énoncé :} Alice (15 ans) découvre un faux profil Facebook à son nom avec ses photos. Le compte publie des insultes et des montages humiliants (non pornographiques). L'auteur demande 50 000 FCFA pour supprimer le compte. L'enquête identifie Bob (16 ans), camarade de classe. Analysez la situation de Bob.
    \tcblower
    \textbf{Solution détaillée (Méthode F.A.R.C.) :}
    \begin{enumerate}
        \item \textbf{Faits :} Création compte faux (usurpation identité), Publication insultes (harcèlement), Tentative obtention argent par menace (chantage/extorsion), Auteur : Bob, 16 ans, mineur.
        \item \textbf{Articles de loi (Loi 2010/012) :}
              \begin{itemize}
                  \item \textbf{Art. 23 :} Usurpation d'identité numérique (1 à 3 ans / 500k à 5M FCFA).
                  \item \textbf{Art. 24 :} Harcèlement par voie électronique (6 mois à 2 ans / 500k à 2M FCFA).
                  \item \textbf{Art. 21 (renvoie Art. 316/317 CPP) :} Extorsion / Escroquerie par voie numérique (1 à 5 ans / 100k à 5M FCFA). \textit{Note :} La tentative (demande argent sans paiement) est punie comme l'infraction (Art. 33).
              \end{itemize}
        \item \textbf{Raisonnement :}
              \begin{itemize}
                  \item \textbf{Usurpation :} Bob a utilisé nom/photos d'Alice sans accord $\rightarrow$ Élément matériel Art. 23 constitué. Volonté de se faire passer pour elle $\rightarrow$ Élément moral constitué.
                  \item \textbf{Harcèlement :} Messages insultants répétés via Facebook $\rightarrow$ Troubles anormaux vie d'Alice $\rightarrow$ Art. 24 constitué.
                  \item \textbf{Extorsion (tentative) :} Menace (garder compte actif) pour obtenir argent $\rightarrow$ Art. 21/317 CPP constitué (tentative punie comme consommée Art. 33).
              \end{itemize}
        \item \textbf{Conclusion -- Sanctions pour Bob (Mineur 16 ans) :}
              \begin{itemize}
                  \item \textbf{Concours d'infractions :} Usurpation + Harcèlement + Tentative extorsion. Principe : \cle{Confusion des peines} (on ne cumule pas les prisons, on prend la plus forte : Extorsion max 5 ans adulte).
                  \item \textbf{Régime mineur (Tribunal pour Enfants) :} Mesures éducatives prioritaires. \textbf{Peine prison max divisée par 2 = 2,5 ans} (exceptionnelle). Mesures probables : Liberté surveillée avec obligations (réparation, stage citoyenneté) ou Placement en centre éducatif fermé si récidive/grande dangerosité.
                  \item \textbf{Amendes :} Encourues (ex: 5M max) mais souvent converties en Travail d'Intérêt Général (TIG) ou payées par parents.
                  \item \textbf{Partie civile :} Alice (via parents) peut demander dommages-intérêts (préjudice moral, atteinte honneur/image).
              \end{itemize}
    \end{enumerate}
\end{exoresolu}

\begin{exercice}[Entraînement BEPC -- Cas pratique]
    \textbf{Sujet :} Dans un cybercafé à Douala, Paul (17 ans) installe un logiciel espion (keylogger) sur les postes publics. Il récupère les identifiants Mobile Money de 5 clients et vole au total 200 000 FCFA. Il est identifié grâce aux logs du cybercafé.
    \textbf{Questions :}
    \begin{enumerate}
        \item Qualifiez juridiquement les faits de Paul en citant les articles exacts de la Loi 2010/012 et/ou du Code Pénal.
        \item Indiquez les sanctions maximales encourues par un adulte pour ces infractions.
        \item Paul a 17 ans. Quelle juridiction est compétente ? Quelles sont les mesures ou peines spécifiques applicables à son âge ?
        \item L'ANTIC peut-elle condamner Paul à de la prison ? Justifiez.
    \end{enumerate}
\end{exercice}


\begin{corrige}[Exercice n°1 -- Entraînement BEPC]
    \begin{enumerate}
        \item \textbf{Qualifications :}
              \begin{itemize}
                  \item \textbf{Introduction frauduleuse de données / Logiciel espion :} \textbf{Art. 13 Loi 2010/012} (« introduit frauduleusement des données... »). Ou \textbf{Art. 11} (Entrave fonctionnement système par programme malveillant).
                  \item \textbf{Interception illicite de correspondances / Données :} \textbf{Art. 10 Loi 2010/012} (« intercepte... des communications électroniques » ou « des données »).
                  \item \textbf{Vol / Escroquerie par voie numérique :} \textbf{Art. 21 Loi 2010/012} renvoyant à \textbf{Art. 316 CPP (Escroquerie)} ou \textbf{Art. 317 CPP (Extorsion/Vol avec violence)} selon qualification exacte (vol des codes $\rightarrow$ vol ; usage codes pour virer argent $\rightarrow$ escroquerie). \textbf{Art. 316 CPP} (Escroquerie) est le plus adapté : manœuvres frauduleuses (keylogger) $\rightarrow$ remise volontaire argent (virement) $\rightarrow$ préjudice.
              \end{itemize}
        \item \textbf{Sanctions maximales (Adulte) :}
              \begin{itemize}
                  \item Art. 10/11/13 : 5 ans prison + 10 000 000 FCFA.
                  \item Art. 21 / Art. 316 CPP : 5 ans prison + 5 000 000 FCFA.
                  \item \textbf{Peine la plus forte retenue (Confusion) :} \textbf{5 ans de prison} et \textbf{10 000 000 FCFA} d'amende (cumule amendes possibles).
              \end{itemize}
        \item \textbf{Régime Mineur (17 ans) :}
              \begin{itemize}
                  \item \textbf{Juridiction :} \textbf{Tribunal pour Enfants} (TPI si pas de TPE spécialisé, mais chambre d'enfants).
                  \item \textbf{Mesures :} Primauté éducative. Liberté surveillée, placement centre éducatif (fermé possible vu gravité/réitération 5 victimes).
                  \item \textbf{Peine prison :} Possible mais exceptionnelle. Max = 5 ans / 2 = \textbf{2,5 ans}.
                  \item \textbf{Amendes :} Maintenues mais TIG ou paiement parents.
              \end{itemize}
        \item \textbf{Rôle de l'ANTIC :} \textbf{NON.} L'ANTIC a un pouvoir de \textbf{sanction administrative} (amende administrative, mise en demeure, suspension service, saisine Procureur). \textbf{Seul le Juge (Tribunal) peut condamner à de la prison et à une amende pénale.} L'ANTIC transmet le dossier au Procureur de la République.
    \end{enumerate}
\end{corrige}

\begin{methode}[Rédiger une réponse argumentée (Question de réflexion type BEPC)]
    \textbf{Sujet type :} « Malgré les lois, la cybercriminalité augmente. Proposez 3 mesures pour améliorer la lutte. »
    \textbf{Structure en 3 parties (Problème - Causes - Solutions) :}
    \begin{enumerate}
        \item \textbf{Introduction :} Constat chiffré (ex: Rapport ANTIC +30\% incidents) $\rightarrow$ Problématique : « L'arsenal juridique est nécessaire mais insuffisant sans moyens d'application et prévention. » $\rightarrow$ Annonce plan.
        \item \textbf{Développement (3 arguments = 3 paragraphes) :}
              \begin{itemize}
                  \item \textbf{1. Formation des acteurs (Police, Gendarmerie, Juges) :} Manque de compétences techniques (preuves numériques, hash, horodatage). $\rightarrow$ \textbf{Mesure :} Cycles de formation continue obligatoires (ENAM, ENSP, ANTIC), création de pôles « cyber » dans les tribunaux (Yaoundé, Douala).
                  \item \textbf{2. Coopération internationale rapide :} 80\%+ auteurs/serveurs à l'étranger. CRI trop lente. $\rightarrow$ \textbf{Mesure :} Renforcer CERT-CM (ANTIC) pour réseau 24/7 (Budapest Art. 35), accords bilatéraux simplifiés (MLAT) avec pays sources (Nigeria, France, USA, Chine).
                  \item \textbf{3. Prévention \& Éducation (Le bouclier amont) :} Victimes par ignorance (phishing, OTP, mots de passe). $\rightarrow$ \textbf{Mesure :} Module « Hygiène \& Éthique Numérique » obligatoire au collège (6ème-3ème). Obligation légale pour plateformes (Facebook, WhatsApp, TikTok) d'avoir un représentant légal au Cameroun et délais de retrait contenus illicites (24h-72h) sous peine d'amende.
              \end{itemize}
        \item \textbf{Conclusion :} Loi = squelette. Formation, Coopération, Éducation = muscles. L'IA générative (deepfakes, arnaques automatisées) impose une veille législative permanente.
    \end{enumerate}
\end{methode}

\begin{exoresolu}[Sujet de réflexion : Efficacité de l'arsenal juridique]
    \textbf{Énoncé :} « Malgré l'arsenal juridique existant, la cybercriminalité progresse au Cameroun. » Discutez et proposez trois axes d'amélioration.
    \tcblower
    \textbf{Introduction :} Le Cameroun dispose d'un cadre juridique robuste (Loi 2010/012, ratification Budapest 2023, Malabo). Pourtant, le rapport ANTIC 2023 signale une hausse sensible des incidents (arnaques MM, piratage, harcèlement). L'arsenal est \textbf{nécessaire mais non suffisant} sans une chaîne pénale opérationnelle et une prévention massive.
    \textbf{Développement :}
    \begin{enumerate}
        \item \textbf{Axe 1 : Professionnalisation de la chaîne pénale (Le maillon faible).}
              \begin{itemize}
                  \item \textbf{Constat :} OPJ et Magistrats manquent souvent de formation technique (preuve numérique, chaîne de custody, horodatage, hash). Preuves souvent irrecevables (Art. 32-37 Loi 2010/012 mal appliqués). Affaires classées sans suite.
                  \item \textbf{Mesure concrète :} Création de \textbf{Sections Cybercriminelles spécialisées} au sein des TGI (Yaoundé, Douala, Bafoussam) avec magistrats/greffiers formés (ENAM + ANTIC). Formation continue obligatoire « Preuve Numérique » pour tous les OPJ (Police/Gendarmerie).
              \end{itemize}
        \item \textbf{Axe 2 : Opérationnalisation de la coopération internationale (Le trou noir transfrontalier).}
              \begin{itemize}
                  \item \textbf{Constat :} Majorité des auteurs/serveurs hors Cameroun. CRI classique = 6 à 18 mois. Réseau 24/7 (CERT-CM $\leftrightarrow$ Points contacts étrangers) sous-doté en personnel linguistique/technique.
                  \item \textbf{Mesure concrète :} Renforcer le \textbf{CERT-CM (ANTIC)} : recrutement analysts bilingues (FR/EN), accords bilatéraux MLAT simplifiés avec pays sources prioritaires. Utilisation plateforme sécurisée (I-24/7 Interpol / Budapest) pour demandes automatisées.
              \end{itemize}
        \item \textbf{Axe 3 : Prévention par l'éducation et la régulation technique (Le bouclier amont).}
              \begin{itemize}
                  \item \textbf{Constat :} Victimes par ignorance (partage OTP, clics phishing, mots de passe faibles). Plateformes lentes à retirer contenus signalés (harcèlement, fausses nouvelles).
                  \item \textbf{Mesure concrète :}
                        \begin{itemize}
                            \item Intégration module \textbf{« Hygiène \& Éthique Numérique »} obligatoire dès la 6ème (Programme MINESEC).
                            \item Modification Loi 2010/013 : Obligation pour \textbf{Plateformes/FAI} d'avoir un \textbf{Point de Contact Légal au Cameroun}. Délais de retrait contraignants : 24h (pédopornographie/terrorisme), 72h (haine/harcèlement/fausses nouvelles) sous peine d'amende journalière lourde.
                            \item Campagnes grand public continues (ANTIC/MINPOSTEL) : « Ne partagez jamais votre OTP », « Vérifiez l'URL avant de payer ».
                        \end{itemize}
              \end{itemize}
    \end{enumerate}
    \textbf{Conclusion :} La loi est le squelette ; la formation des acteurs, la coopération rapide et la prévention éducative sont les muscles. Sans eux, l'arsenal reste « lettre morte ». L'arrivée de l'IA générative (deepfakes, automatisation arnaques) impose une révision législative rapide (amendement Loi 2010/012 prévu 2024/2025).
\end{exoresolu}

\begin{attention}[Les 5 erreurs qui coûtent des points au BEPC]
    \begin{enumerate}
        \item \textbf{Confondre sanction administrative et sanction pénale :} Écrire « L'ANTIC condamne à 2 ans de prison ». \textbf{FAUX.} L'ANTIC/APDP/ARTEL prononcent des sanctions \textbf{administratives} (amendes admin, mise en demeure, suspension). \textbf{Seul le Tribunal (Juge) prononce la prison et les amendes pénales.} $\rightarrow$ \textit{Toujours écrire : « Le Procureur poursuit devant le Tribunal qui peut condamner à... »}.
        \item \textbf{Citer la loi sans l'article :} Écrire « La loi interdit le piratage ». \textbf{INSUFFISANT.} $\rightarrow$ \textit{Obligatoire : « L'article 7 de la Loi n°2010/012 du 21 décembre 2010 réprime l'accès frauduleux... »}.
        \item \textbf{Oublier le régime des mineurs :} Appliquer les peines maximales adultes (ex: 10 ans) à un élève de 3ème (14-15 ans) sans mentionner la division par 2 ou les mesures éducatives. $\rightarrow$ \textit{Systématiquement : « Toutefois, l'auteur étant mineur, le Tribunal pour Enfants applique l'Ordonnance 1962 : mesures éducatives prioritaires, peine max divisée par 2... »}.
        \item \textbf{Mélanger Budapest et Malabo :} Attribuer à Budapest la protection des données personnelles (c'est Malabo / RGPD / Loi 2010/013) ou à Malabo le réseau 24/7 (c'est Budapest Art. 35). $\rightarrow$ \textit{Mémo : Budapest = Cybercriminalité + Entraide + Réseau 24/7 ; Malabo = Cadre Panafricain (Données + Cyber + Coopération UA)}.
        \item \textbf{Réponse « sentiment » sans base légale :} « Je trouve que c'est grave, il faut punir sévèrement. » $\rightarrow$ \textit{ZÉRO POINT.} \textbf{Toute affirmation doit être ancrée :} « L'article X prévoit... », « La Convention Y impose... », « Le principe Z (légalité, proportionnalité) impose... ».
    \end{enumerate}
\end{attention}

\begin{savaistu}[Le saviez-vous ? -- Le Cameroun et la Convention de Budapest]
    Le Cameroun est le \textbf{premier pays d'Afrique centrale} à avoir ratifié la Convention de Budapest (mars 2023). Cela lui permet de siéger au Comité des parties (T-CY) et d'exiger l'entraide judiciaire de 70+ pays signataires (France, USA, Allemagne, Nigeria, Ghana, etc.) pour traquer les cybercriminels hébergés à l'étranger. C'est un outil majeur pour la souveraineté numérique nationale.
\end{savaistu}

\newpage

\coursec{Exercices}

\courssub{Je vérifie mes connaissances}

\begin{exercice}[QCM : le bon texte, la bonne institution]
Pour chaque situation, choisis la bonne réponse en la justifiant par une phrase.

\begin{enumerate}
\item Un élève de Douala pirate la boîte mail de son voisin. Le texte national applicable est :
\begin{enumerate}
\item la loi n° 2010/012 du 21 décembre 2010 relative à la cybersécurité et à la cybercriminalité au Cameroun ;
\item le règlement intérieur du collège ;
\item le code de la route.
\end{enumerate}
\item Une entreprise française collecte les données personnelles de clients européens sans leur accord. Le texte applicable est :
\begin{enumerate}
\item la loi camerounaise de 2010 ;
\item le RGPD (Règlement Général sur la Protection des Données) ;
\item la Convention de Genève.
\end{enumerate}
\item Un réseau de cybercriminels opère depuis trois pays différents. L'instrument international de coopération contre la cybercriminalité est :
\begin{enumerate}
\item la Convention de Budapest ;
\item le traité de Versailles ;
\item la charte de l'UNESCO sur le patrimoine.
\end{enumerate}
\item Au Cameroun, l'institution chargée de la sécurité des systèmes d'information de l'État et de la lutte contre la cybercriminalité est :
\begin{enumerate}
\item l'ANTIC ;
\item la CAMTEL ;
\item la CRTV.
\end{enumerate}
\item L'organisation internationale de police qui facilite la coopération entre les polices des pays membres contre les cybercriminels est :
\begin{enumerate}
\item l'UNICEF ;
\item Interpol ;
\item l'OMC.
\end{enumerate}
\end{enumerate}
\end{exercice}

\begin{exercice}[Vrai ou faux justifié]
Réponds par VRAI ou FAUX, puis justifie chaque réponse en une ou deux phrases en citant un texte ou un principe.

\begin{enumerate}
\item « Internet est un espace de liberté totale : on peut y faire et y dire tout ce qu'on veut. »
\item « Le RGPD protège automatiquement tous les Camerounais dans toutes leurs activités en ligne. »
\item « Partager une fausse information sur WhatsApp n'engage pas ma responsabilité, car je ne l'ai pas écrite moi-même. »
\item « La Convention de Budapest permet aux pays de coopérer pour enquêter sur la cybercriminalité. »
\end{enumerate}
\end{exercice}

\begin{exercice}[Définitions]
Définis en deux ou trois phrases chacun des termes suivants, en donnant un exemple pour chacun :
\begin{enumerate}
\item la cybercriminalité ;
\item le phishing (hameçonnage) ;
\item la diffamation en ligne ;
\item une sanction pénale.
\end{enumerate}
\end{exercice}

\begin{exercice}[Infractions et peines : calcul direct]
La loi camerounaise de 2010 prévoit pour certaines infractions numériques des peines de prison et des amendes. On donne : escroquerie informatique : amende pouvant atteindre 5 000 000 FCFA ; accès frauduleux à un système informatique : amende pouvant atteindre 3 000 000 FCFA.

\begin{enumerate}
\item Un cybercriminel est condamné aux deux peines maximales cumulées. Calcule le montant total des amendes.
\item Une victime d'escroquerie a perdu 850 000 FCFA. L'amende maximale pour escroquerie couvre-t-elle cette perte ? Justifie par un calcul.
\end{enumerate}
\end{exercice}

\courssub{Je m'entraîne}

\begin{exercice}[Tableau à compléter : infractions et sanctions]
Recopie et complète le tableau suivant en associant à chaque infraction sa description et la sanction encourue (peine de prison et/ou amende) selon la loi camerounaise de 2010.

\begin{center}
\begin{tabularx}{\linewidth}{|l|X|X|}
\hline
\rowcolor{popblueL} \textbf{Infraction} & \textbf{Description} & \textbf{Sanction encourue} \\
\hline
Piratage (accès frauduleux) & & \\
\hline
Phishing (hameçonnage) & & \\
\hline
Diffamation en ligne & & \\
\hline
Fausse information & & \\
\hline
Atteinte à la vie privée & & \\
\hline
Escroquerie informatique & & \\
\hline
\end{tabularx}
\end{center}
\end{exercice}

\begin{exercice}[Qui est compétent ?]
Pour chacune des situations suivantes, indique l'institution compétente (ANTIC, police/gendarmerie nationale, Interpol) et justifie ton choix.

\begin{enumerate}
\item Marlyse, élève à Yaoundé, découvre que quelqu'un a créé un faux compte Facebook à son nom.
\item Un groupe installé à l'étranger envoie des SMS frauduleux « Orange Money » à des milliers de Camerounais.
\item Le site web d'un ministère camerounais est défiguré par des pirates.
\item Une enquête sur un réseau de fraude bancaire implique des suspects au Cameroun, au Nigeria et en France.
\end{enumerate}
\end{exercice}

\begin{exercice}[Comparer les sanctions]
Construis un tableau comparant les sanctions nationales (Cameroun) et internationales selon les trois critères suivants : champ d'application, nature des peines, institutions chargées de les appliquer. Rédige ensuite deux phrases de conclusion qui dégagent la principale différence entre les deux niveaux.
\end{exercice}

\begin{exercice}[La bonne conduite à tenir]
Ton petit frère reçoit un message : « Félicitations ! Vous avez gagné 100 000 FCFA. Envoyez 2 000 FCFA de frais par Mobile Money au 6XX XX XX XX pour recevoir votre gain. »

\begin{enumerate}
\item De quel type d'infraction s'agit-il ?
\item Rédige la liste des quatre conseils que tu lui donnes (ne pas payer, conserver les preuves, etc.).
\item Indique à qui il peut signaler cette tentative d'escroquerie.
\end{enumerate}
\end{exercice}

\courssub{Je me prépare au BEPC}

\begin{exercice}[Étude de cas : le faux site de Mobile Money]
À Bafoussam, un jeune de 19 ans crée un faux site web imitant une page officielle de Mobile Money. Il envoie des SMS à des commerçants en leur demandant de « confirmer leur code secret » pour sécuriser leur compte. En trois semaines, douze victimes lui communiquent leurs codes ; il leur soutire au total 1 450 000 FCFA. L'une des victimes, Mme Ngo Bell, découvre la fraude en vérifiant son solde.

\begin{enumerate}
\item Qualifie les infractions commises par ce jeune (au moins deux).
\item Cite le texte de loi national applicable et précise son intitulé complet (numéro et date).
\item Décris la procédure que Mme Ngo Bell doit suivre pour obtenir justice (dépôt de plainte, preuves à conserver, services compétents).
\item Indique les sanctions encourues par le coupable (nature des peines).
\item Explique pourquoi la création d'un faux site imitant un service officiel est une circonstance aggravante pour la confiance dans les services numériques.
\end{enumerate}
\end{exercice}

\begin{exercice}[Question de synthèse : la coopération internationale]
Un réseau de cybercriminels envoie ses messages frauduleux depuis un pays A, héberge ses faux sites dans un pays B et escroque des victimes dans un pays C.

Rédige un paragraphe argumenté de 8 à 10 lignes expliquant pourquoi la lutte contre la cybercriminalité exige une coopération internationale. Tu citeras au moins deux instruments ou organisations internationaux (Convention de Budapest, Interpol, RGPD) et tu montreras ce qui arriverait si chaque pays agissait seul.
\end{exercice}

\begin{exercice}[Sujet complet type BEPC]
\textbf{Partie A — Connaissances (6 points).}
\begin{enumerate}
\item Donne l'intitulé complet de la loi camerounaise qui réprime la cybercriminalité. (1 pt)
\item Cite deux institutions nationales intervenant dans la lutte contre la cybercriminalité et précise le rôle de chacune. (2 pts)
\item Cite deux instruments ou organisations internationaux de lutte contre la cybercriminalité. (1 pt)
\item Définis : phishing ; diffamation en ligne. (2 pts)
\end{enumerate}

\textbf{Partie B — Étude de cas (8 points).}
Sandrine, élève en classe de 3ème, publie sur un réseau social une photo de sa camarade Aïcha avec un commentaire mensonger qui la ridiculise. La publication est partagée 300 fois. Humiliée, Aïcha ne veut plus aller au collège. Par ailleurs, le père d'Aïcha reçoit le même jour un SMS lui annonçant un gain à la loterie et lui demandant ses coordonnées bancaires.

\begin{enumerate}
\item Identifie les deux infractions numériques présentes dans ce récit. (2 pts)
\item Pour chaque infraction, indique le texte applicable et la sanction encourue. (2 pts)
\item Explique pourquoi le fait de « partager » la publication engage aussi la responsabilité des autres élèves. (2 pts)
\item Propose deux conseils que tu donnerais à Sandrine pour un usage responsable des réseaux sociaux. (2 pts)
\end{enumerate}

\textbf{Partie C — Argumentation (6 points).}
Rédige un texte de 8 à 10 lignes répondant à la question : « Internet est-il un espace de liberté totale ? » Tu prendras clairement position, tu appuieras ton argumentation sur au moins un texte de loi national et un principe international, et tu concluras par une phrase donnant ton avis personnel d'élève responsable.
\end{exercice}

\newpage

\coursec{Corrigés des exercices}

\begin{corrige}[Exercice 1 — QCM : le bon texte, la bonne institution]
\begin{enumerate}
\item \textbf{Réponse a} : la loi n° 2010/012 du 21 décembre 2010 relative à la cybersécurité et à la cybercriminalité au Cameroun. Le piratage d'une boîte mail est un accès frauduleux à un système informatique, infraction prévue et punie par ce texte national ; le règlement intérieur et le code de la route ne concernent pas les infractions numériques.
\item \textbf{Réponse b} : le RGPD. Ce règlement européen s'applique à toute entreprise qui traite les données personnelles de personnes se trouvant dans l'Union européenne, et la collecte sans consentement y est interdite.
\item \textbf{Réponse a} : la Convention de Budapest. C'est le principal traité international qui organise la coopération entre États pour enquêter sur la cybercriminalité et poursuivre les cybercriminels, même lorsqu'ils opèrent depuis plusieurs pays.
\item \textbf{Réponse a} : l'ANTIC (Agence Nationale des Technologies de l'Information et de la Communication). Elle est chargée de la sécurité des systèmes d'information de l'État et participe à la lutte contre la cybercriminalité ; la CAMTEL est un opérateur de télécommunications et la CRTV un média.
\item \textbf{Réponse b} : Interpol. Cette organisation de police internationale facilite la coopération entre les polices des pays membres, notamment pour traquer les cybercriminels à travers les frontières.
\end{enumerate}
\end{corrige}

\begin{corrige}[Exercice 2 — Vrai ou faux justifié]
\begin{enumerate}
\item \textbf{FAUX.} Internet est un espace de liberté, mais cette liberté a des limites fixées par la loi : la loi n° 2010/012 du 21 décembre 2010 punit par exemple le piratage, l'escroquerie en ligne et la diffamation. La liberté de chacun s'arrête là où commencent les droits des autres.
\item \textbf{FAUX.} Le RGPD est un règlement européen : il protège les personnes se trouvant dans l'Union européenne. Un Camerounais n'en bénéficie que lorsque ses données sont traitées par une entreprise visée par ce règlement ; au Cameroun, c'est le droit national qui s'applique.
\item \textbf{FAUX.} Celui qui partage une fausse information la diffuse et contribue à sa propagation : il engage donc aussi sa responsabilité. La loi camerounaise de 2010 réprime la diffusion de fausses informations par voie électronique, même si l'on n'en est pas l'auteur.
\item \textbf{VRAI.} La Convention de Budapest est un traité international qui permet aux pays signataires de coopérer : échange d'informations, entraide judiciaire et enquêtes communes contre la cybercriminalité.
\end{enumerate}
\end{corrige}

\begin{corrige}[Exercice 3 — Définitions]
\begin{enumerate}
\item \textbf{La cybercriminalité} désigne l'ensemble des infractions commises au moyen d'un ordinateur, d'un téléphone connecté ou d'un réseau comme Internet. Elle regroupe le piratage, l'escroquerie en ligne, la diffamation numérique, etc. \emph{Exemple :} pirater la boîte mail d'un camarade pour lire ses messages.
\item \textbf{Le phishing (hameçonnage)} est une technique qui consiste à envoyer de faux messages ou à créer de faux sites imitant un service officiel (banque, Mobile Money) pour voler des informations confidentielles : mots de passe, codes secrets, coordonnées bancaires. \emph{Exemple :} un SMS demandant de « confirmer son code Mobile Money » sur un faux site.
\item \textbf{La diffamation en ligne} est le fait de publier sur Internet (réseaux sociaux, blogs, messageries) des propos mensongers qui portent atteinte à l'honneur ou à la réputation d'une personne. \emph{Exemple :} accuser faussement un commerçant de vol sur Facebook.
\item \textbf{Une sanction pénale} est une punition prononcée par un tribunal contre l'auteur d'une infraction. Elle peut consister en une peine de prison (emprisonnement) et/ou en une amende (somme d'argent à payer). \emph{Exemple :} un escroc condamné à 2 ans de prison et à une amende de 1 000 000 FCFA.
\end{enumerate}
\end{corrige}

\begin{corrige}[Exercice 4 — Infractions et peines : calcul direct]
\begin{enumerate}
\item Le montant total des amendes est la somme des deux peines maximales :
\[ 5\,000\,000 + 3\,000\,000 = 8\,000\,000 \text{ FCFA}. \]
Le cybercriminel paie donc au total \textbf{8 000 000 FCFA} d'amendes.
\item On compare l'amende maximale à la perte subie :
\[ 5\,000\,000 - 850\,000 = 4\,150\,000 \text{ FCFA}. \]
L'amende maximale (5 000 000 FCFA) est supérieure à la perte de 850 000 FCFA : elle la couvre donc largement, avec un excédent de \textbf{4 150 000 FCFA}. Attention toutefois : l'amende est payée à l'État ; pour être dédommagée, la victime doit demander des dommages et intérêts au tribunal.
\end{enumerate}
\end{corrige}

\begin{corrige}[Exercice 5 — Tableau à compléter : infractions et sanctions]
\begin{center}
\begin{tabularx}{\linewidth}{|l|X|X|}
\hline
\rowcolor{popblueL} \textbf{Infraction} & \textbf{Description} & \textbf{Sanction encourue} \\
\hline
Piratage (accès frauduleux) & Entrer ou se maintenir sans autorisation dans un système informatique (compte, boîte mail, site). & Peine de prison et amende pouvant atteindre 3 000 000 FCFA. \\
\hline
Phishing (hameçonnage) & Créer de faux messages ou de faux sites pour voler des données confidentielles (mots de passe, codes). & Peine de prison et amende (souvent plusieurs millions de FCFA), car c'est une forme d'escroquerie. \\
\hline
Diffamation en ligne & Publier des propos mensongers qui nuisent à l'honneur ou à la réputation d'une personne. & Peine de prison et/ou amende, plus d'éventuels dommages et intérêts à la victime. \\
\hline
Fausse information & Diffuser volontairement par voie électronique une information que l'on sait fausse. & Peine de prison et/ou amende prévues par la loi de 2010. \\
\hline
Atteinte à la vie privée & Publier sans accord des photos, vidéos ou informations intimes d'une personne. & Peine de prison et/ou amende, avec réparation du préjudice. \\
\hline
Escroquerie informatique & Tromper des personnes par des moyens informatiques pour leur soutirer de l'argent. & Peine de prison et amende pouvant atteindre 5 000 000 FCFA. \\
\hline
\end{tabularx}
\end{center}
\textbf{Point de vigilance :} les montants exacts des amendes et les durées de prison varient selon la gravité des faits et les circonstances ; l'essentiel est de retenir que chaque infraction entraîne \textbf{une peine de prison et/ou une amende}, prononcées par un tribunal.
\end{corrige}

\begin{corrige}[Exercice 6 — Qui est compétent ?]
\begin{enumerate}
\item \textbf{Police / gendarmerie nationale.} La création d'un faux compte est une infraction (usurpation d'identité, atteinte à la vie privée) : Marlyse, accompagnée de ses parents, doit porter plainte au commissariat ou à la brigade de gendarmerie le plus proche, avec des captures d'écran comme preuves.
\item \textbf{ANTIC et police nationale, avec coopération internationale.} L'ANTIC, chargée de la sécurité des systèmes d'information, peut tracer les envois massifs de SMS frauduleux ; comme le groupe est installé à l'étranger, la coopération avec les polices étrangères (via Interpol) est nécessaire.
\item \textbf{ANTIC.} La défiguration du site d'un ministère touche un système d'information de l'État : c'est précisément la mission de l'ANTIC de sécuriser ces systèmes et de contribuer à l'enquête, en lien avec la police judiciaire.
\item \textbf{Interpol.} L'enquête concerne trois pays (Cameroun, Nigeria, France) : seule une organisation de coopération policière internationale comme Interpol permet d'échanger les informations et de coordonner les arrestations entre les polices des trois États.
\end{enumerate}
\end{corrige}

\begin{corrige}[Exercice 7 — Comparer les sanctions]
\begin{center}
\begin{tabularx}{\linewidth}{|l|X|X|}
\hline
\rowcolor{popblueL} \textbf{Critère} & \textbf{Niveau national (Cameroun)} & \textbf{Niveau international} \\
\hline
Champ d'application & Territoire camerounais : loi n° 2010/012 du 21 décembre 2010. & Ensemble des pays liés par un traité (Convention de Budapest) ou un règlement (RGPD pour l'Europe). \\
\hline
Nature des peines & Peines de prison et amendes en FCFA, prononcées par les tribunaux camerounais. & Amendes parfois très élevées (RGPD : jusqu'à des millions d'euros), extradition et poursuites possibles grâce à la coopération. \\
\hline
Institutions chargées de les appliquer & ANTIC, police et gendarmerie nationales, tribunaux camerounais. & Interpol, autorités de protection des données (ex. CNIL en France), justice de chaque pays signataire. \\
\hline
\end{tabularx}
\end{center}
\textbf{Conclusion (deux phrases) :} La principale différence est que les sanctions nationales ne s'appliquent que sur le territoire du Cameroun, alors que les instruments internationaux permettent de poursuivre les cybercriminels au-delà des frontières. Les deux niveaux sont donc complémentaires : la loi nationale punit les infractions commises au pays, et la coopération internationale empêche les criminels de se cacher à l'étranger.
\end{corrige}

\begin{corrige}[Exercice 8 — La bonne conduite à tenir]
\begin{enumerate}
\item Il s'agit d'une \textbf{tentative d'escroquerie informatique} (arnaque au faux gain), souvent réalisée par phishing : on promet un gain fictif pour soutirer de l'argent ou des données personnelles.
\item Les quatre conseils :
\begin{enumerate}
\item \textbf{Ne pas payer} et ne jamais envoyer d'argent ni de code secret : un vrai gain ne demande jamais de frais à l'avance.
\item \textbf{Ne pas répondre} au message et ne pas cliquer sur les liens qu'il contient.
\item \textbf{Conserver les preuves} : garder le SMS, noter le numéro de l'expéditeur, faire des captures d'écran.
\item \textbf{En parler à un adulte} (parents, enseignant) et bloquer le numéro de l'escroc.
\end{enumerate}
\item Il peut signaler cette tentative à la \textbf{police ou à la gendarmerie nationale} (dépôt de plainte), ainsi qu'à l'\textbf{ANTIC} et à son opérateur téléphonique (MTN, Orange) qui peut bloquer le numéro frauduleux.
\end{enumerate}
\end{corrige}

\begin{corrige}[Exercice 9 — Étude de cas : le faux site de Mobile Money]
\begin{enumerate}
\item Le jeune commet au moins deux infractions : le \textbf{phishing} (création d'un faux site imitant un service officiel pour voler les codes secrets) et l'\textbf{escroquerie informatique} (soutirer 1 450 000 FCFA à douze victimes par tromperie). On peut ajouter l'\textbf{accès frauduleux} aux comptes des victimes.
\item Le texte applicable est la \textbf{loi n° 2010/012 du 21 décembre 2010 relative à la cybersécurité et à la cybercriminalité au Cameroun}.
\item Procédure pour Mme Ngo Bell :
\begin{enumerate}
\item conserver toutes les \textbf{preuves} : SMS reçus, captures d'écran du faux site, relevés des transactions Mobile Money ;
\item contacter son \textbf{opérateur} pour signaler la fraude et tenter de sécuriser son compte (changement de code) ;
\item \textbf{porter plainte} au commissariat de police ou à la brigade de gendarmerie de Bafoussam, qui transmettra aux services spécialisés, avec l'appui de l'ANTIC pour l'enquête technique ;
\item se constituer éventuellement \textbf{partie civile} au procès pour obtenir des dommages et intérêts.
\end{enumerate}
\item Le coupable encourt une \textbf{peine de prison} (emprisonnement) et une \textbf{amende} pouvant atteindre plusieurs millions de FCFA (jusqu'à 5 000 000 FCFA pour l'escroquerie informatique), ainsi que l'obligation de \textbf{rembourser les victimes}.
\item Créer un faux site imitant un service officiel est aggravant car cela trompe la confiance légitime des usagers : les victimes croient avoir affaire à un service sérieux. Si de tels actes se multiplient, les citoyens n'osent plus utiliser les services numériques (paiement mobile, administration en ligne), ce qui freine le développement du numérique au Cameroun.
\end{enumerate}
\textbf{Points de vigilance :} citer le numéro ET la date de la loi ; distinguer clairement les deux infractions ; mentionner les preuves à conserver, qui sont indispensables à l'enquête.
\end{corrige}

\begin{corrige}[Exercice 10 — Question de synthèse : la coopération internationale]
\textbf{Corrigé type (8 à 10 lignes) :}

La cybercriminalité ne connaît pas les frontières : un même réseau peut envoyer ses messages frauduleux depuis un pays A, héberger ses faux sites dans un pays B et escroquer des victimes dans un pays C. Si chaque pays agissait seul, sa police s'arrêterait à ses frontières : le pays C pourrait constater les dégâts mais jamais arrêter les criminels installés en A ni fermer les sites hébergés en B. Les cybercriminels profiteraient alors des différences entre les lois pour échapper à toute sanction. C'est pourquoi la coopération internationale est indispensable. La \textbf{Convention de Budapest} permet aux pays signataires d'harmoniser leurs lois, d'échanger rapidement des preuves électroniques et de s'entraider judiciairement. \textbf{Interpol}, de son côté, met en relation les polices des pays membres pour coordonner les enquêtes et les arrestations. Des règlements comme le \textbf{RGPD} imposent enfin des règles communes de protection des données. Grâce à ces instruments, les criminels ne peuvent plus se cacher derrière une frontière : la lutte contre la cybercriminalité est donc forcément un effort partagé entre les nations.

\textbf{Barème indicatif :} problématique comprise (2 pts) ; deux instruments cités et expliqués (2 pts) ; conséquence de l'action isolée montrée (1 pt) ; clarté et correction de la langue (1 pt).
\end{corrige}

\begin{corrige}[Exercice 11 — Sujet complet type BEPC]
\textbf{Partie A — Connaissances (6 points).}
\begin{enumerate}
\item Loi n° 2010/012 du 21 décembre 2010 relative à la cybersécurité et à la cybercriminalité au Cameroun. \emph{(1 pt)}
\item Deux institutions (au choix) : l'\textbf{ANTIC}, chargée de la sécurité des systèmes d'information de l'État et de l'appui technique aux enquêtes ; la \textbf{police / gendarmerie nationale}, qui reçoit les plaintes, enquête et arrête les auteurs ; les \textbf{tribunaux}, qui jugent et prononcent les sanctions. \emph{(2 pts : 1 pt par institution avec son rôle)}
\item La \textbf{Convention de Budapest} (coopération entre États) et \textbf{Interpol} (coopération policière internationale) ; on accepte aussi le RGPD. \emph{(1 pt)}
\item \textbf{Phishing} : technique frauduleuse consistant à imiter un service officiel (faux SMS, faux site) pour voler des données confidentielles. \textbf{Diffamation en ligne} : publication sur Internet de propos mensongers portant atteinte à l'honneur ou à la réputation d'une personne. \emph{(2 pts : 1 pt par définition)}
\end{enumerate}

\textbf{Partie B — Étude de cas (8 points).}
\begin{enumerate}
\item Les deux infractions sont : la \textbf{diffamation en ligne} (photo accompagnée d'un commentaire mensonger ridiculisant Aïcha) et la \textbf{tentative d'escroquerie par phishing} (faux SMS de loterie demandant les coordonnées bancaires du père). \emph{(2 pts)}
\item Les deux infractions tombent sous le coup de la \textbf{loi n° 2010/012 du 21 décembre 2010}. Diffamation en ligne : peine de prison et/ou amende, avec dommages et intérêts pour Aïcha. Escroquerie informatique (même tentée) : peine de prison et amende pouvant atteindre 5 000 000 FCFA. \emph{(2 pts)}
\item Partager la publication, c'est la diffuser à de nouvelles personnes et aggraver l'atteinte à la réputation d'Aïcha : chaque élève qui partage contribue à l'infraction et engage donc sa propre responsabilité, même s'il n'est pas l'auteur du message initial. \emph{(2 pts)}
\item Conseils à Sandrine : réfléchir avant de publier (« est-ce vrai ? est-ce respectueux ? ») et ne jamais publier la photo de quelqu'un sans son accord ; supprimer la publication, présenter des excuses à Aïcha et signaler les contenus blessants au lieu de les partager. \emph{(2 pts)}
\end{enumerate}

\textbf{Partie C — Argumentation (6 points). Corrigé type :}

Internet n'est pas un espace de liberté totale. Certes, il permet de s'informer, d'apprendre et de s'exprimer librement, et cette liberté est précieuse. Mais comme dans la vie réelle, elle a des limites : au Cameroun, la loi n° 2010/012 du 21 décembre 2010 punit le piratage, l'escroquerie en ligne, la diffamation et la diffusion de fausses informations. Sur le plan international, des principes communs, comme ceux de la Convention de Budapest ou du RGPD sur la protection des données, rappellent que les droits de chacun doivent être respectés partout. Dire ou faire en ligne ce qui est interdit hors ligne n'est donc pas une liberté, mais une infraction. La vraie liberté consiste à utiliser Internet de manière responsable, en respectant les autres. En tant qu'élève responsable, je pense qu'Internet est un outil formidable à condition que chacun y respecte la loi et la dignité des personnes.

\textbf{Barème indicatif :} position claire (1 pt) ; texte national cité correctement (1,5 pt) ; principe international cité (1,5 pt) ; conclusion personnelle (1 pt) ; qualité de la rédaction (1 pt).

\textbf{Points de vigilance :} donner l'intitulé complet de la loi (numéro + date) ; en Partie B, bien distinguer les deux infractions et justifier la responsabilité du partage ; en Partie C, prendre clairement position et conclure par un avis personnel, comme demandé.
\end{corrige}

\newpage

\coursec{Fiche bilan}

\begin{popfigure}
\centering
\begin{tikzpicture}[
  mindmap,
  grow cyclic,
  every node/.style={concept, execute at begin node=\hskip0pt},
  concept color=popblue,
  level 1/.append style={level distance=4.5cm, sibling angle=60, font=\small\bfseries},
  level 2/.append style={level distance=3cm, sibling angle=45, font=\footnotesize},
  branch1/.style={concept color=popblue},
  branch2/.style={concept color=poporange},
  branch3/.style={concept color=popgreen},
  branch4/.style={concept color=poppurple},
  branch5/.style={concept color=poppink},
  branch6/.style={concept color=popgold},
]
\node[concept, text width=3.5cm, align=center, font=\bfseries\large, minimum size=2.5cm] {Éthique Numérique\\ \& Cadre Juridique}
  child[branch1] { node[concept, text width=2.8cm, align=center] {1. Cadre National\\ (Lois 2010, 2024)} }
  child[branch2] { node[concept, text width=2.8cm, align=center] {2. Institutions\\ (ANTIC, MINPOSTEL, Justice)} }
  child[branch3] { node[concept, text width=2.8cm, align=center] {3. Infractions \& Sanctions\\ Nationales} }
  child[branch4] { node[concept, text width=2.8cm, align=center] {4. Cadre International\\ (Budapest, Malabo, ONU)} }
  child[branch5] { node[concept, text width=2.8cm, align=center] {5. Comparaison\\ National vs International} }
  child[branch6] { node[concept, text width=2.8cm, align=center] {6. Démarche\\ Argumentée} };
\end{tikzpicture}
\legende{Carte mentale de la leçon 10 : Droit et éthique numérique (niveau 3ème)}
\end{popfigure}

\begin{bilan}
\courssub{Définitions clés}
\begin{itemize}
  \item \cle{Cybercriminalité} : Ensemble des infractions pénales commises à l'aide de systèmes d'information ou de réseaux de télécommunications (Loi n°2010/012).
  \item \cle{Données à caractère personnel (DCP)} : Toute information relative à une personne physique identifiée ou identifiable (Loi n°2024/006).
  \item \cle{Infraction numérique} : Violation d'une disposition légale liée aux TIC (accès frauduleux, interception, contrefaçon logicielle, diffusion de fausses nouvelles, etc.).
  \item \cle{Sanction} : Peine prononcée par l'autorité judiciaire (prison, amende, confiscation) ou administrative (avertissement, suspension de licence) en réponse à une infraction.
  \item \cle{ANTIC} (Agence Nationale des Technologies de l'Information et de la Communication) : Autorité de régulation, de contrôle et de sanction administrative au Cameroun.
\end{itemize}

\courssub{Textes fondateurs à retenir (Cameroun)}
\begin{center}
\begin{tabularx}{\linewidth}{|l|X|}\hline
\rowcolor{popblueL}
\textbf{Texte} & \textbf{Objet principal} \\\hline
Loi n°2010/012 du 21 déc. 2010 & Cybercriminalité et cyber-sécurité (infractions, peines, procédure). \\\hline
Loi n°2010/013 du 21 déc. 2010 & Communications électroniques (régulation, licences, ANTIC). \\\hline
Loi n°2024/006 du 19 juil. 2024 & Protection des données à caractère personnel (consentement, droits, APDP). \\\hline
Décret n°2012/003 & Organisation et fonctionnement de l'ANTIC. \\\hline
Code pénal (Loi n°2016/007) & Infractions classiques transposées au numérique (escroquerie, diffamation). \\\hline
\end{tabularx}
\end{center}

\courssub{Cadre international de référence}
\begin{itemize}
  \item \textbf{Convention de Budapest (2001)} : Premier traité international contraignant sur la cybercriminalité (harmonisation des législations, coopération judiciaire, extradition). Le Cameroun y a adhéré (loi n°2022/015).
  \item \textbf{Convention de Malabo (UA, 2014)} : Cadre africain sur la sécurité informatique et la protection des données personnelles. Ratifiée par le Cameroun.
  \item \textbf{Résolutions ONU / UIT} : Normes pour un comportement responsable des États dans le cyberespace, lutte contre la cybercriminalité transnationale.
  \item \textbf{RGPD (UE, 2016)} : Référence mondiale pour la protection des données (droit à l'oubli, portabilité, DPO), inspire la loi camerounaise 2024.
\end{itemize}

\courssub{Méthode express : Analyser une situation et qualifier l'infraction}
\begin{enumerate}
  \item \textbf{Identifier les faits} : Quel acte ? Quel outil numérique ? Quel préjudice (données, argent, honneur, système) ?
  \item \textbf{Rechercher la qualification légale} : Faire correspondre les faits à un article précis (ex: Art. 7 Loi 2010/012 = Accès frauduleux ; Art. 14 = Escroquerie électronique).
  \item \textbf{Déterminer la juridiction} : Territoriale (lieu de l'infraction ou du résultat) + Nationalité de l'auteur (principe de l'universalité pour certains crimes).
  \item \textbf{Citer la sanction encourue} : Peine de prison (jours/années) + Amende (FCFA) + Peines complémentaires (confiscation matériel, interdiction de droits).
  \item \textbf{Argumenter la responsabilité} : Intentionnalité (dol), négligence, complicité, tentative. Distinguer auteur, coauteur, complice.
\end{enumerate}

\courssub{Tableau comparatif : Sanctions Nationales vs Internationales}
\begin{center}
\begin{tabularx}{\linewidth}{|l|X|X|}\hline
\rowcolor{popblueL}
\textbf{Critère} & \textbf{Niveau National (Cameroun)} & \textbf{Niveau International} \\\hline
\textbf{Base légale} & Lois 2010/012, 2010/013, 2024/006, Code pénal & Convention Budapest, Convention Malabo, Résolutions ONU \\\hline
\textbf{Portée} & Territoriale (principalement) + extraterritoriale limitée (Art. 54 L.2010/012) & Transnationale : coopération, entraide, extradition \\\hline
\textbf{Types de sanctions} & Pénales (prison, amende), Administratives (ANTIC), Civiles (dommages-intérêts) & Principalement obligation de \emph{légiférer} et de \emph{coopérer} pour les États parties \\\hline
\textbf{Exécution} & Tribunaux camerounais (TGI, Cour d'Appel, Cour Suprême) & Dépend de la transposition en droit interne par chaque État \\\hline
\textbf{Exemple concret} & Hacker d'un compte Mobile Money à Yaoundé : 6 mois--5 ans prison + 500\,000--10\,000\,000 FCFA & Pirate russe attaquant une banque camerounaise : Demande d'entraide judiciaire via Budapest \\\hline
\end{tabularx}
\end{center}
\end{bilan}

\begin{aretenir}[Checklist avant le BEPC]
\begin{itemize}
  \item $\square$ Je sais citer les 3 lois camerounaises majeures du numérique (2010/012, 2010/013, 2024/006) et leur objet.
  \item $\square$ Je connais le rôle de l'ANTIC (régulation, contrôle, sanction administrative, CERT-CM).
  \item $\square$ Je sais différencier sanction \textbf{pénale} (juge, prison/amende) et sanction \textbf{administrative} (ANTIC, avertissement/suspension).
  \item $\square$ Je sais qualifier 5 infractions types : accès frauduleux, interception, escroquerie électronique, atteintes aux DCP, diffusion fausses nouvelles.
  \item $\square$ Je connais les peines maximales : jusqu'à 10 ans de prison et 10\,000\,000 FCFA d'amende pour les cas graves (Art. 7, 14, 21 L.2010/012).
  \item $\square$ Je sais expliquer l'apport de la \textbf{Convention de Budapest} (harmonisation, coopération 24/7, double incrimination).
  \item $\square$ Je sais expliquer l'apport de la \textbf{Convention de Malabo} (cadre africain, protection DCP, sécurité SI).
  \item $\square$ Je comprends le principe de \textbf{territorialité} et ses exceptions (effet sur le territoire, nationalité de l'auteur).
  \item $\square$ Je sais structurer une réponse argumentée : Faits $\rightarrow$ Qualification (Article) $\rightarrow$ Sanction $\rightarrow$ Juridiction $\rightarrow$ Responsabilité.
  \item $\square$ Je sais donner un exemple camerounais concret (ex: arnaque Orange Money, piratage site web école, fuite données élèves).
  \item $\square$ Je connais les droits du titulaire des données (accès, rectification, opposition, effacement, portabilité) selon la Loi 2024/006.
\end{itemize}
\end{aretenir}

\findecours

\end{document}
